% Les acides carboxyliques et leurs dérivés — Chimie Tle C — Cours signé OnBuch+
% Programme officiel MINESEC, Terminale C, D, E. Compiler avec : tectonic L02-acides-carboxyliques-et-derives.tex
\newcommand{\DOCMATIERE}{Chimie}
\newcommand{\DOCNIVEAU}{Tle C}
\newcommand{\DOCMODULE}{Module 1 · Chimie organique}
\newcommand{\DOCLECON}{02}
\newcommand{\DOCTITRE}{Les acides carboxyliques et leurs dérivés}
% OnBuch+ — préambule « cours signés » ENRICHI (compilé avec tectonic / XeLaTeX)
% Système visuel « Pop » d'OnBuch+ : fond crème (jamais blanc), bordures encre
% épaisses, ombres « dures » décalées, titres Archivo Black en capitales.
% Version enrichie pour les cours de chimie : chimie (mhchem, chemfig),
% graphes (pgfplots), boîtes pédagogiques supplémentaires (définition,
% propriété, méthode, attention, le savais-tu, exercice résolu, exercice,
% TP, bilan), page de garde refaite et signature OnBuch+ ancrée en bas.
%
% Macros attendues dans main.tex AVANT \input{preamble} :
%   \DOCMATIERE  \DOCNIVEAU  \DOCMODULE  \DOCLECON (numéro)  \DOCTITRE
\documentclass[11pt]{article}

\usepackage{fontspec}
\usepackage[french]{babel}
\usepackage[a4paper,margin=2cm,top=3.4cm,bottom=2.8cm,headheight=1.4cm,footskip=1.3cm]{geometry}
\usepackage{amsmath,amssymb,amsfonts}
\usepackage[table]{xcolor} % [table] : fournit aussi \rowcolors (alternance de lignes)
\usepackage{pagecolor}
\usepackage{tikz}
\usetikzlibrary{arrows.meta,positioning,calc,shapes.geometric,shapes.misc,shapes.arrows,decorations.pathmorphing,decorations.markings,patterns,shadows,mindmap,trees,fit,backgrounds,matrix,intersections}
\usepackage{pgfplots}
\pgfplotsset{compat=1.18}
\usepackage[version=4]{mhchem}
\usepackage{chemfig}
\usepackage{enumitem}
\usepackage{fancyhdr}
% [most] charge toutes les bibliothèques tcolorbox (listings, minted, raster,
% documentation...) et gonfle inutilement la mémoire TeX sur un document long
% (~300 boîtes) : on ne charge que ce qui est réellement utilisé.
\usepackage[skins,breakable]{tcolorbox}
\usepackage{graphicx}
\usepackage{colortbl}
\usepackage{booktabs}
\usepackage{tabularx}
\usepackage{array}
\usepackage{multicol}
\usepackage{siunitx}
% parse-numbers=false : accepte aussi \SI{2,5 \times 10^{-3}}{...}
\sisetup{locale=FR,output-decimal-marker={,},group-minimum-digits=4,parse-numbers=false}
\usepackage{qrcode}
% \degree : commande fréquemment utilisée par les modèles pour un angle ou une
% température en dehors de \SI{}{} (non fournie par les paquets chargés).
\providecommand{\degree}{^{\circ}}
\usepackage{lastpage}
\usepackage{hyperref}
\hypersetup{hidelinks}

% ============================================================
% Palette « Pop » OnBuch+ — mêmes hex que la classe P (plus_ui.dart)
% ============================================================
\definecolor{poporange}{HTML}{F59321}
\definecolor{popdark}{HTML}{E07A0C}
\definecolor{popcream}{HTML}{FFF6E8}
\definecolor{popink}{HTML}{1C1714}
\definecolor{poppurple}{HTML}{7A5AE0}
\definecolor{popgreen}{HTML}{2E9E6A}
\definecolor{poppink}{HTML}{E0457B}
\definecolor{popblue}{HTML}{2F7DE1}
\definecolor{popgold}{HTML}{C98A12}
\definecolor{popmuted}{HTML}{8A7F76}
\definecolor{popline}{HTML}{EADFD2}
\definecolor{popgray}{HTML}{8A7F76}
\colorlet{popgrey}{popgray}
% Alias tolérants (le modèle nomme parfois les couleurs différemment)
\colorlet{popwhite}{white}
\colorlet{popblack}{popink}
\colorlet{popred}{poppink}
\colorlet{popyellow}{popgold}
% Teintes claires pour fonds de tableaux / zones
\colorlet{popblueL}{popblue!10!white}
\colorlet{poporangeL}{poporange!14!white}
\colorlet{popgreenL}{popgreen!12!white}
\colorlet{poppurpleL}{poppurple!12!white}
\colorlet{poppinkL}{poppink!12!white}
\colorlet{popgoldL}{popgold!14!white}

\pagecolor{popcream}

% ============================================================
% Typographie
% ============================================================
\defaultfontfeatures{Ligatures=TeX}
\setmainfont{PlusJakartaSans-Medium.ttf}[Path=../../../fonts/,BoldFont=PlusJakartaSans-Bold.ttf]
% Maths en sans-serif (Fira Math) avec chiffres et lettres droites de Plus
% Jakarta, pour que les formules se fondent dans le texte.
\usepackage[math-style=ISO,bold-style=ISO]{unicode-math}
\setmathfont{FiraMath-Regular.otf}[Path=../../../fonts/]
\setmathfont{PlusJakartaSans-Medium.ttf}[Path=../../../fonts/,range={up/num,up/{Latin,latin},\mathup}]
\usepackage{icomma} % virgule décimale sans espace en mode math
% Caractères absents de la police de texte (sinon carré vide) : rendus par
% leur équivalent mathématique, en texte comme en formule.
\usepackage{newunicodechar}
\newunicodechar{α}{\ensuremath{\alpha}}
\newunicodechar{Α}{\ensuremath{\alpha}}
\newunicodechar{β}{\ensuremath{\beta}}
\newunicodechar{δ}{\ensuremath{\delta}}
\newunicodechar{✓}{\popcheck{popgreen}}
\newfontfamily\popdisplay{ArchivoBlack-Regular.ttf}[Path=../../../fonts/]
\newfontfamily\poplogo{SpaceGrotesk-Bold.ttf}[Path=../../../fonts/]
\newfontfamily\popsemi{PlusJakartaSans-SemiBold.ttf}[Path=../../../fonts/]
\newfontfamily\popxbold{PlusJakartaSans-ExtraBold.ttf}[Path=../../../fonts/]
\newfontfamily\popxboldls{PlusJakartaSans-ExtraBold.ttf}[Path=../../../fonts/,LetterSpace=3.0]

\setlist[enumerate]{leftmargin=1.7em,itemsep=5pt,topsep=4pt}
\setlist[itemize]{leftmargin=1.7em,itemsep=4pt,topsep=4pt,label={\color{poporange}\textbullet}}
\setlength{\parindent}{0pt}
\setlength{\parskip}{5pt}
\renewcommand{\arraystretch}{1.25}

% chemfig : liaisons un peu plus épaisses, encre Pop
\setchemfig{atom sep=2em,bond style={line width=0.9pt,popink},cram width=3pt}

% pgfplots : style Pop par défaut
\pgfplotsset{
  every axis/.append style={axis line style={popink,line width=1pt},tick style={popink},
    grid style={popline,line width=0.5pt},label style={font=\small},tick label style={font=\footnotesize}},
  popcurve/.style={poporange,line width=1.8pt,no markers,smooth},
}

% ============================================================
% Pill / squiggle
% ============================================================
\newcommand{\poppill}[3]{%
  \tikz[baseline=-0.55ex]\node[draw=popink,line width=1.2pt,rounded corners=2.6pt,%
    fill=#2,inner xsep=6.5pt,inner ysep=3pt]{%
    \popxboldls\fontsize{7}{7}\selectfont\color{#3}\MakeUppercase{#1}};%
}
\newcommand{\popsquiggle}[1]{%
  \tikz[baseline=-0.4ex]{
    \draw[#1,line width=2.6pt,line cap=round]
      plot [smooth] coordinates {(0,0.09) (0.34,0.01) (0.68,0.21) (1.02,0.01) (1.36,0.10)};
  }%
}
% Mot-clé surligné (marqueur orange sous le texte)
\newcommand{\cle}[1]{{\popxbold\color{popdark}#1}}

% ============================================================
% En-tête / pied
% ============================================================
\pagestyle{fancy}
\fancyhf{}
\renewcommand{\headrulewidth}{0pt}
\renewcommand{\footrulewidth}{0pt}
\lhead{\raisebox{-0.15ex}{\poplogo\fontsize{13}{13}\selectfont\color{popink}OnBuch\color{poporange}+}}
\rhead{\poppill{\DOCMATIERE\ \textbullet\ \DOCNIVEAU\ \textbullet\ Leçon \DOCLECON}{white}{popink}}
\lfoot{\popxbold\fontsize{7.6}{7.6}\selectfont\color{popmuted}\MakeUppercase{Cours signé OnBuch+}\ \textbullet\ {\popsemi\DOCTITRE}}
\rfoot{\tikz[baseline=-0.6ex]\node[rounded corners=8.5pt,fill=popink,minimum height=17pt,minimum width=17pt,inner xsep=5pt,inner sep=0]{\popxbold\fontsize{8}{8}\selectfont\color{popcream}\thepage\,/\,\pageref*{LastPage}};}

% ============================================================
% Titres
% ============================================================
\newcommand{\chaptitre}[2]{%
  \begin{tcolorbox}[enhanced,colback=poporange,colframe=popink,boxrule=2.6pt,arc=11pt,
      drop shadow=popink,left=15pt,right=15pt,top=12pt,bottom=11pt,before skip=3pt,after skip=18pt]
    \poppill{#2}{popink}{popcream}\par\vspace{9pt}
    {\raggedright\hyphenpenalty=10000\exhyphenpenalty=10000\popdisplay\fontsize{22}{24}\selectfont\color{popcream}\MakeUppercase{#1}\par}
  \end{tcolorbox}
}
% Section numérotée : pastille orange avec le numéro + titre Archivo
\newcounter{popsec}
\newcommand{\coursec}[1]{%
  \stepcounter{popsec}\setcounter{popsub}{0}%
  \par\vspace{16pt}\noindent
  \begin{minipage}{\linewidth}
  \tikz[baseline=-0.7ex]\node[draw=popink,line width=1.6pt,fill=poporange,rounded corners=4pt,minimum size=22pt,inner sep=0]{\popdisplay\fontsize{12}{12}\selectfont\color{popcream}\thepopsec};%
  \hspace{8pt}{\popdisplay\fontsize{15}{17}\selectfont\color{popink}\MakeUppercase{#1}}\par
  \vspace{4pt}\noindent{\color{poporange}\rule{36pt}{3.6pt}}
  \end{minipage}\par\nopagebreak\vspace{8pt}
}
\newcounter{popsub}
\newcommand{\courssub}[1]{%
  \stepcounter{popsub}%
  \par\vspace{9pt}\noindent{\popxbold\fontsize{11.5}{13}\selectfont\color{popink}{\color{poporange}\thepopsec.\thepopsub}\hspace{6pt}#1}\par\nopagebreak\vspace{4pt}
}

% ============================================================
% Boîtes pédagogiques — même recette : fond blanc, bordure épaisse colorée,
% ombre dure encre, badge plein. Titre optionnel [#1].
% ============================================================
\tcbset{popbase/.style={breakable,enhanced,colback=white,boxrule=2.3pt,arc=9pt,
  shadow={3pt}{-3pt}{0pt}{popink},left=14pt,right=14pt,top=11pt,bottom=11pt,before skip=11pt,after skip=15pt,
  fontupper=\popsemi\fontsize{10.6}{14.5}\selectfont\color{popink},
  fontlower=\popsemi\fontsize{10.6}{14.5}\selectfont\color{popink}}}
% Coche dessinée (le glyphe ✓ n'existe pas dans les polices du document)
\newcommand{\popcheck}[1]{\tikz[baseline=-0.5ex]\draw[#1,line width=1.6pt,line cap=round,line join=round](-0.13,0.0)--(-0.04,-0.09)--(0.13,0.1);}
\newcommand{\popboxhead}[3]{\poppill{#1}{#2}{white}\ifx\relax#3\relax\else\hspace{6pt}{\popxbold\fontsize{10.6}{12}\selectfont\color{popink}#3}\fi\par\vspace{7pt}}

\newtcolorbox{aretenir}[1][]{popbase,colframe=popgreen,before upper={\popboxhead{À retenir}{popgreen}{#1}}}
\newtcolorbox{exemplebox}[1][]{popbase,colframe=poporange,before upper={\popboxhead{Exemple}{poporange}{#1}}}
\newtcolorbox{definition}[1][]{popbase,colframe=popblue,before upper={\popboxhead{Définition}{popblue}{#1}}}
\newtcolorbox{propriete}[1][]{popbase,colframe=poppurple,before upper={\popboxhead{Propriété}{poppurple}{#1}}}
\newtcolorbox{methode}[1][]{popbase,colframe=popgold,before upper={\popboxhead{Méthode}{popgold}{#1}}}
\newtcolorbox{attention}[1][]{popbase,colframe=poppink,before upper={\popboxhead{Attention}{poppink}{#1}}}
\newtcolorbox{experience}[1][]{popbase,colframe=popblue,colback=popblueL,before upper={\popboxhead{Expérience}{popblue}{#1}}}
% « Le savais-tu ? » : carte violette pleine (contexte, culture, Cameroun)
\newtcolorbox{savaistu}[1][]{popbase,colback=poppurple,colframe=popink,
  fontupper=\popsemi\fontsize{10.4}{14.2}\selectfont\color{white},
  before upper={\poppill{Le savais-tu\,?}{white}{poppurple}\ifx\relax#1\relax\else\hspace{6pt}{\popxbold\color{white}#1}\fi\par\vspace{7pt}}}
% Exercice résolu : énoncé en haut, \tcblower puis la solution sur fond crème
\newcounter{popexo}\newcounter{popexr}
\newtcolorbox{exoresolu}[1][]{popbase,colframe=poporange,colbacklower=popcream,
  segmentation style={popink,line width=1.2pt,dashed},
  before upper={\stepcounter{popexr}\popboxhead{Exercice résolu \thepopexr}{poporange}{#1}},
  before lower={\poppill{Solution}{popink}{popcream}\par\vspace{6pt}}}
% Exercice (non résolu en place) — la correction est dans la partie Corrigés
\newtcolorbox{exercice}[1][]{popbase,colframe=popink,
  before upper={\stepcounter{popexo}\popboxhead{Exercice \thepopexo}{popink}{#1}}}
% Corrigé
\newtcolorbox{corrige}[1][]{popbase,colframe=popgreen,colback=popgreenL,
  before upper={\popboxhead{Corrigé}{popgreen}{#1}}}
% Objectifs / prérequis / bilan
\newtcolorbox{objectifs}{popbase,colframe=popink,colback=poporangeL,
  before upper={\popboxhead{Objectifs de la leçon}{popink}{}}}
\newtcolorbox{prerequis}{popbase,colframe=popink,colback=popblueL,
  before upper={\popboxhead{Prérequis}{popblue}{}}}
\newtcolorbox{bilan}{popbase,colframe=popink,colback=white,boxrule=2.8pt,
  before upper={\popboxhead{Fiche bilan}{poporange}{L'essentiel à connaître pour l'examen}}}
% Figure encadrée avec légende
\newtcolorbox{popfigure}[1][]{enhanced,breakable=false,colback=white,colframe=popline,boxrule=1.4pt,arc=8pt,
  left=10pt,right=10pt,top=10pt,bottom=8pt,before skip=10pt,after skip=12pt,halign=center,
  lower separated=false,halign lower=center,
  fontlower=\popsemi\fontsize{9}{11.5}\selectfont\color{popmuted},
  #1}
\newcommand{\legende}[1]{\par\vspace{6pt}{\popsemi\fontsize{9}{11.5}\selectfont\color{popmuted}#1}}

% ============================================================
% Signature OnBuch+
% ============================================================
\newcommand{\poplogoinline}[1]{{\poplogo\fontsize{#1}{#1}\selectfont\color{popink}OnBuch\color{poporange}+}}
\newcommand{\signatureonbuch}{%
  \begin{tcolorbox}[enhanced,colback=popink,colframe=popink,boxrule=2.2pt,arc=10pt,
      drop shadow=poporange,left=14pt,right=14pt,top=10pt,bottom=10pt,before skip=0pt,after skip=0pt]
    \tikz[baseline=-0.7ex]\node[circle,fill=popgreen,draw=popcream,line width=1.3pt,minimum size=22pt,inner sep=0]{\popcheck{white}};%
    \hspace{9pt}\begin{minipage}[c]{0.62\linewidth}
      {\popxbold\fontsize{12}{13}\selectfont\color{popcream}Signé\ \textbullet\ {\poplogo OnBuch\color{poporange}+}}\par\vspace{2pt}
      {\raggedright\popsemi\fontsize{8}{10.5}\selectfont\color{popline}\MakeUppercase{Équipe pédagogique OnBuch+}\\\MakeUppercase{Conforme au programme officiel MINESEC}\par}
    \end{minipage}\hfill
    {\poplogo\fontsize{22}{22}\selectfont\color{popcream}OnBuch\color{poporange}+}
  \end{tcolorbox}%
}

% ============================================================
% Page de garde
% #1 titre · #2 matière · #3 niveau · #4 module · #5 n° leçon · #6 durée ·
% #7 description / objectifs (texte ou itemize)
% ============================================================
\newcommand{\pagegarde}[7]{%
  \thispagestyle{empty}
  \newgeometry{a4paper,margin=1.7cm,top=1.6cm,bottom=1.4cm}%
  \begin{tikzpicture}[remember picture,overlay]
    \fill[poporange!18] ([xshift=-3.2cm,yshift=-3.6cm]current page.north east) circle (4.6cm);
    \fill[poppurple!14] ([xshift=2.4cm,yshift=7.6cm]current page.south west) circle (3.8cm);
  \end{tikzpicture}%
  {\poplogo\fontsize{30}{30}\selectfont\color{popink}OnBuch\color{poporange}+}\hfill
  \raisebox{0.6ex}{\poppill{Cours signé \textbullet\ Premium}{popink}{popcream}}\par
  \vspace{1pt}\popsquiggle{poppurple}\par
  \vspace{8pt}
  \begin{tcolorbox}[enhanced,colback=poporange,colframe=popink,boxrule=2.8pt,arc=13pt,
      drop shadow=popink,left=18pt,right=18pt,top=12pt,bottom=13pt,after skip=10pt]
    \poppill{#2 \textbullet\ #3}{popink}{popcream}\hspace{5pt}\poppill{#4}{white}{popink}\par\vspace{8pt}
    {\popdisplay\fontsize{14}{14}\selectfont\color{popink}LEÇON #5}\par\vspace{4pt}
    {\raggedright\hyphenpenalty=10000\exhyphenpenalty=10000\popdisplay\fontsize{25}{27}\selectfont\color{popcream}\MakeUppercase{#1}\par}
    \vspace{8pt}
    {\popxbold\fontsize{9.5}{9.5}\selectfont\color{popink}\MakeUppercase{Durée conseillée : #6}}
  \end{tcolorbox}
  \begin{tcolorbox}[enhanced,colback=white,colframe=popink,boxrule=2.2pt,arc=10pt,
      drop shadow=popink,left=16pt,right=16pt,top=10pt,bottom=10pt,after skip=8pt]
    \poppill{Dans cette leçon}{poppurple}{white}\par\vspace{5pt}
    \popsemi\fontsize{9.3}{12.6}\selectfont\color{popink}#7
  \end{tcolorbox}
  \noindent\begin{minipage}[t]{0.487\linewidth}
    \begin{tcolorbox}[enhanced,colback=popgreen,colframe=popink,boxrule=2.2pt,arc=10pt,
        drop shadow=popink,left=11pt,right=11pt,top=10pt,bottom=10pt,height=3.1cm,valign=center]
      \tcbox[colback=white,colframe=popink,boxrule=1.3pt,arc=5pt,boxsep=0pt,nobeforeafter,
        left=3pt,right=3pt,top=3pt,bottom=3pt,valign=center]{\qrcode[height=1.9cm]{https://play.google.com/store/apps/details?id=com.luvvix.onbuch&hl=fr}}%
      \hspace{8pt}\begin{minipage}[c]{0.5\linewidth}
        {\popdisplay\fontsize{11}{12.5}\selectfont\color{white}\MakeUppercase{Télécharge OnBuch}}\par\vspace{4pt}
        {\raggedright\popsemi\fontsize{8.4}{11}\selectfont\color{white}Gratuit \textbullet\ Google Play\\Tuteur IA, annales, concours, résultats\par}
      \end{minipage}
    \end{tcolorbox}
  \end{minipage}\hfill
  \begin{minipage}[t]{0.487\linewidth}
    \begin{tcolorbox}[enhanced,colback=poppurple,colframe=popink,boxrule=2.2pt,arc=10pt,
        drop shadow=popink,left=11pt,right=11pt,top=10pt,bottom=10pt,height=3.1cm,valign=center]
      \tikz[baseline=-0.7ex]\node[circle,fill=white,draw=popink,line width=1.3pt,minimum size=1.6cm,inner sep=0]{\popdisplay\fontsize{24}{24}\selectfont\color{poppurple}+};%
      \hspace{8pt}\begin{minipage}[c]{0.6\linewidth}
        {\popdisplay\fontsize{11}{12.5}\selectfont\color{white}\MakeUppercase{Passe à OnBuch+}}\par\vspace{4pt}
        {\raggedright\popsemi\fontsize{8.4}{11}\selectfont\color{white}Tous les cours signés, Tuteur IA illimité, fiches PDF — onglet OnBuch+ de l'app\par}
      \end{minipage}
    \end{tcolorbox}
  \end{minipage}
  \vspace{8pt}
  \popcontenu
  \vfill
  \signatureonbuch
  \restoregeometry
  \newpage
}

% Rangée « Ce cours contient » (page de garde)
\newcommand{\poptile}[3]{% #1 couleur #2 gros texte #3 légende
  \begin{tcolorbox}[enhanced,colback=#1,colframe=popink,boxrule=2pt,arc=9pt,shadow={2.5pt}{-2.5pt}{0pt}{popink},
      left=6pt,right=6pt,top=9pt,bottom=9pt,halign=center,valign=center,height=1.75cm,width=\linewidth,nobeforeafter]
    {\popdisplay\fontsize{12.5}{13}\selectfont\color{white}\MakeUppercase{#2}}\par\vspace{3pt}
    {\centering\popsemi\fontsize{7.8}{9.5}\selectfont\color{white}#3\par}
  \end{tcolorbox}}
\newcommand{\popcontenu}{%
  \par\noindent{\popxboldls\fontsize{7.5}{7.5}\selectfont\color{popmuted}CE COURS SIGNÉ CONTIENT}\par\vspace{7pt}
  \noindent\begin{minipage}[t]{0.235\linewidth}\poptile{popblue}{Cours}{complet, illustré, pas à pas}\end{minipage}\hfill
  \begin{minipage}[t]{0.235\linewidth}\poptile{poporange}{Méthodes}{et exercices résolus}\end{minipage}\hfill
  \begin{minipage}[t]{0.235\linewidth}\poptile{poppink}{Exercices}{type examen + corrigés}\end{minipage}\hfill
  \begin{minipage}[t]{0.235\linewidth}\poptile{popgreen}{Fiche bilan}{l'essentiel à retenir}\end{minipage}\par}

% Sommaire
\newtcolorbox{sommaire}{enhanced,colback=white,colframe=popink,boxrule=2.2pt,arc=10pt,
  drop shadow=popink,left=18pt,right=18pt,top=13pt,bottom=13pt,before skip=6pt,after skip=20pt,
  before upper={\poppill{Sommaire}{poppurple}{white}\par\vspace{9pt}\popsemi\fontsize{10.6}{14.5}\selectfont\color{popink}}}

% Signature de fin de cours — poussée tout en bas de la dernière page
\newcommand{\findecours}{%
  \par\vfill\nopagebreak
  \begin{center}\popsquiggle{poporange}\end{center}\vspace{4pt}
  \signatureonbuch
}

%% FIGFIX-BEGIN v3 — correctifs globaux des figures (docs/onbuch-generation/tools/figfix)
\usepackage{adjustbox}\usepackage{etoolbox}
% toute figure TikZ/pgfplots plus large que la ligne est réduite à la largeur disponible
\BeforeBeginEnvironment{tikzpicture}{\begin{adjustbox}{max width=\linewidth}}
\AfterEndEnvironment{tikzpicture}{\end{adjustbox}}
% légendes pgfplots lisibles par-dessus les courbes
\pgfplotsset{every axis/.append style={legend style={fill=white,fill opacity=0.92,text opacity=1,draw=black!15,inner sep=3pt}}}
% étiquettes d'arêtes (syntaxe edge["..."]) avec fond blanc
\tikzset{every edge quotes/.append style={fill=white,inner sep=1.2pt}}
% bulles de cartes mentales : texte à la ligne dans la bulle, bulle un peu plus grande
\tikzset{every concept/.append style={text width=2.0cm,align=center,minimum size=2.3cm,inner sep=2pt}}
%% FIGFIX-END
\begin{document}

\pagegarde{Les acides carboxyliques et leurs dérivés}{Chimie}{Tle C}{Module 1 · Chimie organique}{02}{5 h (cours + TP)}{Cette leçon étudie les acides carboxyliques : structure, nomenclature, propriétés physiques et acides, puis leurs dérivés (chlorures d'acyle, anhydrides, esters, amides). Elle aboutit aux applications industrielles majeures : saponification, synthèse des esters, et polycondensations conduisant au nylon 6,6 et au tergal.
\vspace{4pt}
\begin{multicols}{2}\small
\begin{itemize}
\item À la fin de cette leçon, je sais écrire la formule générale d'un acide carboxylique et décrire la structure du groupe carboxyle
\item À la fin de cette leçon, je sais nommer les monoacides saturés et insaturés, les polyacides et leurs dérivés (chlorures d'acyle, anhydrides, esters, amides)
\item À la fin de cette leçon, je sais expliquer les propriétés physiques des acides carboxyliques par les liaisons hydrogène et la formation de dimères
\item À la fin de cette leçon, je sais écrire l'équation-bilan de l'action de l'eau sur un acide carboxylique et identifier le couple RCOOH/RCOO⁻
\item À la fin de cette leçon, je sais écrire les équations de formation des chlorures d'acyle (PCl5, SOCl2) et des anhydrides d'acide (P4O10)
\item À la fin de cette leçon, je sais écrire l'équation d'estérification directe et celles des synthèses d'esters à partir des chlorures d'acyle et des anhydrides, en comparant leurs caractéristiques
\end{itemize}\end{multicols}}

\begin{sommaire}
\begin{multicols}{2}
\begin{enumerate}
\item Structure et nomenclature des acides carboxyliques
\item Propriétés physiques des acides carboxyliques
\item Propriété acide et formation des chlorures d'acyle et anhydrides
\item Les esters : estérification et synthèses rapides
\item Hydrolyse, saponification et fabrication du savon
\item Amides et polycondensation : nylon 6,6 et tergal
\item Travaux pratiques
\item Méthodes et exercices résolus
\item Exercices
\item Corrigés des exercices
\item Fiche bilan
\end{enumerate}
\end{multicols}
\end{sommaire}

\begin{objectifs}
\begin{itemize}
\item À la fin de cette leçon, je sais écrire la formule générale d'un acide carboxylique et décrire la structure du groupe carboxyle
\item À la fin de cette leçon, je sais nommer les monoacides saturés et insaturés, les polyacides et leurs dérivés (chlorures d'acyle, anhydrides, esters, amides)
\item À la fin de cette leçon, je sais expliquer les propriétés physiques des acides carboxyliques par les liaisons hydrogène et la formation de dimères
\item À la fin de cette leçon, je sais écrire l'équation-bilan de l'action de l'eau sur un acide carboxylique et identifier le couple RCOOH/RCOO⁻
\item À la fin de cette leçon, je sais écrire les équations de formation des chlorures d'acyle (PCl5, SOCl2) et des anhydrides d'acide (P4O10)
\item À la fin de cette leçon, je sais écrire l'équation d'estérification directe et celles des synthèses d'esters à partir des chlorures d'acyle et des anhydrides, en comparant leurs caractéristiques
\item À la fin de cette leçon, je sais définir l'hydrolyse et la saponification d'un ester et décrire la fabrication du savon à partir d'un triglycéride
\item À la fin de cette leçon, je sais écrire les équations de polycondensation conduisant au nylon 6,6 et au tergal (PET)
\end{itemize}
\end{objectifs}

\begin{prerequis}
\begin{itemize}
\item Nomenclature des hydrocarbures et des composés à fonction oxygénée (alcools, aldéhydes, cétones) vue en Première
\item Notion de fonction organique et de formule semi-développée
\item Couples acide/base, pH et indicateurs colorés (leçon 01 de Terminale)
\item Réactions d'oxydation ménagée des alcools (Première)
\item Notion de réaction lente/rapide, totale/limitée (cinétique, début d'année)
\end{itemize}
\end{prerequis}

\newpage

\coursec{Structure et nomenclature des acides carboxyliques}

\courssub{Situation d'accroche : une même famille, mille usages}

À Douala comme à Yaoundé, le savon utilisé chaque matin, le tissu tergal des uniformes scolaires, le nylon des cordes de pêche et même l'aspirine vendue en pharmacie ont un point commun : ils proviennent tous d'une même famille de composés organiques, les \cle{acides carboxyliques} et leurs dérivés. Comment l'huile de palme, abondante au Cameroun, se transforme-t-elle en savon dans les savonneries artisanales de l'Ouest ? Comment deux petites molécules s'assemblent-elles en une fibre textile solide comme le nylon ? Cette leçon donne les clés chimiques de ces fabrications industrielles et artisanales.

\textbf{Problématique.} Qu'est-ce qui caractérise la structure d'un acide carboxylique, et comment nommer rigoureusement ces molécules et leurs dérivés pour pouvoir ensuite comprendre leurs réactions (saponification, estérification, polycondensation) ?

\courssub{Le groupe carboxyle : définition et structure}

\textbf{Intuition d'abord.} En chimie organique, une \emph{famille} de composés est définie par un \cle{groupe caractéristique} : un petit motif d'atomes, toujours identique, qui commande les propriétés de toute la molécule. Pour les alcools, c'était le groupe hydroxyle \ce{-OH}. Pour les acides carboxyliques, c'est un groupe un peu plus riche : un carbonyle \ce{C=O} et un hydroxyle \ce{-OH} portés par le \emph{même} atome de carbone. Ce groupe s'appelle le \cle{groupe carboxyle} et s'écrit \ce{-COOH}.

\begin{definition}[Acide carboxylique]
Un \cle{acide carboxylique} est un composé organique possédant le groupe caractéristique \cle{carboxyle} \ce{-COOH}. Sa formule générale est \ce{R-COOH}, où \ce{R} est un atome d'hydrogène, un groupe alkyle (\ce{-CnH2n+1}) ou un groupe aryle (comme le phényle \ce{-C6H5}). Le carbone du groupe carboxyle est appelé \cle{carbone fonctionnel}.
\end{definition}

\textbf{Structure du groupe carboxyle.} Le carbone fonctionnel est lié à trois atomes (le carbone ou l'hydrogène de \ce{R}, l'oxygène du carbonyle, l'oxygène de l'hydroxyle) par trois directions de liaisons : il est \cle{trigonal}, hybridé $sp^2$. Les trois directions forment entre elles des angles voisins de \SI{120}{\degree}, et le groupe carboxyle est \emph{plan}. Les deux atomes d'oxygène, très électronégatifs, attirent les doublets de liaison : la liaison \ce{C=O} et la liaison \ce{O-H} sont fortement \cle{polarisées}. Le carbone fonctionnel porte une charge partielle $\delta^+$ et les oxygènes des charges partielles $\delta^-$ ; l'hydrogène de l'hydroxyle porte lui-même une charge partielle $\delta^+$, ce qui explique à la fois le caractère acide de la fonction et sa capacité à former des liaisons hydrogène.

\begin{popfigure}
\begin{tikzpicture}[scale=1.1,>=Stealth]
  % Atomes
  \node[font=\small\bfseries,text=popdark] (C) at (0,0) {C};
  \node[font=\small\bfseries,text=popdark] (O1) at (120:1.4) {O};
  \node[font=\small\bfseries,text=popdark] (O2) at (60:1.4) {O};
  \node[font=\small\bfseries,text=popdark] (R) at (180:1.4) {R};
  \node[font=\small\bfseries,text=popdark] (H) at (0:2.4) {H};
  % Liaisons
  \draw[popline,very thick,double distance=2pt] (C) -- (O1) node[midway,above left,font=\scriptsize,text=poporange]{$\delta^-$};
  \draw[popline,very thick] (C) -- (O2) -- (H);
  \draw[popline,very thick] (C) -- (R);
  % Angles
  \draw[poporange,thick,<->] (180:0.6) arc[start angle=180,end angle=120,radius=0.6] node[midway,left,font=\scriptsize,text=poporange]{$\approx\SI{120}{\degree}$};
  \draw[poporange,thick,<->] (120:0.6) arc[start angle=120,end angle=60,radius=0.6] node[midway,above,font=\scriptsize,text=poporange]{$\approx\SI{120}{\degree}$};
  \draw[poporange,thick,<->] (60:0.6) arc[start angle=60,end angle=0,radius=0.6] node[midway,right,font=\scriptsize,text=poporange]{$\approx\SI{120}{\degree}$};
  % Charges partielles
  \node[font=\scriptsize,text=popink] at (0.15,-0.35) {$\delta^+$};
  \node[font=\scriptsize,text=popblue] at ($(O1)+(0.3,0.25)$) {$\delta^-$};
  \node[font=\scriptsize,text=popblue] at ($(O2)+(0.3,0.25)$) {$\delta^-$};
  \node[font=\scriptsize,text=popink] at ($(H)+(0.35,0.1)$) {$\delta^+$};
  % Polarisation labels
  \draw[popink,thick,->] ($(C)!0.5!(O1)$) -- ++(0.5,0.5) node[right,font=\small\itshape,text=popink]{Liaison \ce{C=O} polarisée};
  \draw[popink,thick,->] ($(C)!0.5!(O2)$) -- ++(-0.5,0.5) node[left,font=\small\itshape,text=popink]{Liaison \ce{O-H} polarisée};
\end{tikzpicture}
\legende{Le groupe carboxyle : carbone trigonal $sp^2$, géométrie plane (angles voisins de \SI{120}{\degree}), liaisons \ce{C=O} et \ce{O-H} fortement polarisées.}
\end{popfigure}

\begin{attention}[Le groupe carboxyle n'est ni un carbonyle ni un alcool]
Ne confonds pas le groupe carboxyle \ce{-COOH} avec un simple carbonyle \ce{C=O} (aldéhydes, cétones) ou un simple hydroxyle \ce{-OH} (alcools). Dans \ce{-COOH}, les deux motifs sont portés par le \emph{même} carbone et se modifient mutuellement : c'est une \textbf{nouvelle fonction}, avec des propriétés propres (acidité notamment), qui ne sont ni celles d'un alcool ni celles d'une cétone.
\end{attention}

\courssub{Formule générale des monoacides saturés}

Pour les \cle{monoacides carboxyliques saturés acycliques} (une seule fonction \ce{-COOH}, chaîne carbonée sans double liaison ni cycle), le groupe alkyle a pour formule \ce{CnH2n+1} avec $n \geq 0$. On en déduit deux écritures équivalentes de la formule générale :

\begin{aretenir}[Formules générales à connaître]
\begin{itemize}
  \item Formule en fonction du nombre $n$ d'atomes de carbone de la chaîne alkyle : \ce{CnH2n+1-COOH} ($n \geq 0$) ;
  \item Formule brute globale, en comptant le carbone fonctionnel : \ce{C_nH_{2n}O_2} ($n \geq 1$), où $n$ est le nombre total d'atomes de carbone.
\end{itemize}
\end{aretenir}

\begin{exemplebox}[Vérification sur l'acide éthanoïque]
L'acide éthanoïque (vinaigre) a pour formule semi-développée \ce{CH3-COOH}. Sa chaîne alkyle est \ce{CH3-}, soit $n = 1$ dans \ce{CnH2n+1COOH}. Sa formule brute est \ce{C2H4O2} : on vérifie bien \ce{C_nH_{2n}O_2} avec $n = 2$ atomes de carbone au total ($2 \times 2 = 4$ hydrogènes).
\end{exemplebox}

\courssub{Nomenclature des acides carboxyliques}

\textbf{Intuition.} Nommer une molécule organique, c'est donner son adresse : la longueur de la rue (la chaîne carbonée), le numéro des habitants (les ramifications) et le type de maison (la fonction). Pour les acides carboxyliques, la maison est toujours au bout de la rue : le carbone du carboxyle est nécessairement en extrémité de chaîne.

\begin{methode}[Nommer un acide carboxylique]
\begin{enumerate}
  \item \textbf{Repérer} la chaîne carbonée la plus longue contenant le carbone du groupe carboxyle \ce{-COOH} ;
  \item \textbf{Numéroter} en partant obligatoirement du carbone fonctionnel, qui porte toujours le numéro 1 ;
  \item \textbf{Nommer} la chaîne avec le suffixe \og -oïque \fg{} précédé du préfixe \og acide \fg{} : acide + (racine de l'alcane) + oïque ;
  \item \textbf{Indiquer}, s'il y a lieu, les ramifications avec leur position, et la position des doubles liaisons éventuelles.
\end{enumerate}
\end{methode}

Ainsi, l'alcane à 4 carbones est le butane : l'acide correspondant est l'\cle{acide butanoïque}. Le \og e \fg{} final de l'alcane est remplacé par le suffixe \og -oïque \fg.

\begin{attention}[Erreur fréquente : oublier le carbone fonctionnel dans la chaîne]
Le carbone du groupe \ce{-COOH} \textbf{fait partie de la chaîne principale} et porte \textbf{toujours le numéro 1}. Ainsi \ce{CH3-CH2-COOH} compte \emph{trois} carbones : c'est l'acide \textbf{propanoïque}, et non un dérivé de l'éthane. De même, on ne précise jamais la position du groupe carboxyle : écrire \og acide butan-1-oïque \fg{} est inutile, le carboxyle est forcément en position 1.
\end{attention}

\begin{popfigure}
\begin{center}
\begin{tabular}{|c|c|c|}
\hline
\rowcolor{popblueL} \textbf{$n$ total} & \textbf{Formule semi-développée} & \textbf{Nom officiel (nom usuel)} \\
\hline
1 & \chemfig{H-C(=[1]O)-[7]OH} & acide méthanoïque (acide formique) \\
\hline
2 & \chemfig{CH_3-C(=[1]O)-[7]OH} & acide éthanoïque (acide acétique) \\
\hline
3 & \chemfig{CH_3-CH_2-C(=[1]O)-[7]OH} & acide propanoïque \\
\hline
4 & \chemfig{CH_3-CH_2-CH_2-C(=[1]O)-[7]OH} & acide butanoïque (acide butyrique) \\
\hline
5 & \chemfig{CH_3-(CH_2)_3-C(=[1]O)-[7]OH} & acide pentanoïque \\
\hline
6 & \chemfig{CH_3-(CH_2)_4-C(=[1]O)-[7]OH} & acide hexanoïque \\
\hline
\end{tabular}
\end{center}
\legende{Les six premiers monoacides carboxyliques saturés acycliques.}
\end{popfigure}

\courssub{Acides insaturés et polyacides}

\textbf{Acides insaturés.} Si la chaîne comporte une ou plusieurs doubles liaisons \ce{C=C}, le suffixe devient \og -ènoïque \fg{} et la position de la double liaison est indiquée par le plus petit indice possible, la numérotation partant toujours du carboxyle. Exemple célèbre : l'\cle{acide prop-2-ènoïque} \ce{CH2=CH-COOH}, appelé couramment \emph{acide acrylique}, matière première des peintures acryliques. L'\cle{acide oléique} \ce{C17H33COOH}, constituant majeur de l'huile d'olive et présent dans l'huile de palme, possède 18 carbones et une double liaison en position 9 : c'est l'acide octadéc-9-ènoïque.

\begin{exemplebox}[Nommer un acide insaturé]
Nommons \ce{CH3-CH=CH-COOH}.
\begin{itemize}
  \item Chaîne contenant le carboxyle : 4 carbones $\Rightarrow$ racine \og but \fg{} ;
  \item Numérotation depuis le carboxyle : la double liaison part du carbone 2 ;
  \item Nom : \textbf{acide but-2-ènoïque}.
\end{itemize}
\end{exemplebox}

\textbf{Polyacides.} Un \cle{polyacide} possède plusieurs groupes carboxyles. Les plus importants sont les \cle{diacides}, nommés avec le suffixe \og -dioïque \fg{} : l'\cle{acide éthanedioïque} \ce{HOOC-COOH} (acide oxalique, présent dans les feuilles d'oseille), l'\cle{acide hexanedioïque} \ce{HOOC-(CH2)4-COOH} (acide adipique, monomère du nylon 6,6) et l'\cle{acide benzène-1,4-dicarboxylique} (acide téréphtalique, monomère du tergal). Retiens bien ces deux derniers : ils reviendront dans la section 6 sur la polycondensation.

\begin{popfigure}
\begin{center}
\chemfig{CH_3-[:30](CH_2)_7-[:30]CH=[::-60]CH-[:30](CH_2)_7-[:30]C(=[2]O)-[:-30]OH}
\hspace{1.5cm}
\chemfig{HO-[:30]C(=[2]O)-[:-30](CH_2)_4-[:30]C(=[2]O)-[:-30]OH}
\end{center}
\legende{À gauche : l'acide oléique (acide octadéc-9-ènoïque), double liaison entre C9 et C10 en comptant depuis le carboxyle. À droite : l'acide adipique (acide hexanedioïque), diacide à 6 carbones, monomère du nylon 6,6.}
\end{popfigure}

\courssub{Noms usuels à connaître}

De nombreux acides carboxyliques sont connus depuis des siècles et conservent des noms usuels liés à leur origine naturelle. Le programme exige de savoir faire le lien entre nom usuel et nom officiel.

\begin{center}
\begin{tabularx}{\linewidth}{|l|X|l|}
\hline
\rowcolor{popgreenL} \textbf{Nom usuel} & \textbf{Formule} & \textbf{Nom officiel / origine} \\
\hline
Acide formique & \ce{HCOOH} & acide méthanoïque (fourmis) \\
\hline
Acide acétique & \ce{CH3COOH} & acide éthanoïque (vinaigre) \\
\hline
Acide butyrique & \ce{C3H7COOH} & acide butanoïque (beurre rance) \\
\hline
Acide oxalique & \ce{HOOC-COOH} & acide éthanedioïque (oseille) \\
\hline
Acide lactique & \ce{CH3-CH(OH)-COOH} & acide 2-hydroxypropanoïque (lait caillé, muscle) \\
\hline
Acide citrique & \ce{C6H8O7} & triacide des agrumes (citron) \\
\hline
Acide benzoïque & \ce{C6H5-COOH} & conservateur alimentaire (E210) \\
\hline
\end{tabularx}
\end{center}

\begin{savaistu}[Des fourmis au vinaigre]
L'\cle{acide méthanoïque} doit son nom usuel, \emph{acide formique} (du latin \emph{formica}, fourmi), aux piqûres de fourmis rousses qui en injectent une solution : c'est lui le responsable de la sensation de brûlure. L'\cle{acide éthanoïque}, ou acide acétique (du latin \emph{acetum}, vinaigre), donne son acidité au vinaigre, qui en contient environ \SI{6}{\percent} en masse. Au Cameroun, la fermentation acétique transforme naturellement le vin de palme laissé à l'air en vinaigre : les bactéries du genre \emph{Acetobacter} oxydent l'éthanol en acide éthanoïque.
\end{savaistu}

\begin{exoresolu}[Nommer et écrire des acides carboxyliques]
\textbf{Énoncé.} 1) Donner le nom officiel de l'acide de formule semi-développée \ce{CH3-CH(CH3)-CH2-COOH}. 2) Écrire la formule semi-développée de l'acide 2-méthylpropanoïque et vérifier que sa formule brute suit bien \ce{C_nH_{2n}O_2}.
\tcblower
\textbf{Solution détaillée.}

1) La chaîne la plus longue contenant le carboxyle compte 4 carbones (ne pas oublier le carbone fonctionnel, numéro 1). Un groupe méthyle est fixé sur le carbone 3. Le nom est donc : \textbf{acide 3-méthylbutanoïque}.

2) L'acide 2-méthylpropanoïque possède une chaîne de 3 carbones avec un méthyle en position 2 :
\[ \chemfig{CH_3-CH(-[2]CH_3)-C(=[1]O)-[7]OH} \]
Formule brute : \ce{C4H8O2}. Vérification : $n = 4$ atomes de carbone au total, $2n = 8$ hydrogènes, 2 oxygènes. La formule \ce{C_nH_{2n}O_2} est bien respectée.

\begin{attention}[Vigilance sur le comptage des carbones]
L'acide 2-méthylpropanoïque (\ce{C4H8O2}) est isomère de l'acide butanoïque. L'acide 3-méthylbutanoïque (\ce{C5H10O2}) est isomère de l'acide pentanoïque. Compte toujours le carbone fonctionnel comme le carbone n°1 de la chaîne principale.
\end{attention}
\end{exoresolu}

\begin{exercice}[À toi de jouer]
1) Écris la formule semi-développée de l'acide 2,2-diméthylpropanoïque. 2) Nomme l'acide \ce{HOOC-CH2-CH2-COOH}. 3) L'acide sorbique, conservateur alimentaire, a pour formule \ce{CH3-CH=CH-CH=CH-COOH} : donne son nom officiel.
\end{exercice}

\begin{corrige}[Exercice]
1) \chemfig{H_3C-C(-[2]CH_3)(-[6]CH_3)-C(=[1]O)-[7]OH} : chaîne de 3 carbones, deux méthyles sur C2. 2) Chaîne de 4 carbones avec deux groupes carboxyles aux extrémités : \textbf{acide butanedioïque} (nom usuel : acide succinique). 3) Chaîne de 6 carbones, doubles liaisons en positions 2 et 4 : \textbf{acide hexa-2,4-diènoïque}.
\end{corrige}

\begin{aretenir}[L'essentiel de la section]
\begin{itemize}
  \item Un acide carboxylique est un composé \ce{R-COOH} ; le groupe carboxyle \ce{-COOH} est plan, avec un carbone trigonal $sp^2$ et des liaisons fortement polarisées ;
  \item Formule générale des monoacides saturés acycliques : \ce{CnH2n+1COOH}, soit \ce{C_nH_{2n}O_2} au total ;
  \item Nomenclature : chaîne la plus longue incluant le carboxyle, carbone fonctionnel toujours numéro 1, suffixe \og -oïque \fg{} (\og -ènoïque \fg{} si insaturation, \og -dioïque \fg{} pour les diacides) ;
  \item Noms usuels à connaître : formique, acétique, butyrique, oxalique, lactique, citrique, benzoïque.
\end{itemize}
\end{aretenir}

\coursec{Propriétés physiques des acides carboxyliques}

Dans la section précédente, nous avons appris à reconnaître et à nommer les acides carboxyliques grâce à leur groupe caractéristique \ce{-COOH}. Mais pourquoi le vinaigre (solution d'acide éthanoïque) est-il un liquide qui se mélange parfaitement à l'eau, alors que l'acide palmitique de l'huile de palme est un solide gras ? Pourquoi l'acide éthanoïque bout-il à une température plus élevée qu'un alcool de même masse ? Toutes ces propriétés physiques s'expliquent par un phénomène unique : les \cle{liaisons hydrogène} que le groupe carboxyle peut former. C'est l'objet de cette section.

\courssub{État physique des acides carboxyliques à température ambiante}

Observe une série d'acides carboxyliques dans des flacons au laboratoire : les premiers sont liquides, les derniers sont des solides blancs cireux. Ce n'est pas un hasard.

\begin{propriete}[État physique]
À température ambiante (environ \SI{25}{\celsius}) :
\begin{itemize}
\item les acides carboxyliques de \ce{C1} à \ce{C9} (acide méthanoïque \ce{HCOOH}, acide éthanoïque \ce{CH3COOH}, ..., acide nonanoïque) sont des \textbf{liquides} incolores, souvent visqueux ;
\item les acides carboxyliques à partir de \ce{C10} (acide décanoïque, acide palmitique \ce{C15H31COOH}, acide stéarique \ce{C17H35COOH}...) sont des \textbf{solides} blancs, d'aspect cireux.
\end{itemize}
\end{propriete}

L'interprétation est simple : quand la chaîne carbonée s'allonge, la masse molaire augmente, les interactions de van der Waals entre les longues chaînes hydrocarbonées se renforcent, et il faut de plus en plus d'énergie pour séparer les molécules. Les températures de fusion et d'ébullition augmentent donc régulièrement avec le nombre d'atomes de carbone. Les acides gras de l'huile de palme (acides palmitique et oléique, \ce{C16} et \ce{C18}) sont ainsi solides ou pâteux à la température d'une cuisine camerounaise le matin.

\begin{propriete}[Odeur]
Les acides carboxyliques à chaîne courte possèdent une odeur \textbf{piquante et désagréable} : l'acide éthanoïque donne l'odeur âcre du vinaigre ; l'acide butanoïque (ou acide butyrique) est responsable de l'odeur du beurre rance et de certaines sueurs. Les acides à longue chaîne, non volatils, sont pratiquement inodores.
\end{propriete}

\courssub{Les liaisons hydrogène : la clé de toutes les propriétés physiques}

\textbf{Intuition.} Le groupe carboxyle \ce{-COOH} possède un atome d'oxygène doublement lié (\ce{C=O}) et un groupe hydroxyle \ce{-OH}. L'atome d'oxygène est très électronégatif : il attire vers lui les électrons de la liaison \ce{O-H}, laissant l'hydrogène partiellement positif ($\delta^+$) et l'oxygène partiellement négatif ($\delta^-$). Deux molécules voisines vont donc s'attirer électrostatiquement, comme de petits aimants.

\begin{definition}[Liaison hydrogène]
Une \cle{liaison hydrogène} est une interaction attractive qui s'établit entre :
\begin{itemize}
\item un atome d'hydrogène lié par covalence à un atome très électronégatif (O ou N), cet hydrogène portant une charge partielle $\delta^+$ ;
\item et un \textbf{doublet non liant} d'un autre atome électronégatif (O ou N) porteur d'une charge partielle $\delta^-$.
\end{itemize}
On la représente par des pointillés : \ce{O-H ... O}. Son énergie (environ \SI{20}{kJ.mol^{-1}}) est intermédiaire entre les interactions de van der Waals et une liaison covalente.
\end{definition}

\begin{attention}[Liaison hydrogène $\neq$ liaison covalente \ce{O-H}]
Erreur très fréquente au Baccalauréat ! Ne confonds pas :
\begin{itemize}
\item la \textbf{liaison covalente} \ce{O-H}, \emph{intramoléculaire} (à l'intérieur de la molécule), forte (environ \SI{460}{kJ.mol^{-1}}), qui relie réellement l'atome H à l'atome O ;
\item la \textbf{liaison hydrogène} \ce{O-H ... O}, \emph{intermoléculaire} (entre deux molécules), environ 20 fois plus faible, qui associe les molécules entre elles sans créer de nouvelle espèce chimique.
\end{itemize}
Dire « la liaison hydrogène relie H et O dans la molécule » est une faute grave : la liaison hydrogène n'existe qu'\textbf{entre} molécules.
\end{attention}

\courssub{Des températures d'ébullition anormalement élevées : les dimères}

\textbf{Intuition.} Pour faire bouillir un liquide, il faut arracher les molécules les unes aux autres. Si elles sont « collées » par des liaisons hydrogène, il faut fournir davantage d'énergie, donc chauffer plus fort : la température d'ébullition augmente.

Or l'acide carboxylique fait mieux que l'alcool : chaque molécule possède \textbf{à la fois} un donneur de liaison hydrogène (le \ce{O-H}) et un excellent accepteur (le \ce{C=O}). Deux molécules s'associent alors par \textbf{deux} liaisons hydrogène simultanées, formant un \cle{dimère} cyclique à 8 atomes, si stable qu'il persiste même en phase vapeur.

\begin{popfigure}
\begin{center}
\begin{tikzpicture}[scale=1.1]
% Molécule 1 (en haut)
\node (M1) at (0,0) {\chemfig{CH_3-C(=[2]O)-[6]O-H}};
% Molécule 2 (en bas, inversée)
\node (M2) at (0,-2.8) {\chemfig{H-O-[2]C(=[6]O)-CH_3}};
% Liaisons hydrogène (pointillés bleus)
\draw[popblue, very thick, dashed] (M1.south west) -- (M2.north west);
\draw[popblue, very thick, dashed] (M1.south east) -- (M2.north east);
% Étiquettes
\node[popblue, font=\small] at (-1.8,-1.4) {liaison hydrogène};
\node[popblue, font=\small] at (1.8,-1.4) {liaison hydrogène};
% Cadre cycle 8 atomes
\draw[popink, thick, rounded corners=3pt] (-2.2,0.5) rectangle (2.2,-3.3);
\node[popink, font=\small] at (0,-3.7) {Cycle à 8 atomes (dimère)};
\end{tikzpicture}
\end{center}
\legende{Dimère de l'acide éthanoïque : deux molécules associées par deux liaisons hydrogène (pointillés bleus), formant un cycle stable à 8 atomes.}
\end{popfigure}

\begin{propriete}[Conséquence sur les températures d'ébullition]
Les acides carboxyliques ont des températures d'ébullition \textbf{plus élevées} que celles des alcools de masse molaire voisine, car le dimère se comporte comme une molécule de masse double. L'ordre général, à masse molaire comparable, est :
\[ \theta_{\text{éb}}(\text{acide}) > \theta_{\text{éb}}(\text{alcool}) > \theta_{\text{éb}}(\text{aldéhyde}) > \theta_{\text{éb}}(\text{alcane}) \]
Les aldéhydes (molécules polaires sans \ce{O-H}) n'ont que des interactions dipôle-dipôle ; les alcanes (apolaires) n'ont que des interactions de van der Waals.
\end{propriete}

\begin{exemplebox}[Comparaison chiffrée à masse molaire égale]
L'acide éthanoïque \ce{CH3COOH} et le propan-1-ol \ce{CH3CH2CH2OH} ont la \textbf{même masse molaire} $M = \SI{60}{g.mol^{-1}}$. Pourtant :
\[ \theta_{\text{éb}}(\ce{CH3COOH}) = \SI{118}{\celsius} \qquad \theta_{\text{éb}}(\text{propan-1-ol}) = \SI{97}{\celsius} \]
L'écart de \SI{21}{\celsius} s'explique par la formation des \textbf{dimères} : l'acide éthanoïque se vaporise sous forme d'agrégats de deux molécules ($M_{\text{apparente}} \approx \SI{120}{g.mol^{-1}}$), ce qui demande plus d'énergie. À titre de comparaison, le butane ($M = \SI{58}{g.mol^{-1}}$) bout à \SI{-0,5}{\celsius} et le propanal ($M = \SI{58}{g.mol^{-1}}$) à \SI{49}{\celsius} !
\end{exemplebox}

\begin{popfigure}
\begin{center}
\begin{tikzpicture}
\begin{axis}[
width=0.95\linewidth, height=7cm,
xlabel={Nombre d'atomes de carbone $n$},
ylabel={$\theta_{\text{éb}}$ (\si{\celsius})},
xmin=0.5, xmax=6.5, ymin=-170, ymax=140,
grid=both, grid style={popline!40},
legend style={at={(0.03,0.97)}, anchor=north west, font=\small, fill=white, fill opacity=0.8},
tick label style={font=\small}, label style={font=\small}
]
% Acides : C1 à C6
\addplot[popcurve, popink, mark=*, mark size=1.5pt] coordinates {
(1,101) (2,118) (3,141) (4,164) (5,186) (6,205)
};
\addlegendentry{Acides carboxyliques}
% Alcools : C1 à C6
\addplot[popcurve, poporange, mark=square*, mark size=1.5pt] coordinates {
(1,65) (2,78) (3,97) (4,118) (5,137) (6,157)
};
\addlegendentry{Alcools}
% Alcanes : C1 à C6 (valeurs corrigées)
\addplot[popcurve, popgreen, mark=triangle*, mark size=1.8pt] coordinates {
(1,-161.5) (2,-88.6) (3,-42.1) (4,-0.5) (5,36.1) (6,68.7)
};
\addlegendentry{Alcanes}
\end{axis}
\end{tikzpicture}
\end{center}
\legende{Températures d'ébullition des acides carboxyliques, alcools et alcanes en fonction du nombre d'atomes de carbone. À $n$ donné, l'acide bout toujours le plus haut. Notez l'échelle étendue pour inclure le méthane ($n=1$).}
\end{popfigure}

\begin{savaistu}[L'acide acétique glacial]
L'acide éthanoïque pur cristallise dès $\theta_{\text{fus}} = \SI{16,6}{\celsius}$. En hiver, dans les laboratoires européens non chauffés, le flacon d'acide éthanoïque se solidifie en cristaux ressemblant à de la glace : c'est pourquoi on l'appelle \emph{acide acétique glacial}. Au Cameroun, il reste liquide toute l'année ! Attention : c'est un liquide corrosif qui provoque de graves brûlures.
\end{savaistu}

\courssub{Solubilité dans l'eau}

\textbf{Intuition.} L'eau est elle-même une championne des liaisons hydrogène. Une molécule organique qui sait, elle aussi, former des liaisons hydrogène avec l'eau sera « acceptée » dans le réseau des molécules d'eau : elle sera soluble.

\begin{propriete}[Solubilité des acides carboxyliques dans l'eau]
\begin{itemize}
\item Les acides à chaîne courte (\ce{C1} à \ce{C4} : méthanoïque, éthanoïque, propanoïque, butanoïque) sont \textbf{miscibles à l'eau en toutes proportions} : leur groupe \ce{-COOH} forme des liaisons hydrogène avec les molécules d'eau (partie \emph{hydrophile}).
\item La solubilité \textbf{diminue rapidement} lorsque la chaîne carbonée s'allonge : la partie hydrocarbonée \ce{R-}, apolaire, est \emph{hydrophobe} et « chassée » par l'eau. L'acide pentanoïque est faiblement soluble ; l'acide hexanoïque l'est encore moins ; les acides gras (\ce{C12} et plus) sont pratiquement \textbf{insolubles} dans l'eau (mais solubles dans les solvants organiques comme l'éthanol ou l'éther).
\end{itemize}
\end{propriete}

\begin{popfigure}
\begin{center}
\begin{tikzpicture}[scale=1.1]
% Molécule d'acide éthanoïque
\node (acid) at (0,0) {\chemfig{CH_3-C(=[2]O)-[6]O-H}};
% Molécules d'eau autour
\node (w1) at (3.5,1.2) {\chemfig{H-[:-30]O-[:30]H}};
\node (w2) at (3.5,-1.2) {\chemfig{H-[:-30]O-[:30]H}};
\node (w3) at (-2.5,1.2) {\chemfig{H-[:-30]O-[:30]H}};
\node (w4) at (-2.5,-1.2) {\chemfig{H-[:-30]O-[:30]H}};
% Liaisons hydrogène acide-eau (pointillés bleus)
\draw[popblue, very thick, dashed] (acid.east) -- (w1.west);
\draw[popblue, very thick, dashed] (acid.south east) -- (w2.north west);
\draw[popblue, very thick, dashed] (acid.north west) -- (w3.south east);
\draw[popblue, very thick, dashed] (acid.west) -- (w4.east);
% Annotations
\node[popdark, font=\small] at (-1.2,0.8) {partie hydrophile};
\draw[popdark, ->, shorten >=2pt] (-0.8,0.5) -- (acid.north west);
\node[popmuted, font=\small] at (1.5,-0.8) {partie hydrophobe ($\ce{CH3}$)};
\draw[popmuted, ->, shorten >=2pt] (1.0,-0.5) -- (acid.south east);
\node[popblue, font=\small, align=center] at (3.5,-2.2) {liaisons hydrogène\\acide--eau};
\end{tikzpicture}
\end{center}
\legende{Solubilité des petits acides carboxyliques : le groupe carboxyle établit des liaisons hydrogène avec l'eau (accepteur via le \protect\ce{C=O}, donneur via le \protect\ce{O-H}), tandis que la courte chaîne \protect\ce{CH3} reste peu perturbante.}
\end{popfigure}

C'est l'équilibre entre les deux parties de la molécule qui décide : tant que la « tête » hydrophile \ce{-COOH} domine (petits acides), l'acide se mélange à l'eau ; quand la « queue » hydrophobe $R$ devient longue, c'est elle qui impose son comportement, et l'acide devient insoluble. Cette dualité hydrophile/hydrophobe sera décisive dans la section 5 pour comprendre l'action \emph{détergente} du savon !

\begin{methode}[Expliquer une propriété physique d'un acide carboxylique]
Face à une question du type « Expliquer la température d'ébullition élevée » ou « Justifier la solubilité », rédige en trois étapes :
\begin{enumerate}
\item \textbf{Repérer} dans la molécule les groupes concernés : le \ce{O-H} (donneur de liaison hydrogène) et le \ce{C=O} (accepteur, doublets de l'oxygène).
\item \textbf{Invoquer} les liaisons hydrogène intermoléculaires : entre molécules d'acide (dimères, pour $\theta_{\text{éb}}$) ou entre l'acide et l'eau (pour la solubilité).
\item \textbf{Conclure} : liaisons hydrogène nombreuses $\Rightarrow$ cohésion forte $\Rightarrow$ $\theta_{\text{éb}}$ élevée ; liaisons hydrogène avec l'eau + chaîne courte $\Rightarrow$ bonne solubilité ; chaîne longue hydrophobe $\Rightarrow$ solubilité faible.
\end{enumerate}
\end{methode}

\begin{exoresolu}[Justifier et comparer des températures d'ébullition]
\textbf{Énoncé.} On donne trois composés de masse molaire voisine : l'acide méthanoïque \ce{HCOOH} ($M = \SI{46}{g.mol^{-1}}$), l'éthanol \ce{CH3CH2OH} ($M = \SI{46}{g.mol^{-1}}$) et le propane \ce{C3H8} ($M = \SI{44}{g.mol^{-1}}$). Leurs températures d'ébullition, dans le désordre, sont : \SI{-42}{\celsius}, \SI{78}{\celsius} et \SI{101}{\celsius}. Attribuer à chaque composé sa température d'ébullition en justifiant.
\tcblower
\textbf{Solution détaillée.}
\begin{enumerate}
\item \textbf{Repérage des groupes.} L'acide méthanoïque possède un groupe carboxyle \ce{-COOH} : un \ce{O-H} donneur et un \ce{C=O} accepteur. L'éthanol possède un seul groupe \ce{-OH}. Le propane, alcane apolaire, ne possède aucun site polaire.
\item \textbf{Liaisons hydrogène.} L'acide méthanoïque forme des \textbf{dimères} maintenus par deux liaisons hydrogène : sa cohésion est la plus forte. L'éthanol ne forme qu'\textbf{une} liaison hydrogène par molécule en moyenne : cohésion intermédiaire. Le propane n'établit que de faibles interactions de van der Waals.
\item \textbf{Conclusion.} La cohésion décroît dans l'ordre acide $>$ alcool $>$ alcane, donc :
\[ \theta_{\text{éb}}(\ce{HCOOH}) = \SI{101}{\celsius} \quad ; \quad \theta_{\text{éb}}(\ce{C2H5OH}) = \SI{78}{\celsius} \quad ; \quad \theta_{\text{éb}}(\ce{C3H8}) = \SI{-42}{\celsius}. \]
\end{enumerate}
Remarque : bien que leurs masses molaires soient quasi identiques, l'écart entre l'acide et l'alcane atteint \SI{143}{\celsius} ! C'est toute la puissance des liaisons hydrogène.
\end{exoresolu}

\begin{aretenir}[L'essentiel de la section]
\begin{itemize}
\item Les acides \ce{C1} à \ce{C9} sont liquides à \SI{25}{\celsius} ; au-delà, ce sont des solides cireux. Les acides courts ont une odeur piquante (beurre rance : acide butanoïque).
\item Le groupe \ce{-COOH} permet des \textbf{liaisons hydrogène} : interaction attractive entre un H lié à O et un doublet d'un autre O.
\item Deux molécules d'acide s'associent en \textbf{dimère} cyclique à 8 atomes (deux liaisons hydrogène), même en phase vapeur $\Rightarrow$ températures d'ébullition très élevées : acide $>$ alcool $>$ aldéhyde $>$ alcane à masse voisine.
\item Les acides \ce{C1} à \ce{C4} sont miscibles à l'eau (liaisons hydrogène acide--eau) ; la solubilité chute quand la chaîne hydrophobe s'allonge.
\end{itemize}
\end{aretenir}

Ces liaisons hydrogène, qui rendent le groupe carboxyle si cohésif, sont aussi à l'origine de son comportement en solution aqueuse : l'hydrogène du groupe \ce{-OH}, rendu très mobile, peut être cédé à l'eau. Nous étudierons dans la section suivante cette \textbf{propriété acide}, puis la manière de transformer les acides carboxyliques en dérivés plus réactifs : chlorures d'acyle et anhydrides.

\coursec{Propriété acide et formation des chlorures d'acyle et anhydrides}

Dans la section précédente, nous avons vu que les liaisons hydrogène expliquent les températures d'ébullition élevées et la bonne solubilité dans l'eau des petits acides carboxyliques. Mais que se passe-t-il, chimiquement, lorsqu'un acide carboxylique se dissout dans l'eau ? Pourquoi le vinaigre pique-t-il la langue, et pourquoi peut-on fabriquer à partir de l'acide éthanoïque des composés bien plus réactifs, indispensables à la synthèse des médicaments ? C'est tout l'objet de cette section : comprendre d'abord le \cle{caractère acide} du groupe carboxyle, puis apprendre à transformer le groupe \ce{-OH} du carboxyle en groupes plus réactifs, le \ce{-Cl} des \cle{chlorures d'acyle} et le \ce{-O-CO-R} des \cle{anhydrides d'acide}.

\courssub{Action de l'eau : les acides carboxyliques sont des acides faibles}

\textbf{Intuition.} Verse un peu de vinaigre (solution d'acide éthanoïque) dans l'eau : le mélange reste acide, il fait virer le bleu de bromothymol au jaune. Pourtant, si l'acide éthanoïque se comportait comme l'acide chlorhydrique, une solution à \SI{0,10}{mol.L^{-1}} aurait un pH de 1. Or on mesure un pH voisin de 2,9 ! L'acide ne s'est donc pas totalement transformé : seule une petite fraction des molécules a cédé son proton à l'eau.

\begin{definition}[Acide carboxylique, acide faible]
Un acide carboxylique \ce{R-COOH} est un \cle{acide faible} : sa réaction avec l'eau est une réaction \emph{limitée} (un équilibre), conduisant à l'ion \cle{carboxylate} \ce{R-COO-} et à l'ion oxonium \ce{H3O+} :
\[ \ce{R-COOH + H2O <=> R-COO- + H3O+} \]
Le couple acide/base mis en jeu est le couple \cle{\ce{RCOOH}/\ce{RCOO-}}.
\end{definition}

L'ion carboxylate est la \emph{base conjuguée} de l'acide carboxylique. Sa stabilité relative explique pourquoi l'acide peut céder son proton : la charge négative portée par l'ion n'est pas localisée sur un seul atome d'oxygène, elle est \cle{délocalisée} sur les deux atomes d'oxygène du groupe carboxylate. On dit que l'ion est \emph{stabilisé par résonance} (mésomérie).

\begin{popfigure}
\begin{center}
\begin{tikzpicture}[node distance=1.2cm, every node/.style={anchor=center}]
\node (m1) {\chemfig{R-C(=[1]O)(-[7]O^{-})}};
\node[right=of m1] (eq1) {$\longleftrightarrow$};
\node[right=of eq1] (m2) {\chemfig{R-C(-[1]O^{-})(=[7]O)}};
\node[right=of m2] (eq2) {$\equiv$};
\node[right=of eq2] (m3) {\chemfig{R-C(-[1]O^{\scriptstyle\delta^{-}})(-[7]O^{\scriptstyle\delta^{-}})}};
\end{tikzpicture}
\end{center}
\legende{Formes mésomères de l'ion carboxylate \ce{R-COO-} : la charge négative est répartie équitablement sur les deux oxygènes (hybride de résonance à droite). Cette délocalisation stabilise l'ion et favorise la libération du proton \ce{H+}.}
\end{popfigure}

\begin{aretenir}[L'essentiel sur l'acidité]
\begin{itemize}
\item Tout acide carboxylique est un acide \textbf{faible} : écrire systématiquement la double flèche \ce{<=>} dans son équation avec l'eau.
\item Couple acide/base : \ce{RCOOH}/\ce{RCOO-}.
\item L'ion carboxylate est stabilisé par résonance : les deux liaisons C--O ont la même longueur, intermédiaire entre simple et double liaison.
\end{itemize}
\end{aretenir}

\begin{attention}[Piège classique au Baccalauréat]
Ne jamais écrire \ce{RCOOH + H2O -> RCOO- + H3O+} avec une flèche simple ! Une flèche simple signifie une réaction \emph{totale}, ce qui est faux pour un acide carboxylique. Seuls les acides forts (comme \ce{HCl}) s'ionisent totalement. L'oubli de la double flèche coûte des points à chaque session.
\end{attention}

\begin{exoresolu}[Le pH du vinaigre prouve l'ionisation partielle]
On dissout de l'acide éthanoïque \ce{CH3COOH} dans l'eau distillée de façon à obtenir une solution de concentration $C = \SI{0,10}{mol.L^{-1}}$. Le pH-mètre indique $\text{pH} = 2,9$.
\begin{enumerate}
\item Quel serait le pH si l'acide était fort ?
\item Calculer la concentration réelle en ions \ce{H3O+} et conclure.
\end{enumerate}
\tcblower
\textbf{1.} Si l'acide éthanoïque était un acide fort, il serait totalement ionisé : $[\ce{H3O+}] = C = \SI{0,10}{mol.L^{-1}}$, d'où $\text{pH} = -\log(0,10) = 1,0$.

\textbf{2.} La mesure donne :
\[ [\ce{H3O+}] = 10^{-\text{pH}} = 10^{-2,9} \approx \SI{1,3e-3}{mol.L^{-1}}. \]
Le taux d'ionisation est donc :
\[ \tau = \frac{[\ce{H3O+}]}{C} = \frac{1,3\times 10^{-3}}{0,10} \approx 1,3\,\%. \]
Seulement environ \textbf{1 molécule sur 80} a cédé son proton : l'ionisation est bien \emph{partielle}, l'acide éthanoïque est un \cle{acide faible}. L'équilibre \ce{CH3COOH + H2O <=> CH3COO- + H3O+} est très déplacé vers la gauche.
\end{exoresolu}

\courssub{Conséquences expérimentales du caractère acide}

Le caractère acide des acides carboxyliques se manifeste par plusieurs tests simples, réalisables au laboratoire du lycée.

\begin{itemize}
\item \textbf{pH et indicateurs colorés :} toute solution aqueuse d'acide carboxylique a un $\text{pH} < 7$ ; elle fait virer l'hélianthine au rouge-rose et le bleu de bromothymol au jaune.
\item \textbf{Action sur les bases :} un acide carboxylique réagit avec l'hydroxyde de sodium (soude) selon une réaction quasi totale (réaction acido-basique de titrage) :
\[ \ce{R-COOH + HO- -> R-COO- + H2O} \]
On obtient une solution de \emph{carboxylate de sodium} \ce{R-COO- Na+}. C'est exactement ce type de sel que l'on retrouve dans les savons (section 5) !
\item \textbf{Action sur les carbonates et hydrogénocarbonates : test caractéristique.} Au contact du carbonate de sodium \ce{Na2CO3} ou de l'hydrogénocarbonate de sodium \ce{NaHCO3}, un acide carboxylique provoque un \cle{dégagement gazeux} de dioxyde de carbone (effervescence) :
\[ \ce{2 R-COOH + Na2CO3 -> 2 R-COO- Na+ + H2O + CO2 ^} \]
Le gaz dégagé trouble l'eau de chaux : c'est bien \ce{CO2}. Ce test permet de distinguer un acide carboxylique d'un alcool ou d'un phénol.
\end{itemize}

\begin{experience}[Test au carbonate de sodium]
\textbf{Protocole :} introduire dans un tube à essais \SI{2}{mL} d'acide éthanoïque dilué, puis ajouter une pointe de spatule de carbonate de sodium solide. Boucher le tube avec un tuyau débouchant dans un second tube contenant de l'eau de chaux.

\textbf{Observations :} effervescence immédiate dans le premier tube ; l'eau de chaux se trouble dans le second.

\textbf{Interprétation :} le gaz dégagé est \ce{CO2}, formé selon \ce{2 CH3COOH + Na2CO3 -> 2 CH3COO- Na+ + H2O + CO2 ^}. L'effervescence est la signature expérimentale de la fonction acide carboxylique.
\end{experience}

\begin{popfigure}
\begin{center}
\begin{tikzpicture}[scale=0.9]
% Tube 1
\draw[popdark, thick] (0,0) -- (0,2.6);
\draw[popdark, thick] (1,0) -- (1,2.6);
\draw[popdark, thick] (0,0) arc (180:360:0.5);
\fill[popblueL] (0.03,0.15) -- (0.03,1.2) -- (0.97,1.2) -- (0.97,0.15) arc (180:360:0.47);
% bulles
\foreach \x/\y/\r in {0.25/0.4/0.06, 0.6/0.7/0.08, 0.8/0.35/0.05, 0.4/0.95/0.07, 0.7/1.05/0.05, 0.2/0.8/0.05}
  \draw[popblue, fill=white] (\x,\y) circle (\r);
\foreach \x/\y in {0.3/1.5, 0.55/1.8, 0.75/1.45, 0.45/2.1}
  \draw[popblue] (\x,\y) circle (0.05);
\node[below, align=center, font=\small] at (0.5,-0.35) {\ce{CH3COOH} + \ce{Na2CO3}\\ effervescence de \ce{CO2}};
% flèche
\draw[->, poporange, very thick] (1.4,1.3) -- (2.4,1.3);
% Tube 2
\draw[popdark, thick] (2.8,0) -- (2.8,2.6);
\draw[popdark, thick] (3.8,0) -- (3.8,2.6);
\draw[popdark, thick] (2.8,0) arc (180:360:0.5);
\fill[popgoldL, opacity=0.8] (2.83,0.15) -- (2.83,1.5) -- (3.77,1.5) -- (3.77,0.15) arc (180:360:0.47);
\foreach \x/\y in {3.0/0.5, 3.3/0.9, 3.6/0.4, 3.15/1.2, 3.5/1.1}
  \fill[popgold!60] (\x,\y) circle (0.04);
\node[below, align=center, font=\small] at (3.3,-0.35) {eau de chaux\\ qui se trouble};
\end{tikzpicture}
\end{center}
\legende{Test caractéristique des acides carboxyliques : l'effervescence de \ce{CO2} au contact d'un carbonate trouble l'eau de chaux.}
\end{popfigure}

\courssub{Nomenclature des ions carboxylates}

\begin{definition}[Nommer un ion carboxylate]
Le nom d'un ion carboxylate s'obtient à partir du nom de l'acide correspondant en remplaçant la terminaison \emph{-oïque} par le suffixe \cle{-oate}, précédé du mot « ion » :
\begin{itemize}
\item \ce{HCOO-} : ion \textbf{méthanoate} (ou formiate) ;
\item \ce{CH3COO-} : ion \textbf{éthanoate} (ou acétate) ;
\item \ce{C6H5COO-} : ion \textbf{benzoate} ;
\item \ce{CH3(CH2)16COO-} : ion \textbf{stéarate} (composant des savons).
\end{itemize}
Les sels correspondants se nomment de la même façon : \ce{CH3COO- Na+} est l'\textbf{éthanoate de sodium}.
\end{definition}

\courssub{Formation des chlorures d'acyle}

\textbf{Pourquoi transformer l'acide ?} Le groupe \ce{-OH} du carboxyle est un mauvais « partant » : les réactions de l'acide carboxylique (comme l'estérification, section 4) sont lentes et limitées. Pour obtenir des réactions \emph{rapides et totales}, on remplace \ce{-OH} par un atome de chlore, excellent groupe partant : on obtient un \cle{chlorure d'acyle}.

\begin{definition}[Chlorure d'acyle]
Un \cle{chlorure d'acyle} est un dérivé d'acide carboxylique de formule générale \ce{R-CO-Cl}, dans lequel le groupe hydroxyle \ce{-OH} du groupe carboxyle a été remplacé par un atome de chlore. Son nom s'obtient en remplaçant la terminaison \emph{-oïque} de l'acide par \cle{chlorure de ...oyle} :
\begin{itemize}
\item \ce{CH3-CO-Cl} : chlorure d'\textbf{éthanoyle} ;
\item \ce{CH3-CH2-CO-Cl} : chlorure de \textbf{propanoyle} ;
\item \ce{C6H5-CO-Cl} : chlorure de \textbf{benzoyle}.
\end{itemize}
\end{definition}

Deux réactifs chlorurants sont au programme : le pentachlorure de phosphore \ce{PCl5} et le chlorure de thionyle \ce{SOCl2}.

\textbf{Avec le pentachlorure de phosphore \ce{PCl5} :}
\[ \ce{R-COOH + PCl5 -> R-COCl + POCl3 + HCl} \]
La réaction est vive et \textbf{totale}. Sous-produits : l'oxychlorure de phosphore \ce{POCl3} (liquide) et le chlorure d'hydrogène \ce{HCl} (gaz).

\textbf{Avec le chlorure de thionyle \ce{SOCl2} :}
\[ \ce{R-COOH + SOCl2 -> R-COCl + SO2 ^ + HCl ^} \]
La réaction est également totale, et elle présente un avantage décisif : les deux sous-produits, \ce{SO2} et \ce{HCl}, sont \emph{gazeux} et s'échappent spontanément du milieu réactionnel. Le chlorure d'acyle est donc obtenu quasiment pur, sans purification délicate.

\begin{savaistu}[Pourquoi l'industrie préfère \ce{SOCl2}]
Dans l'industrie pharmaceutique, la pureté des intermédiaires de synthèse est cruciale. Avec \ce{PCl5}, le sous-produit \ce{POCl3} est un liquide difficile à séparer du chlorure d'acyle. Avec \ce{SOCl2}, \ce{SO2} et \ce{HCl} s'éliminent d'eux-mêmes sous forme gazeuse : on dit que la réaction est « propre ». C'est pourquoi \ce{SOCl2} est le réactif chlorurant de choix, par exemple dans la fabrication de l'aspirine à partir de l'acide salicylique.
\end{savaistu}

\courssub{Formation des anhydrides d'acide}

\begin{definition}[Anhydride d'acide]
Un \cle{anhydride d'acide} est un dérivé d'acide carboxylique de formule générale \ce{(R-CO)2O}, c'est-à-dire \ce{R-CO-O-CO-R}, obtenu par \emph{élimination d'une molécule d'eau entre deux molécules d'acide carboxylique}. Son nom s'obtient en remplaçant le mot « acide » par \cle{anhydride} et en gardant la terminaison \emph{-oïque} :
\begin{itemize}
\item \ce{(CH3-CO)2O} : anhydride \textbf{éthanoïque} ;
\item \ce{(CH3-CH2-CO)2O} : anhydride \textbf{propanoïque}.
\end{itemize}
Lorsque les deux groupes R sont différents, on parle d'\textbf{anhydride mixte}, nommé en citant les deux chaînes par ordre alphabétique (ex. : anhydride éthanoïque propanoïque).
\end{definition}

La déshydratation est réalisée à chaud en présence d'un agent déshydratant puissant, le décaoxyde de tétraphosphore \ce{P4O10} (ancien nom : anhydride phosphorique) :
\[ \ce{2 R-COOH ->[P4O10,\ \Delta] (R-CO)2O + H2O} \]
Le \ce{P4O10} « piège » l'eau formée, ce qui déplace l'équilibre vers l'anhydride. Pour l'acide éthanoïque :
\[ \ce{2 CH3-COOH ->[P4O10] (CH3-CO)2O + H2O} \]

\begin{popfigure}
\begin{center}
\begin{tikzpicture}[font=\small, node distance=0.35cm]
\node (a1) {\chemfig{CH_3-C(=[1]O)-[7]OH}};
\node[right=of a1] (p1) {$+$};
\node[right=of p1] (r1) {\ce{PCl5}};
\node[right=of r1] (f1) {$\longrightarrow$};
\node[right=of f1] (c1) {\chemfig{CH_3-C(=[1]O)-[7]Cl}};
\node[right=of c1] (p2) {$+$ \ce{POCl3} $+$ \ce{HCl}};
\node[below=0.5cm of a1] (a2) {\chemfig{CH_3-C(=[1]O)-[7]OH}};
\node[right=of a2] (p3) {$+$};
\node[right=of p3] (r2) {\ce{SOCl2}};
\node[right=of r2] (f2) {$\longrightarrow$};
\node[right=of f2] (c2) {\chemfig{CH_3-C(=[1]O)-[7]Cl}};
\node[right=of c2] (p4) {$+$ \ce{SO2} $\uparrow$ $+$ \ce{HCl} $\uparrow$};
\node[below=0.5cm of a2] (a3) {2 \chemfig{CH_3-C(=[1]O)-[7]OH}};
\node[right=0.55cm of a3] (f3) {$\xrightarrow{\ \ce{P4O10},\ \Delta\ }$};
\node[right=0.55cm of f3] (c3) {\chemfig{CH_3-C(=[1]O)-[7]O-C(=[1]O)-[7]CH_3}};
\node[right=of c3] (p5) {$+$ \ce{H2O}};
\end{tikzpicture}
\end{center}
\legende{Synthèses des dérivés de l'acide éthanoïque : chlorure d'éthanoyle par \ce{PCl5} (ligne 1) ou \ce{SOCl2} (ligne 2, sous-produits gazeux), et anhydride éthanoïque par déshydratation avec \ce{P4O10} (ligne 3).}
\end{popfigure}

\begin{methode}[Écrire l'équation d'une réaction de dérivation]
Pour écrire correctement l'équation de formation d'un dérivé d'acide carboxylique :
\begin{enumerate}
\item \textbf{Repérer} le groupe \ce{-OH} du carboxyle : c'est lui qui sera éliminé et remplacé.
\item \textbf{Identifier} le réactif : \ce{PCl5} ou \ce{SOCl2} conduisent au chlorure d'acyle (remplacement par \ce{-Cl}) ; \ce{P4O10} conduit à l'anhydride (deux acides se rejoignent par un pont \ce{-O-}).
\item \textbf{Écrire} le dérivé obtenu : \ce{R-CO-Cl} ou \ce{R-CO-O-CO-R}.
\item \textbf{Équilibrer} les sous-produits : \ce{POCl3 + HCl} avec \ce{PCl5} ; \ce{SO2 + HCl} avec \ce{SOCl2} ; \ce{H2O} avec \ce{P4O10} (penser au coefficient 2 devant l'acide !).
\item \textbf{Vérifier} la conservation des atomes des deux côtés de la flèche.
\end{enumerate}
\end{methode}

\begin{attention}[Erreurs fréquentes]
\begin{itemize}
\item Oublier le coefficient \textbf{2} devant l'acide dans la formation de l'anhydride : il faut \emph{deux} molécules d'acide pour une molécule d'anhydride.
\item Confondre les sous-produits : avec \ce{PCl5} on obtient \ce{POCl3} (et non \ce{PCl3}) ; avec \ce{SOCl2} on obtient \ce{SO2} (et non \ce{S} ou \ce{SO3}).
\item Écrire \ce{R-CO-OH-Cl} ou placer le chlore sur la chaîne carbonée : le chlore remplace \emph{uniquement} le \ce{-OH} du carboxyle, jamais un hydrogène de la chaîne.
\end{itemize}
\end{attention}

\begin{exoresolu}[Équations à compléter]
Compléter et équilibrer les équations suivantes, puis nommer tous les produits organiques :
\begin{enumerate}
\item \ce{CH3-CH2-COOH + SOCl2 -> ...}
\item \ce{C6H5-COOH + PCl5 -> ...}
\item \ce{2 CH3-CH2-COOH ->[P4O10] ...}
\end{enumerate}
\tcblower
\textbf{1.} \ce{CH3-CH2-COOH + SOCl2 -> CH3-CH2-COCl + SO2 ^ + HCl ^}. Le produit organique est le \textbf{chlorure de propanoyle}.

\textbf{2.} \ce{C6H5-COOH + PCl5 -> C6H5-COCl + POCl3 + HCl}. Le produit organique est le \textbf{chlorure de benzoyle}.

\textbf{3.} \ce{2 CH3-CH2-COOH ->[P4O10] (CH3-CH2-CO)2O + H2O}. Le produit organique est l'\textbf{anhydride propanoïque}.
\end{exoresolu}

\begin{aretenir}[Bilan de la section]
\begin{itemize}
\item Les acides carboxyliques sont des \textbf{acides faibles} : \ce{RCOOH + H2O <=> RCOO- + H3O+} ; l'ion carboxylate est stabilisé par résonance.
\item Tests du caractère acide : pH $< 7$, indicateurs colorés, réaction avec les bases, \textbf{effervescence de \ce{CO2}} avec les carbonates.
\item Ions carboxylates : suffixe \textbf{-oate} (éthanoate, benzoate).
\item \ce{RCOOH + PCl5 -> RCOCl + POCl3 + HCl} ; \ce{RCOOH + SOCl2 -> RCOCl + SO2 ^ + HCl ^} (réaction propre, sous-produits gazeux).
\item \ce{2 RCOOH ->[P4O10] (RCO)2O + H2O} : anhydride d'acide.
\item Nomenclature : chlorure de \textbf{...oyle}, anhydride \textbf{...oïque}.
\end{itemize}
\end{aretenir}

Ces deux familles de dérivés — chlorures d'acyle et anhydrides — ne sont pas de simples curiosités de laboratoire : ce sont des \emph{intermédiaires réactionnels} d'une grande efficacité. Nous verrons dans la prochaine section qu'ils permettent de préparer les esters par des réactions \textbf{rapides et totales}, là où l'estérification directe entre un acide et un alcool reste lente et limitée.

\coursec{Les esters : estérification et synthèses rapides}

Dans la section précédente, nous avons vu que le groupe carboxyle \ce{-COOH} confère aux acides carboxyliques une \cle{propriété acide} (couple \ce{RCOOH}/\ce{RCOO-}), et que l'on peut transformer un acide en dérivés beaucoup plus réactifs : les \cle{chlorures d'acyle} \ce{RCOCl} et les \cle{anhydrides d'acide} \ce{(RCO)2O}. À quoi servent ces dérivés ? Principalement à fabriquer des \cle{esters}, ces molécules aux odeurs fruitées que l'on retrouve dans la banane, l'ananas, les parfums... et même dans l'aspirine ! C'est l'objet de cette section : comprendre ce qu'est un ester, comment on le prépare \emph{lentement} à partir d'un acide, et comment on le prépare \emph{rapidement et totalement} grâce aux dérivés d'acide.

\courssub{Définition et nomenclature des esters}

\begin{definition}[Ester et estérification]
Un \cle{ester} est un composé organique oxygéné de formule générale \ce{R-COO-R$'$}, où \ce{R} est un atome d'hydrogène ou un groupe alkyle, et \ce{R$'$} un groupe alkyle. Le groupe caractéristique \ce{-COO-} est appelé \cle{groupe ester}.

On appelle \cle{estérification} toute réaction entre un acide carboxylique (ou l'un de ses dérivés : chlorure d'acyle, anhydride) et un alcool, conduisant à un ester et à de l'eau (ou à \ce{HCl}, ou à un acide carboxylique).
\end{definition}

La structure de l'ester révèle sa « double origine » : la partie \ce{R-CO-} (appelée groupe \cle{acyle}) provient de l'acide carboxylique, tandis que le groupe \ce{-O-R$'$} provient de l'alcool. C'est précisément cette double origine qui commande la nomenclature.

\begin{methode}[Nommer un ester : « alkanoate d'alkyle »]
\begin{enumerate}
\item \textbf{Identifier la partie acyle} \ce{R-CO-} issue de l'acide : compter ses carbones (y compris le carbone du \ce{C=O}), nommer l'alcane correspondant, remplacer la terminaison « -ane » par « \textbf{-anoate} ».
\item \textbf{Identifier le groupe alkyle} \ce{R$'$} fixé sur l'oxygène, issu de l'alcool : le nommer comme un alkyle (méthyle, éthyle, propyle...).
\item \textbf{Assembler} : le nom de l'ester est « alkanoate \textbf{d'}alkyle », dans cet ordre : d'abord la partie acide, ensuite la partie alcool.
\end{enumerate}
\end{methode}

\begin{exemplebox}[Deux esters à nommer]
\begin{itemize}
\item \chemfig{CH_3-C(=[1]O)-[7]O-CH_2-CH_3} : la partie acyle \ce{CH3-CO-} a 2 carbones $\to$ « éthanoate » ; le groupe sur l'oxygène est \ce{-CH2-CH3} $\to$ « d'éthyle ». Nom : \textbf{éthanoate d'éthyle}.
\item \chemfig{H-C(=[1]O)-[7]O-CH_2-CH_2-CH_3} : la partie acyle \ce{H-CO-} a 1 carbone $\to$ « méthanoate » ; le groupe alkyle est \ce{-CH2CH2CH3} $\to$ « de propyle ». Nom : \textbf{méthanoate de propyle}.
\end{itemize}
\end{exemplebox}

\begin{attention}[Erreur fréquente : inverser les deux parties de l'ester]
L'\textbf{éthanoate d'éthyle} s'écrit \ce{CH3-COO-CH2-CH3} : le groupe \textbf{éthyle} est du côté de l'\textbf{oxygène}, et le groupe \textbf{éthanoate} contient le \ce{C=O}. Écrire \ce{C2H5-COO-CH3} serait faux : cette formule correspond au \emph{propanoate de méthyle}, un ester isomère mais différent ! Retiens : \og le nom se lit dans le même sens que la molécule, en partant du \ce{C=O} \fg.
\end{attention}

\courssub{L'estérification directe : une réaction lente et limitée}

\begin{experience}[Réaction entre l'acide éthanoïque et l'éthanol]
\textbf{Dispositif et protocole :} dans un tube à essai surmonté d'un réfrigérant, on mélange \SI{1}{mol} d'acide éthanoïque, \SI{1}{mol} d'éthanol et quelques gouttes d'acide sulfurique concentré. On chauffe à reflux et on prélève régulièrement des échantillons que l'on dose.

\textbf{Observations :} une odeur fruitée apparaît progressivement (ester formé), mais le dosage montre que la réaction \emph{semble s'arrêter} alors qu'il reste de l'acide et de l'alcool : au bout de plusieurs heures, on n'obtient qu'environ \SI{0,67}{mol} d'ester.

\textbf{Interprétation :} la réaction est \cle{lente} et \cle{limitée} par la réaction inverse (l'hydrolyse de l'ester) : c'est un \cle{équilibre chimique}.
\end{experience}

L'équation-bilan générale de l'estérification directe s'écrit :

\[ \ce{RCOOH + R$'$OH <=> RCOO-R$'$ + H2O} \]

\begin{popfigure}
\begin{center}
\chemfig{[:-30]R-C(=[2]@{c1}\color{popblue}O)-[6]@{o1}\color{popblue}OH}
\hspace{0.6em}{\Large +}\hspace{0.6em}
\chemfig{[:-30]@{h1}\color{poporange}H-[6]@{o2}\color{poporange}O-R'}
\hspace{0.6em}{\Large $\rightleftharpoons$}\hspace{0.6em}
\chemfig{[:-30]R-C(=[2]O)-[6]O-R'}
\hspace{0.6em}{\Large +}\hspace{0.6em}
\chemfig{H_2O}
\chemmove{
\draw[popblue,thick,rounded corners] ([xshift=-2pt,yshift=-2pt]c1.south west) rectangle ([xshift=2pt,yshift=2pt]o1.north east);
\draw[poporange,thick,rounded corners] ([xshift=-2pt,yshift=-2pt]h1.south west) rectangle ([xshift=2pt,yshift=2pt]o2.north east);
}
\end{center}
\legende{Équation générale de l'estérification directe. En bleu : le groupe \ce{-OH} cédé par l'acide ; en orange : l'hydrogène \ce{-H} cédé par l'alcool. Ces deux fragments s'éliminent ensemble pour former l'eau \ce{H2O}.}
\end{popfigure}

\begin{propriete}[Caractéristiques de l'estérification directe]
La réaction entre un acide carboxylique et un alcool est :
\begin{itemize}
\item \textbf{lente} : l'équilibre n'est atteint qu'au bout de plusieurs heures, voire plusieurs jours à température ambiante ;
\item \textbf{limitée} (équilibre) : la réaction inverse, l'\cle{hydrolyse} de l'ester, se produit simultanément ; on n'obtient jamais tout l'ester attendu ;
\item \textbf{athermique} : elle ne dégage ni n'absorbe pratiquement de chaleur ; la température modifie la \emph{vitesse} d'atteinte de l'équilibre mais pas (ou très peu) sa position ;
\item \textbf{catalysée par les ions \ce{H+}} : quelques gouttes d'acide sulfurique concentré accélèrent la réaction sans modifier l'état d'équilibre.
\end{itemize}
\end{propriete}

\begin{aretenir}[Rendement à l'équilibre]
Pour un mélange \textbf{équimolaire} (1 mol d'acide + 1 mol d'alcool), le rendement à l'équilibre dépend de la classe de l'alcool :
\begin{center}
\begin{tabularx}{\linewidth}{|l|X|X|X|}
\hline
\rowcolor{popblueL} \textbf{Alcool} & \textbf{Primaire} & \textbf{Secondaire} & \textbf{Tertiaire} \\
\hline
Rendement en ester & environ 67 \% & environ 60 \% & environ 5 \% \\
\hline
\end{tabularx}
\end{center}
L'estérification directe est donc peu intéressante industriellement, surtout avec les alcools tertiaires.
\end{aretenir}

\begin{exoresolu}[Rendement de l'estérification directe]
\textbf{Énoncé.} On chauffe à reflux, en présence d'acide sulfurique, un mélange de $n_0 = \SI{1,0}{mol}$ d'acide éthanoïque et $n_0 = \SI{1,0}{mol}$ d'éthanol. À l'équilibre, le dosage montre qu'il reste \SI{0,33}{mol} d'acide éthanoïque.
\begin{enumerate}
\item Écrire l'équation-bilan de la réaction.
\item Dresser le tableau d'avancement et déterminer la quantité d'ester formée.
\item Calculer le rendement de la réaction.
\end{enumerate}
\tcblower
\textbf{Solution détaillée.}
\begin{enumerate}
\item Équation-bilan :
\[ \ce{CH3COOH + C2H5OH <=> CH3COOC2H5 + H2O} \]
\item Tableau d'avancement (quantités en mol) :
\begin{center}
\begin{tabular}{lcccc}
\hline
 & \ce{CH3COOH} & \ce{C2H5OH} & \ce{CH3COOC2H5} & \ce{H2O} \\
\hline
État initial & 1,0 & 1,0 & 0 & 0 \\
État d'équilibre & $1{,}0 - x_{éq}$ & $1{,}0 - x_{éq}$ & $x_{éq}$ & $x_{éq}$ \\
\hline
\end{tabular}
\end{center}
Il reste \SI{0,33}{mol} d'acide, donc $1{,}0 - x_{éq} = 0{,}33$, d'où $x_{éq} = \SI{0,67}{mol}$.
La quantité d'ester formée est donc $n_{ester} = \SI{0,67}{mol}$.
\item Si la réaction était totale, on obtiendrait $n_{max} = \SI{1,0}{mol}$ d'ester. Le rendement est :
\[ \eta = \frac{n_{ester}}{n_{max}} = \frac{0{,}67}{1{,}0} = 0{,}67 \quad \text{soit} \quad \eta \approx 67\,\%. \]
On retrouve la valeur caractéristique d'un mélange équimolaire acide + alcool primaire.
\end{enumerate}
\end{exoresolu}

\begin{attention}[Ne pas confondre vitesse et équilibre]
Le catalyseur \ce{H+} et la température permettent d'atteindre l'équilibre \emph{plus vite}, mais ne changent \textbf{pas} le rendement final. Pour améliorer le rendement, on peut utiliser un réactif en excès (alcool bon marché) ou éliminer l'eau au fur et à mesure (déplacement d'équilibre, loi de Le Chatelier).
\end{attention}

\courssub{Synthèses rapides et totales à partir des dérivés d'acide}

Pour contourner la lenteur et la limitation de l'estérification directe, les chimistes remplacent l'acide carboxylique par l'un de ses \cle{dérivés}, beaucoup plus réactifs : le chlorure d'acyle ou l'anhydride d'acide. L'idée est simple : dans \ce{RCOOH}, le groupe \ce{-OH} part difficilement ; dans \ce{RCOCl} et \ce{(RCO)2O}, le groupe partant (\ce{-Cl}, \ce{-OCOR}) part facilement, ce qui rend la réaction rapide et irréversible.

\textbf{1. À partir d'un chlorure d'acyle.} Le chlorure d'acyle réagit avec l'alcool selon :

\[ \ce{RCOCl + R$'$OH -> RCOO-R$'$ + HCl} \]

Cette réaction est \textbf{rapide} (quelques minutes), \textbf{totale} (rendement proche de 100 \%) et fortement exothermique. Elle dégage du chlorure d'hydrogène \ce{HCl}, un gaz acide et corrosif : on travaille donc sous hotte, souvent en présence d'une base (pyridine) qui piège le \ce{HCl} formé.

\begin{exemplebox}[Synthèse de l'éthanoate d'éthyle par le chlorure d'éthanoyle]
\[ \ce{CH3COCl + C2H5OH -> CH3COOC2H5 + HCl} \]
À partir de \SI{1}{mol} de chlorure d'éthanoyle et \SI{1}{mol} d'éthanol, on obtient pratiquement \SI{1}{mol} d'ester en quelques minutes, contre \SI{0,67}{mol} en plusieurs heures par voie directe.
\end{exemplebox}

\textbf{2. À partir d'un anhydride d'acide.} L'anhydride réagit avec l'alcool selon :

\[ \ce{(RCO)2O + R$'$OH -> RCOO-R$'$ + RCOOH} \]

Cette réaction est également \textbf{rapide} et \textbf{totale}, un peu moins violente que celle du chlorure d'acyle, et elle ne dégage pas de gaz corrosif : le sous-produit est simplement un acide carboxylique. C'est pourquoi les anhydrides sont très utilisés en synthèse, notamment pharmaceutique.

\begin{exemplebox}[Synthèse de l'éthanoate d'éthyle par l'anhydride éthanoïque]
\[ \ce{(CH3CO)2O + C2H5OH -> CH3COOC2H5 + CH3COOH} \]
\end{exemplebox}

\begin{popfigure}
\begin{center}
\begin{tikzpicture}
\begin{axis}[
width=0.95\linewidth, height=6.2cm,
xlabel={Temps $t$ (min)}, ylabel={Avancement $x$ (mol)},
xmin=0, xmax=60, ymin=0, ymax=1.15,
xtick={0,10,20,30,40,50,60},
ytick={0,0.33,0.67,1.0},
yticklabels={0, 0,33, 0,67, 1,0},
grid=major, grid style={popline!40},
legend style={at={(0.97,0.35)}, anchor=east, draw=popline, fill=white, font=\small},
axis lines=left, thick
]
\addplot[popcurve, popblue, domain=0:60, samples=200] {0.67*(1-exp(-x/12))};
\addlegendentry{Estérification directe (lente, limitée à 67 \%)}
\addplot[popcurve, poporange, domain=0:60, samples=200] {1-exp(-x/1.5)};
\addlegendentry{Chlorure d'acyle ou anhydride (rapide, total)}
\addplot[popline, dashed, domain=0:60] {0.67};
\addplot[popline, dashed, domain=0:60] {1.0};
\node[font=\footnotesize, popmuted, anchor=west] at (axis cs:38,0.60) {asymptote $x_{éq} = 0{,}67$ mol};
\node[font=\footnotesize, popmuted, anchor=west] at (axis cs:38,1.03) {$x_{max} = 1$ mol};
\end{axis}
\end{tikzpicture}
\end{center}
\legende{Évolution de l'avancement $x$ en fonction du temps pour un mélange équimolaire (1 mol + 1 mol). En bleu : l'estérification directe, lente, plafonne à environ \SI{0,67}{mol} (équilibre). En orange : la réaction avec un chlorure d'acyle ou un anhydride, quasi instantanée et totale ($x \to \SI{1}{mol}$).}
\end{popfigure}

\begin{aretenir}[Comparaison des trois méthodes de synthèse d'un ester]
\begin{center}
\begin{tabularx}{\linewidth}{|X|X|X|X|}
\hline
\rowcolor{popblueL} \textbf{Réactif} & \textbf{Vitesse} & \textbf{Caractère} & \textbf{Sous-produit} \\
\hline
Acide carboxylique \ce{RCOOH} & lente & limitée (équilibre, $\approx$ 67 \%) & eau \ce{H2O} \\
\hline
Chlorure d'acyle \ce{RCOCl} & rapide & totale & \ce{HCl} (corrosif) \\
\hline
Anhydride \ce{(RCO)2O} & rapide & totale & acide \ce{RCOOH} \\
\hline
\end{tabularx}
\end{center}
Industriellement, on préfère les dérivés d'acide : le gain de temps et le rendement de 100 \% compensent largement le coût de la préparation préalable du dérivé.
\end{aretenir}

\courssub{Complément Bac : la synthèse de l'aspirine}

L'\cle{aspirine}, ou \cle{acide acétylsalicylique}, est l'un des médicaments les plus consommés au monde, y compris au Cameroun où elle figure dans la trousse à pharmacie de nombreuses familles. Sa synthèse illustre parfaitement l'intérêt des anhydrides.

La molécule de départ est l'\cle{acide salicylique}, qui possède sur le même cycle benzénique un groupe carboxyle \ce{-COOH} et un groupe hydroxyle phénolique \ce{-OH}. Pour fabriquer l'aspirine, il faut transformer ce \ce{-OH} en groupe ester \ce{-O-CO-CH3}. On pourrait utiliser l'acide éthanoïque, mais la réaction serait lente et limitée : on utilise donc l'\textbf{anhydride éthanoïque}, en chauffant légèrement en présence de quelques gouttes d'acide phosphorique (catalyseur) :

\begin{popfigure}
\begin{center}
\chemfig{[:30]*6((-HO)-=-(-[:60]C(=[2]O)-[:-60]OH)=-=)}
\hspace{0.8em}{\Large +}\hspace{0.8em}
\chemfig{CH_3-C(=[1]O)-[7]O-C(=[5]O)-[1]CH_3}
\hspace{0.8em}{\Large $\longrightarrow$}\hspace{0.8em}
\chemfig{[:30]*6((-O-C(=[2]O)-[:-60]CH_3)-=-(-[:60]C(=[2]O)-[:-60]OH)=-=)}
\hspace{0.8em}{\Large +}\hspace{0.8em}
\chemfig{CH_3-C(=[1]O)-[7]OH}
\end{center}
\legende{Synthèse de l'aspirine : l'acide salicylique (à gauche) réagit avec l'anhydride éthanoïque pour donner l'acide acétylsalicylique (aspirine) et l'acide éthanoïque. Réaction rapide et totale, vers \SI{60}{\celsius}, catalysée par \ce{H3PO4}.}
\end{popfigure}

\[ \ce{C7H6O3 + (CH3CO)2O -> C9H8O4 + CH3COOH} \]

Cette réaction est \textbf{rapide et totale} : le rendement est excellent, ce qui explique le faible coût de l'aspirine. Remarque au passage que le groupe \ce{-COOH} de l'acide salicylique n'est \emph{pas} touché : seul le groupe hydroxyle est estérifié.

\begin{savaistu}[Les esters, des odeurs fruitées aux arômes du quotidien]
De nombreux esters possèdent des odeurs fruitées agréables : l'\textbf{éthanoate d'isoamyle} (éthanoate de 3-méthylbutyle) sent la \textbf{banane}, l'éthanoate d'éthyle rappelle l'ananas, le butanoate d'éthyle évoque la pomme. Industriellement, ces esters servent d'\textbf{arômes alimentaires} (bonbons, sirops, yaourts) et entrent dans la composition des \textbf{parfums}. Les fruits eux-mêmes doivent une partie de leur parfum naturel à des esters ! Quant à l'aspirine, son principe actif a été découvert à partir de l'écorce de saule, utilisée depuis l'Antiquité contre la fièvre — un beau pont entre pharmacopée traditionnelle et chimie moderne.
\end{savaistu}

\begin{exercice}[Pour t'entraîner]
On veut préparer de l'éthanoate de propyle \ce{CH3COO-CH2CH2CH3}.
\begin{enumerate}
\item Nommer et écrire les formules semi-développées des trois couples de réactifs possibles : (acide + alcool), (chlorure d'acyle + alcool), (anhydride + alcool).
\item Écrire les trois équations-bilans correspondantes.
\item On part de \SI{0,50}{mol} de chaque réactif. Quelle quantité maximale d'ester obtient-on dans chaque cas ? Conclure.
\end{enumerate}
\end{exercice}

\begin{corrige}[Exercice : synthèse de l'éthanoate de propyle]
\begin{enumerate}
\item Acide éthanoïque \ce{CH3COOH} + propan-1-ol \ce{CH3CH2CH2OH} ; chlorure d'éthanoyle \ce{CH3COCl} + propan-1-ol ; anhydride éthanoïque \ce{(CH3CO)2O} + propan-1-ol.
\item $\ce{CH3COOH + C3H7OH <=> CH3COOC3H7 + H2O}$ ; $\ce{CH3COCl + C3H7OH -> CH3COOC3H7 + HCl}$ ; $\ce{(CH3CO)2O + C3H7OH -> CH3COOC3H7 + CH3COOH}$.
\item Voie directe (alcool primaire) : $n \approx 0{,}67 \times 0{,}50 = \SI{0,335}{mol}$. Avec le chlorure ou l'anhydride : réaction totale, $n = \SI{0,50}{mol}$. Les dérivés d'acide donnent donc environ 50 \% d'ester en plus, et beaucoup plus vite.
\end{enumerate}
\end{corrige}

\begin{aretenir}[L'essentiel de la section]
\begin{itemize}
\item Un ester a pour formule \ce{R-COO-R$'$} et se nomme « \textbf{alkanoate d'alkyle} » : la partie \ce{R-CO-} vient de l'acide, le groupe \ce{R$'$} vient de l'alcool.
\item L'estérification directe \ce{RCOOH + R$'$OH <=> RCOOR$'$ + H2O} est \textbf{lente}, \textbf{limitée} ($\approx$ 67 \% avec un alcool primaire), \textbf{athermique} et \textbf{catalysée par \ce{H+}}.
\item Avec un chlorure d'acyle ($\to$ \ce{HCl}) ou un anhydride ($\to$ \ce{RCOOH}), la réaction est \textbf{rapide et totale} : c'est la voie industrielle, illustrée par la synthèse de l'\textbf{aspirine}.
\end{itemize}
\end{aretenir}

\coursec{Hydrolyse, saponification et fabrication du savon}

Dans la section précédente, tu as découvert comment un ester naît de la rencontre entre un acide carboxylique et un alcool : l'estérification, réaction lente et limitée. Une question se pose alors naturellement : peut-on faire le chemin inverse, c'est-à-dire \emph{casser} un ester pour retrouver l'acide et l'alcool ? Et surtout, peut-on rendre cette destruction \emph{totale} ? C'est précisément l'enjeu de cette section, et la réponse a une application que tu connais bien : chaque fois que tu te laves les mains avec un savon fabriqué à partir d'huile de palme, tu utilises le produit d'une réaction de saponification. Suivons le raisonnement pas à pas.

\courssub{Hydrolyse des esters : la réaction inverse de l'estérification}

\textbf{Le principe}

L'estérification consomme un acide et un alcool pour produire un ester et de l'eau. Si l'on met un ester en présence d'un \emph{excès d'eau}, la réaction peut s'inverser : c'est l'\cle{hydrolyse} de l'ester (du grec \emph{hydro}, eau, et \emph{lusis}, rupture : « rupture par l'eau »).

\begin{definition}[Hydrolyse d'un ester]
L'\cle{hydrolyse} d'un ester est la réaction de l'eau sur un ester, qui régénère l'acide carboxylique et l'alcool :
\[ \ce{RCOO-R$'$ + H2O <=> RCOOH + R$'$-OH} \]
C'est la réaction \emph{inverse} de l'estérification directe. Comme elle, elle est \textbf{lente}, \textbf{limitée} (elle conduit à un équilibre) et \textbf{faiblement endothermique}. Elle est catalysée par les ions \ce{H3O+} (acide sulfurique dilué) et favorisée par la chaleur.
\end{definition}

\begin{exemplebox}[Hydrolyse de l'éthanoate d'éthyle]
\[ \ce{CH3COO-CH2-CH3 + H2O <=> CH3COOH + CH3-CH2-OH} \]
Si l'on part de \SI{1,0}{mol} d'éthanoate d'éthyle et \SI{1,0}{mol} d'eau, on obtient à l'équilibre environ \SI{0,33}{mol} d'acide éthanoïque et \SI{0,33}{mol} d'éthanol : les deux tiers de l'ester n'ont pas réagi. La limite de l'estérification est aussi la limite de l'hydrolyse.
\end{exemplebox}

\begin{attention}[Hydrolyse et estérification : même équilibre]
Ne crois pas que l'hydrolyse « marche mieux » que l'estérification. Les deux réactions conduisent au \emph{même état d'équilibre} : seul le point de départ change. Pour déplacer l'équilibre vers l'hydrolyse, il faut utiliser un \emph{large excès d'eau} (loi de Le Chatelier). Mais même ainsi, la réaction reste lente. Pour casser un ester de façon \emph{totale}, il faut changer de stratégie : travailler en milieu basique. C'est la saponification.
\end{attention}

\courssub{La saponification : une hydrolyse basique totale}

\textbf{De l'intuition à la définition}

Reprenons le problème : l'hydrolyse acide est limitée parce que l'acide carboxylique formé peut réagir à nouveau avec l'alcool. L'idée géniale consiste à remplacer l'eau par une \emph{base forte} (soude \ce{NaOH} ou potasse \ce{KOH}). L'ion hydroxyde \ce{HO-} attaque l'ester et libère l'ion \cle{carboxylate} \ce{RCOO-}. Or cet ion, base conjuguée d'un acide faible, est \emph{incapable} de réagir avec l'alcool pour reformer l'ester : la réaction inverse devient impossible. La réaction est donc \textbf{totale}.

\begin{definition}[Saponification]
La \cle{saponification} est l'hydrolyse basique d'un ester par une solution concentrée de soude (\ce{NaOH}) ou de potasse (\ce{KOH}), à chaud :
\[ \ce{RCOO-R$'$ + NaOH -> RCOO-Na+ + R$'$-OH} \]
Elle est \textbf{lente} (d'où le chauffage) mais \textbf{totale}. Elle produit un \emph{carboxylate de sodium} (ou de potassium) et un alcool. Lorsque l'ester est un corps gras, le carboxylate obtenu est un \cle{savon}.
\end{definition}

\begin{exemplebox}[Saponification de l'éthanoate d'éthyle]
\[ \ce{CH3COO-CH2-CH3 + NaOH -> CH3COO-Na+ + CH3-CH2-OH} \]
On obtient l'éthanoate de sodium et l'éthanol. Remarque la \emph{flèche simple} : la réaction est totale. À chaud, quelques dizaines de minutes suffisent, alors que l'hydrolyse acide demanderait des jours sans jamais être complète.
\end{exemplebox}

\begin{attention}[L'erreur classique du baccalauréat]
L'erreur la plus fréquente est d'écrire la saponification avec une \emph{double flèche} \ce{<=>}, par analogie avec l'estérification. C'est faux : la saponification est \textbf{totale}, elle s'écrit avec une \textbf{flèche simple} \ce{->}. Autre piège : ne jamais écrire \ce{RCOOH} comme produit de la saponification — en milieu basique, l'acide est immédiatement neutralisé en ion carboxylate \ce{RCOO-}. Enfin, la soude est un \emph{réactif} (elle est consommée), pas un catalyseur : il faut l'écrire dans l'équation.
\end{attention}

\courssub{Les corps gras naturels : les triglycérides}

Pour fabriquer un savon, il faut un ester dont la chaîne carbonée est \emph{longue} : c'est elle qui donnera au savon ses propriétés lavantes. La nature nous fournit de tels esters en abondance : les \cle{corps gras} (huiles et graisses), qui sont des \cle{triglycérides}.

\begin{definition}[Triglycéride]
Un \cle{triglycéride} est un \emph{triester} formé à partir d'une molécule de \cle{glycérol} (propane-1,2,3-triol, \ce{HO-CH2-CH(OH)-CH2-OH}) et de \emph{trois} molécules d'\cle{acides gras}, c'est-à-dire des acides carboxyliques à longue chaîne carbonée non ramifiée (de 12 à 18 atomes de carbone environ).
\end{definition}

Les acides gras les plus courants dans les huiles camerounaises sont :

\begin{center}
\begin{tabularx}{\linewidth}{|l|X|l|}
\hline
\rowcolor{popblueL} \textbf{Acide gras} & \textbf{Formule semi-développée} & \textbf{Présence principale} \\
\hline
Acide palmitique & \ce{CH3-(CH2)14-COOH} (saturé, C16) & Huile de palme \\
\hline
Acide stéarique & \ce{CH3-(CH2)16-COOH} (saturé, C18) & Graisses animales, beurre de karité \\
\hline
Acide oléique & \ce{CH3-(CH2)7-CH=CH-(CH2)7-COOH} (insaturé, C18:1) & Huile d'olive, huile de palme \\
\hline
\end{tabularx}
\end{center}

Le glycérol possède trois groupes hydroxyle \ce{-OH} ; chacun peut estérifier un acide gras. Un triglycéride a donc la structure générale suivante, où \ce{R}, \ce{R$'$} et \ce{R$''$} sont de longues chaînes carbonées (souvent identiques ou très voisines) :

\[ \chemfig{CH_2(-[2]O-C(=[2]O)-R)-CH(-[2]O-C(=[2]O)-R')-CH_2(-[2]O-C(=[2]O)-R'')} \]

\begin{savaistu}[Huile de palme, or rouge du Cameroun]
Au Cameroun, l'\emph{huile de palme} (extraite de la pulpe du fruit du palmier) et l'\emph{huile de palmiste} (extraite de la noix) sont les matières premières traditionnelles de la fabrication artisanale du savon, notamment dans les savonneries villageoises du Littoral et du Sud-Ouest. Industriellement, la SOCOPALM et la SAFACAM fournissent l'huile aux savonneries modernes. Le fameux « savon noir » artisanal doit sa couleur à la cuisson prolongée de l'huile de palme rouge avec la potasse issue de cendres de bois.
\end{savaistu}

\courssub{Fabrication du savon : la saponification d'un triglycéride}

\textbf{L'équation-bilan}

Chauffons un triglycéride avec une solution concentrée de soude. Les trois fonctions ester sont saponifiées : on obtient \emph{trois} molécules de carboxylate de sodium (le savon) et une molécule de glycérol.

\begin{methode}[Écrire l'équation d'une saponification de triglycéride]
\begin{enumerate}
\item Écrire la formule du triglycéride en identifiant les trois chaînes \ce{R-COO} issues des acides gras.
\item Ajouter \textbf{3} moles de soude : \ce{+ 3 NaOH} (une par fonction ester).
\item Écrire les produits : \textbf{3} moles de carboxylate de sodium \ce{3 RCOO-Na+} (le savon) \textbf{+ 1} mole de glycérol \ce{HO-CH2-CH(OH)-CH2-OH}.
\item Vérifier l'équilibre des atomes : même nombre de C, H, O et Na de chaque côté de la flèche.
\end{enumerate}
\end{methode}

\begin{exemplebox}[Saponification de la tripalmitine]
La \emph{tripalmitine} est le triglycéride de l'acide palmitique \ce{C15H31-COOH}, abondant dans l'huile de palme. Sa saponification s'écrit :
\[ \ce{C51H98O6 + 3 NaOH -> 3 C15H31COO-Na+ + C3H8O3} \]
soit, en nommant les espèces : tripalmitine + soude $\longrightarrow$ 3 palmitate de sodium (savon) + glycérol. Vérifions le sodium : 3 à gauche, $3 \times 1 = 3$ à droite. L'équation est équilibrée.
\end{exemplebox}

\begin{popfigure}
\centering
\begin{tikzpicture}[scale=0.9, transform shape]
% Triglycéride unique (représentation schématique)
\node[font=\small] (tg) at (-4.5,1.2) {\chemfig{CH_2(-[2]O-C(=[2]O)-C_{15}H_{31})-CH(-[2]O-C(=[2]O)-C_{15}H_{31})-CH_2(-[2]O-C(=[2]O)-C_{15}H_{31})}};
\node[font=\footnotesize\itshape, popmuted] at (-4.5,-0.5) {tripalmitine (corps gras)};
% + soude
\node[font=\small\bfseries] at (-0.5,1.2) {$+$ \quad \ce{3 NaOH}};
\node[font=\footnotesize\itshape, popmuted] at (-0.5,0.5) {soude concentrée, à chaud};
% flèche
\draw[-{Stealth}, very thick, poporange] (1.0,1.2) -- (2.6,1.2);
% produits
\node[font=\small] at (5.0,1.9) {\ce{3 C15H31-COO-Na+}};
\node[font=\footnotesize\itshape, popmuted] at (5.0,1.2) {palmitate de sodium (savon)};
\node[font=\small\bfseries] at (3.8,0.1) {$+$};
\node[font=\small] at (5.6,0.1) {\chemfig{HO-CH_2-CH(-[2]OH)-CH_2-OH}};
\node[font=\footnotesize\itshape, popmuted] at (5.6,-0.7) {glycérol};
\end{tikzpicture}
\legende{Saponification de la tripalmitine par la soude à chaud : 1 mole de triglycéride donne 3 moles de savon et 1 mole de glycérol. Réaction totale (flèche simple).}
\end{popfigure}

\textbf{Les étapes pratiques de la fabrication artisanale}

\begin{experience}[Fabrication d'un savon à l'huile de palme]
\textbf{Protocole :}
\begin{enumerate}
\item \textbf{Chauffage (saponification)} : dans une marmite, chauffer l'huile de palme avec une solution concentrée de soude en agitant régulièrement. Le mélange s'épaissit progressivement : c'est la « trace ».
\item \textbf{Relargage} : ajouter une solution concentrée de chlorure de sodium \ce{NaCl} (sel de cuisine). Le savon, \emph{très peu soluble dans l'eau salée}, précipite et surnage. Le glycérol, lui, reste dissous dans la phase aqueuse.
\item \textbf{Filtration} : recueillir le savon solide par filtration ou écumage.
\item \textbf{Lavage} : laver le savon à l'eau froide pour éliminer l'excès de soude (corrosive !).
\item \textbf{Séchage et moulage} : presser le savon dans des moules et le laisser sécher à l'air libre plusieurs jours.
\end{enumerate}
\textbf{Observation :} le produit obtenu mousse au contact de l'eau et dissout les graisses : c'est bien un savon.
\end{experience}

\begin{popfigure}
\centering
\begin{tikzpicture}[
  etape/.style={rectangle, rounded corners=4pt, draw=popblue, fill=popblueL, thick, minimum width=2.6cm, minimum height=1.1cm, align=center, font=\footnotesize\bfseries},
  detail/.style={font=\tiny\itshape, text=popmuted, align=center},
  fleche/.style={-{Stealth}, very thick, poporange}]
\node[etape] (e1) at (0,0) {Chauffage\\huile + soude};
\node[etape] (e2) at (3.4,0) {Relargage\\ajout de \ce{NaCl}};
\node[etape] (e3) at (6.8,0) {Filtration\\et lavage};
\node[etape] (e4) at (10.2,0) {Séchage\\et moulage};
\draw[fleche] (e1) -- (e2);
\draw[fleche] (e2) -- (e3);
\draw[fleche] (e3) -- (e4);
\node[detail] at (0,-1.0) {saponification\\(réaction totale)};
\node[detail] at (3.4,-1.0) {le savon précipite\\le glycérol reste en solution};
\node[detail] at (6.8,-1.0) {élimination de\\l'excès de soude};
\node[detail] at (10.2,-1.0) {savon dur\\prêt à l'emploi};
\end{tikzpicture}
\legende{Organigramme des étapes de la fabrication artisanale du savon.}
\end{popfigure}

\begin{attention}[Sécurité et choix de la base]
La soude et la potasse concentrées sont \emph{corrosives} : port de gants et de lunettes obligatoire, et jamais d'eau versée sur la soude solide (toujours l'inverse, lentement). Par ailleurs, la nature de la base conditionne la texture du savon :
\begin{itemize}
\item avec la \textbf{soude} \ce{NaOH} : savon \emph{dur} (savon de Marseille, savon en pain) ;
\item avec la \textbf{potasse} \ce{KOH} : savon \emph{mou} ou \emph{liquide} (savon noir, savon liquide), car les carboxylates de potassium sont plus solubles.
\end{itemize}
\end{attention}

\courssub{Structure et action détergente du savon}

Pourquoi le savon lave-t-il ? Tout tient dans la structure de l'ion carboxylate \ce{R-COO-}, qui est \cle{amphiphile} : il possède deux parties de comportement opposé vis-à-vis de l'eau.

\begin{definition}[Ion carboxylate amphiphile]
L'ion carboxylate d'un savon comporte :
\begin{itemize}
\item une \cle{tête hydrophile} : le groupe ionique \ce{-COO-}, qui aime l'eau (polaire) ;
\item une longue \cle{queue hydrophobe} (et \emph{lipophile}) : la chaîne carbonée \ce{R}, qui fuit l'eau mais se lie aux graisses.
\end{itemize}
\end{definition}

Au contact d'une salissure grasse, les queues hydrophobes se plantent dans la goutte de graisse tandis que les têtes hydrophiles restent dans l'eau : la graisse se fragmente en gouttelettes entourées d'ions carboxylate, appelées \cle{micelles}. Les têtes négatives se repoussant entre elles, les micelles ne se réunissent pas et sont emportées par l'eau de rinçage : la graisse est éliminée.

\begin{popfigure}
\centering
\begin{tikzpicture}[scale=1.0]
% Ion carboxylate seul
\draw[very thick, popdark] (-6.5,2.2) -- (-4.4,2.2);
\draw[popdark, thick, decorate, decoration={coil, segment length=4pt, amplitude=1.5pt}] (-6.5,2.2) -- (-4.4,2.2);
\fill[popblue] (-4.2,2.2) circle (0.28);
\node[white, font=\scriptsize\bfseries] at (-4.2,2.2) {\ce{COO-}};
\node[font=\scriptsize\itshape, popmuted, align=center] at (-5.6,1.6) {queue hydrophobe\\(chaîne \ce{R})};
\node[font=\scriptsize\itshape, popmuted, align=center] at (-3.4,1.6) {tête\\hydrophile};
\node[font=\small\bfseries, popdark] at (-5.2,3.0) {Ion carboxylate};
% Micelle
\node[font=\small\bfseries, popdark] at (2.5,3.0) {Micelle autour d'une salissure grasse};
\fill[popgoldL] (2.5,0.6) circle (0.55);
\node[font=\scriptsize\itshape, popdark] at (2.5,0.6) {graisse};
\foreach \a in {0,30,60,90,120,150,180,210,240,270,300,330}{
  \draw[popdark, thick] ({2.5+0.62*cos(\a)},{0.6+0.62*sin(\a)}) -- ({2.5+1.15*cos(\a)},{0.6+1.15*sin(\a)});
  \fill[popblue] ({2.5+1.32*cos(\a)},{0.6+1.32*sin(\a)}) circle (0.13);
}
\node[font=\scriptsize\itshape, popmuted, align=center] at (5.3,0.6) {têtes hydrophiles\\vers l'eau};
\node[font=\scriptsize\itshape, popmuted, align=center] at (-0.4,0.6) {queues hydrophobes\\dans la graisse};
\end{tikzpicture}
\legende{Structure amphiphile de l'ion carboxylate et formation d'une micelle : les queues hydrophobes emprisonnent la graisse, les têtes hydrophiles assurent la dispersion dans l'eau.}
\end{popfigure}

\begin{exoresolu}[Masse de savon obtenue à partir de la trioléine]
\textbf{Énoncé.} La trioléine est le triglycéride de l'acide oléique \ce{C17H33-COOH}, de masse molaire $M(\text{trioléine}) \approx \SI{885}{g.mol^{-1}}$. On saponifie \SI{885}{g} de trioléine par un excès de soude. On donne $M(\text{oléate de sodium } \ce{C17H33COONa}) \approx \SI{304}{g.mol^{-1}}$.
\begin{enumerate}
\item Écrire l'équation-bilan de la réaction.
\item Calculer la masse de savon obtenue.
\end{enumerate}
\tcblower
\textbf{Solution détaillée.}
\begin{enumerate}
\item La saponification est totale (flèche simple) :
\[ \ce{C57H104O6 + 3 NaOH -> 3 C17H33COO-Na+ + C3H8O3} \]
(trioléine + soude $\rightarrow$ 3 oléate de sodium + glycérol).
\item Quantité de trioléine :
\[ n(\text{trioléine}) = \frac{m}{M} = \frac{885}{885} = \SI{1,0}{mol} \]
D'après les coefficients stœchiométriques, 1 mole de triglycéride donne 3 moles de savon :
\[ n(\text{savon}) = 3 \times 1,0 = \SI{3,0}{mol} \]
D'où la masse de savon :
\[ m(\text{savon}) = n \times M = 3,0 \times 304 = \SI{912}{g} \]
\end{enumerate}
\textbf{Conclusion :} à partir de \SI{885}{g} de trioléine, on obtient environ \SI{912}{g} d'oléate de sodium. La masse de savon est \emph{supérieure} à la masse de corps gras de départ, car le sodium s'est incorporé au produit. C'est pourquoi la savonnerie est économiquement intéressante !
\end{exoresolu}

\begin{aretenir}[L'essentiel de la section]
\begin{itemize}
\item L'\cle{hydrolyse} d'un ester (\ce{RCOO-R$'$ + H2O <=> RCOOH + R$'$-OH}) est lente et limitée : réaction inverse de l'estérification.
\item La \cle{saponification} est l'hydrolyse \emph{basique} d'un ester par \ce{NaOH} ou \ce{KOH} concentré, à chaud : lente mais \textbf{totale} (flèche simple), elle donne un carboxylate (savon) et un alcool.
\item Les corps gras sont des \cle{triglycérides} : triesters du glycérol et d'acides gras. 1 triglycéride + 3 \ce{NaOH} $\rightarrow$ 3 savons + 1 glycérol.
\item Étapes de fabrication : chauffage, \cle{relargage} au sel, filtration, lavage, séchage. Soude $\rightarrow$ savon dur ; potasse $\rightarrow$ savon mou.
\item Le savon est \cle{amphiphile} : tête hydrophile \ce{-COO-}, queue hydrophobe \ce{R} ; il forme des \cle{micelles} qui emprisonnent les graisses.
\end{itemize}
\end{aretenir}

\begin{exercice}[Savon de karité]
Le beurre de karité contient de la tristéarine, triglycéride de l'acide stéarique \ce{C17H35-COOH}, de masse molaire $M = \SI{891}{g.mol^{-1}}$. On traite \SI{1,782}{kg} de tristéarine par la potasse \ce{KOH} en excès. On donne $M(\ce{C17H35COOK}) = \SI{322}{g.mol^{-1}}$.
\begin{enumerate}
\item Écrire l'équation-bilan de la saponification.
\item Calculer la masse de savon obtenu. S'agit-il d'un savon dur ou mou ? Justifier.
\item Calculer la masse de glycérol formé ($M(\ce{C3H8O3}) = \SI{92}{g.mol^{-1}}$).
\end{enumerate}
\end{exercice}

\begin{corrige}[Exercice : Savon de karité]
\begin{enumerate}
\item $\displaystyle \ce{C57H110O6 + 3 KOH -> 3 C17H35COO-K+ + C3H8O3}$ (réaction totale, flèche simple).
\item $n(\text{tristéarine}) = \dfrac{1782}{891} = \SI{2,0}{mol}$, donc $n(\text{savon}) = 3 \times 2,0 = \SI{6,0}{mol}$ et $m = 6,0 \times 322 = \SI{1932}{g} \approx \SI{1,93}{kg}$. La base utilisée est la \emph{potasse} \ce{KOH} : c'est donc un savon \emph{mou}, car les carboxylates de potassium sont plus solubles dans l'eau.
\item $n(\text{glycérol}) = n(\text{tristéarine}) = \SI{2,0}{mol}$, donc $m = 2,0 \times 92 = \SI{184}{g}$.
\end{enumerate}
\end{corrige}

Tu sais désormais transformer un corps gras en savon et expliquer son pouvoir lavant. Il reste un dernier dérivé des acides carboxyliques à explorer, aux applications spectaculaires : les amides, portes d'entrée vers les grandes polymères de notre quotidien — le nylon 6,6 et le tergal.

\coursec{Amides et polycondensation : nylon 6,6 et tergal}

Dans la section précédente, nous avons vu comment un ester peut être hydrolysé ou saponifié : la liaison \ce{-CO-O-} se casse sous l'action de l'eau ou de la soude. Remplace maintenant l'atome d'oxygène de cette liaison par un groupe \ce{-NH-} : tu obtiens une nouvelle famille, les \cle{amides}, dont la liaison \ce{-CO-NH-} est remarquablement solide. C'est cette solidité qui explique pourquoi la nature l'a choisie pour bâtir les protéines, et pourquoi l'industrie l'exploite pour fabriquer des fibres textiles exceptionnellement résistantes comme le \cle{nylon}. Dans cette dernière section, tu vas apprendre à former les amides, à les nommer, puis à assembler des milliers de petites molécules en chaînes géantes : c'est la \cle{polycondensation}, la réaction qui produit le nylon 6,6 et le tergal qui habillent chaque jour des millions de Camerounais.

\courssub{Les amides : formation et nomenclature}

\textbf{Intuition.} Un acide carboxylique \ce{R-COOH} est acide ; l'ammoniac \ce{NH3} est une base. Quand on les mélange, la première chose qui se produit est donc une simple réaction acide-base : on obtient un sel. Mais si l'on chauffe ce sel, il perd une molécule d'eau et se transforme en un composé stable, l'amide. Voilà l'idée ; passons à la rigueur.

\begin{definition}[Amide]
Un \cle{amide} est un dérivé d'acide carboxylique dans lequel le groupe \chemfig{-OH} du groupe carboxyle est remplacé par un groupe \chemfig{-NH_2}, \chemfig{-N(-[2]H)-R'} ou \chemfig{-N(-[2]R')(-[6]R'')}. Sa formule générale est \chemfig{R-C(=[1]O)-NH_2} (amide non substitué). Le groupe caractéristique \chemfig{-C(=[1]O)-NH-} est appelé \cle{groupe amide} ; lorsqu'il relie deux chaînes carbonées, on parle de \cle{liaison amide} (ou liaison peptique en biochimie).
\end{definition}

\textbf{Formation à partir de l'acide carboxylique (en deux étapes).}

\emph{Étape 1 : réaction acide-base, rapide et totale, à température ordinaire.} L'acide éthanoïque réagit avec l'ammoniac pour donner l'éthanoate d'ammonium :
\[ \ce{CH3-COOH + NH3 -> CH3-COO- + NH4+} \]
Le produit s'écrit \ce{CH3-COO^{-}NH4+} : c'est un sel d'ammonium, l'éthanoate d'ammonium.

\emph{Étape 2 : déshydratation thermique.} En chauffant fortement le sel (vers \SI{150}{\celsius} à \SI{200}{\celsius}), on élimine une molécule d'eau :
\[ \ce{CH3-COO^{-}NH4+ ->[\Delta] CH3-CO-NH2 + H2O} \]
Le produit obtenu, \ce{CH3-CO-NH2}, est l'\cle{éthanamide}.

\begin{attention}[Une réaction en deux temps]
Ne confonds pas les deux étapes ! À froid, \ce{RCOOH + NH3} donne un \textbf{sel} (\ce{RCOO^{-}NH4+}), pas un amide. L'amide n'apparaît qu'\textbf{à chaud}, par élimination d'eau. À l'examen, écrire directement \ce{RCOOH + NH3 -> RCONH2 + H2O} sans préciser le chauffage est une imprécision qui coûte des points.
\end{attention}

\textbf{Formation rapide et totale à partir des dérivés d'acide.} Comme pour les esters, on préfère en pratique utiliser les dérivés réactifs de l'acide : le \cle{chlorure d'acyle} ou l'\cle{anhydride d'acide}. Ces réactions sont \textbf{rapides et totales}, à température ordinaire :
\[ \ce{CH3-COCl + NH3 -> CH3-CO-NH2 + HCl} \]
\[ \ce{CH3-COCl + CH3-NH2 -> CH3-CO-NH-CH3 + HCl} \]
Avec une amine primaire ($R'-\ce{NH2}$) ou secondaire ($R'R''\ce{NH}$), on obtient un amide \cle{N-substitué}. L'équation générale avec un anhydride s'écrit :
\[ \ce{(R-CO)2O + 2 NH3 -> R-CO-NH2 + R-COO- + NH4+} \]

\textbf{Nomenclature des amides.} La règle est simple : on remplace la terminaison « oïque » du nom de l'acide par « \cle{amide} ».

\begin{center}
\begin{tabularx}{\linewidth}{|l|X|l|}
\hline
\rowcolor{popblueL} \textbf{Formule} & \textbf{Acide parent} & \textbf{Nom de l'amide} \\
\hline
\ce{H-CO-NH2} & acide méthanoïque & méthanamide \\
\hline
\ce{CH3-CO-NH2} & acide éthanoïque & éthanamide \\
\hline
\ce{C6H5-CO-NH2} & acide benzoïque & benzamide \\
\hline
\ce{CH3-CO-NH-C2H5} & acide éthanoïque & N-éthyléthanamide \\
\hline
\end{tabularx}
\end{center}

Pour les amides substitués sur l'azote, on fait précéder le nom du ou des groupes alkyles de la lettre \cle{N} (en italique), qui indique que la substitution porte sur l'atome d'azote : \ce{CH3-CO-NH-CH3} est le \emph{N}-méthyléthanamide.

\begin{savaistu}[La liaison amide, brique du vivant]
La liaison \ce{-CO-NH-} est appelée \cle{liaison peptique} lorsqu'elle relie des acides $\alpha$-aminés : c'est elle qui assemble les protéines de ton corps, de tes muscles à tes enzymes. L'urée \ce{H2N-CO-NH2}, amide de l'acide carbamique, est à la fois un déchet du métabolisme humain et l'un des engrais azotés les plus utilisés dans les champs de maïs du Cameroun. Les amides sont partout !
\end{savaistu}

\courssub{La polycondensation : définition et principe}

\textbf{Intuition.} Imagine des perles qui possèdent chacune \emph{deux} crochets, un à chaque extrémité. En accrochant les perles bout à bout, tu fabriques un collier aussi long que tu veux. En chimie, une molécule qui porte \textbf{deux} groupes fonctionnels réactifs (on dit qu'elle est \cle{bifonctionnelle}) peut jouer ce rôle de perle : chaque extrémité réagit avec une autre molécule, et la chaîne grandit. Si chaque maillon formé s'accompagne du départ d'une petite molécule (eau ou chlorure d'hydrogène), on parle de polycondensation.

\begin{definition}[Polycondensation]
Une \cle{polycondensation} est une polymérisation au cours de laquelle des molécules \textbf{bifonctionnelles} (ou polyfonctionnelles) réagissent entre elles par réactions successives, avec \textbf{élimination de petites molécules} (eau \ce{H2O}, chlorure d'hydrogène \ce{HCl}, alcool). Le macromolécule obtenue est un \cle{polymère de condensation} ; selon la nature de la liaison formée, on obtient un \cle{polyamide} (liaison \ce{-CO-NH-}) ou un \cle{polyester} (liaison \ce{-CO-O-}).
\end{definition}

\begin{methode}[Écrire l'équation d'une polycondensation]
\begin{enumerate}
\item \textbf{Repérer les deux monomères bifonctionnels} : vérifier que chacun porte bien \emph{deux} fonctions réactives (diacide + diamine, ou diacide + diol).
\item \textbf{Faire réagir les fonctions deux à deux} : un groupe \ce{-COOH} réagit avec un groupe \ce{-NH2} (formation d'une liaison amide) ou avec un groupe \ce{-OH} (formation d'une liaison ester), en éliminant \ce{H2O}.
\item \textbf{Écrire le motif} de la chaîne entre crochets, affecté de l'indice $n$, et faire apparaître les molécules d'eau éliminées ($n$ \ce{H2O} pour $n$ motifs formés à partir de $n$ molécules de chaque monomère).
\end{enumerate}
\end{methode}

\begin{attention}[Les trois pièges classiques]
\begin{itemize}
\item \textbf{Oublier l'eau éliminée} : une équation de polycondensation sans les $n$ \ce{H2O} (ou $n$ \ce{HCl}) est \emph{fausse} car non équilibrée.
\item \textbf{Utiliser un monomère monofonctionnel} : un acide à un seul \ce{-COOH} (comme \ce{CH3COOH}) ne peut former qu'une seule liaison, puis la croissance s'arrête : \emph{pas de chaîne possible}. La bifonctionnalité est indispensable.
\item \textbf{Mal écrire le motif} : le motif doit être l'unité qui se \emph{répète} ; vérifie qu'en le recopiant bout à bout tu reconstitues bien la chaîne.
\end{itemize}
\end{attention}

\courssub{Le nylon 6,6 : un polyamide}

Le nylon 6,6 est obtenu par polycondensation de deux monomères :
\begin{itemize}
\item l'\cle{acide hexanedioïque} (ou acide adipique) : \ce{HOOC-(CH2)4-COOH}, un diacide à 6 atomes de carbone ;
\item l'\cle{hexane-1,6-diamine} : \ce{H2N-(CH2)6-NH2}, une diamine à 6 atomes de carbone.
\end{itemize}

Chaque groupe \ce{-COOH} réagit avec un groupe \ce{-NH2} en éliminant une molécule d'eau ; la réaction est conduite à chaud (environ \SI{250}{\celsius} à \SI{280}{\celsius}), sous pression contrôlée. Pour retenir l'ordre dans le nom (nylon 6,6), on écrit d'abord la diamine puis le diacide. L'équation-bilan s'écrit :

\[ n\,\ce{H2N-(CH2)6-NH2} + n\,\ce{HOOC-(CH2)4-COOH} \ce{->[\Delta]} \ce{[-NH-(CH2)6-NH-CO-(CH2)4-CO-]}_n + n\,\ce{H2O} \]

\begin{popfigure}
\begin{center}
\chemfig{H_2N-[,,,]CH_2-CH_2-CH_2-CH_2-CH_2-CH_2-NH_2}
\hspace{1cm}
\chemfig{HO-C(=[2]O)-[,,,]CH_2-CH_2-CH_2-CH_2-C(=[2]O)-OH}

\vspace{0.4cm}
$\xrightarrow{\ \Delta,\ \text{polycondensation}\ }\ n\,\ce{H2O}\ \text{éliminées}$
\vspace{0.4cm}

\fbox{\chemfig{-[2]N(-[3]H)-CH_2-CH_2-CH_2-CH_2-CH_2-CH_2-N(-[1]H)-C(=[2]O)-CH_2-CH_2-CH_2-CH_2-C(=[6]O)-[2]}}\,${}_n$
\end{center}
\legende{Formation du nylon 6,6 : l'hexane-1,6-diamine et l'acide hexanedioïque s'assemblent avec élimination d'eau ; le motif encadré (diamine puis diacide) se répète $n$ fois dans la macromolécule.}
\end{popfigure}

\textbf{Pourquoi « 6,6 » ?} La nomenclature des nylons indique le nombre d'atomes de carbone de chaque monomère : le \textbf{premier chiffre} correspond à la \textbf{diamine}, le \textbf{second} au \textbf{diacide}. Ici : diamine à 6 C, diacide à 6 C, d'où « nylon 6,6 ».

\begin{exemplebox}[Le nylon 6,10]
Appliquons la règle : le \cle{nylon 6,10} est obtenu à partir d'une diamine à 6 atomes de carbone, l'hexane-1,6-diamine \ce{H2N-(CH2)6-NH2}, et d'un diacide à 10 atomes de carbone, l'acide décanedioïque \ce{HOOC-(CH2)8-COOH}. Le motif de la chaîne est \ce{[-NH-(CH2)6-NH-CO-(CH2)8-CO-]}$_n$. Retiens l'ordre : \textbf{diamine d'abord, diacide ensuite}.
\end{exemplebox}

\courssub{Le tergal (PET) : un polyester}

Le \cle{tergal}, nom commercial du \cle{PET} (polyéthylène téréphtalate), est un polyester obtenu par polycondensation de :
\begin{itemize}
\item l'\cle{acide benzène-1,4-dicarboxylique} (acide téréphtalique) : \ce{HOOC-C6H4-COOH}, un diacide aromatique ;
\item l'\cle{éthane-1,2-diol} (éthylène glycol) : \ce{HO-CH2-CH2-OH}, un diol.
\end{itemize}

Chaque groupe \ce{-COOH} réagit avec un groupe \ce{-OH} par estérification, avec élimination d'eau, à chaud (environ \SI{250}{\celsius}) en présence d'un catalyseur :
\[ n\,\ce{HOOC-C6H4-COOH} + n\,\ce{HO-CH2-CH2-OH} \ce{->[\Delta][\text{catalyseur}]} \ce{[-CO-C6H4-CO-O-CH2-CH2-O-]}_n + n\,\ce{H2O} \]

\begin{popfigure}
\begin{center}
\chemfig{HO-C(=[2]O)-[:-30]*6(---C(=[2]O)-OH---)}
\hspace{0.8cm}
\chemfig{HO-CH_2-CH_2-OH}

\vspace{0.4cm}
$\xrightarrow{\ \Delta,\ \text{catalyseur}\ }\ n\,\ce{H2O}\ \text{éliminées}$
\vspace{0.4cm}

\fbox{\chemfig{-C(=[2]O)-[:-30]*6(---C(=[6]O)-O-CH_2-CH_2-O-[2]---)}}\,${}_n$
\end{center}
\legende{Formation du tergal (PET) : l'acide téréphtalique et l'éthylène glycol forment des liaisons ester \ce{-CO-O-} répétées ; le motif encadré (diacide puis diol) constitue l'unité de la chaîne. Le cycle benzénique est représenté par \texttt{*6} (cercle aromatique).}
\end{popfigure}

\begin{exoresolu}[Masse de motif et degré de polymérisation du PET]
Le motif du PET a pour formule brute \ce{C10H8O4}. Une bouteille en plastique est constituée de macromolécules de masse molaire moyenne $M = \SI{38400}{g.mol^{-1}}$. Données : $M(\ce{C}) = 12$, $M(\ce{H}) = 1$, $M(\ce{O}) = \SI{16}{g.mol^{-1}}$. Calculer la masse molaire du motif, puis le degré de polymérisation moyen $n$ (nombre de motifs par chaîne).
\tcblower
\textbf{1. Masse molaire du motif \ce{C10H8O4} :}
\[ M_{\text{motif}} = 10 \times 12 + 8 \times 1 + 4 \times 16 = 120 + 8 + 64 = \SI{192}{g.mol^{-1}}. \]
\textbf{2. Degré de polymérisation :} c'est le nombre de motifs par macromolécule :
\[ n = \frac{M}{M_{\text{motif}}} = \frac{38400}{192} = 200. \]
Chaque chaîne de PET de cette bouteille contient donc en moyenne \textbf{200 motifs}, c'est-à-dire qu'elle résulte de la condensation d'environ 200 molécules d'acide téréphtalique et 200 molécules d'éthylène glycol, avec élimination d'environ 200 molécules d'eau par chaîne.
\end{exoresolu}

\courssub{Polycondensation ou polyaddition ?}

En Première, tu as rencontré la \cle{polyaddition} : des monomères à \emph{double liaison} \ce{C=C} (comme l'éthylène \ce{CH2=CH2}) s'additionnent les uns aux autres \textbf{sans éliminer aucune molécule} ; tous les atomes des monomères se retrouvent dans le polymère (polyéthylène, PVC). La polycondensation, elle, exige des monomères \emph{bifonctionnels} et s'accompagne \emph{toujours} du départ d'une petite molécule.

\begin{popfigure}
\begin{center}
\begin{tikzpicture}[scale=0.85, every node/.style={transform shape}]
% Polyaddition
\node[font=\bfseries, popblue] at (2.2,3.4) {Polyaddition};
\draw[rounded corners, draw=popblue, fill=popblueL] (0,1.6) rectangle (1.1,2.6) node[midway] {\ce{M}};
\draw[rounded corners, draw=popblue, fill=popblueL] (1.5,1.6) rectangle (2.6,2.6) node[midway] {\ce{M}};
\draw[rounded corners, draw=popblue, fill=popblueL] (3.0,1.6) rectangle (4.1,2.6) node[midway] {\ce{M}};
\draw[->, very thick, popdark] (4.4,2.1) -- (5.6,2.1);
\draw[rounded corners, draw=popblue, fill=popblue!15] (5.9,1.6) rectangle (10.1,2.6);
\node at (8.0,2.1) {\ce{-M-M-M-}${}_n$};
\node[align=center, popmuted] at (2.2,0.9) {monomère à double liaison\\ \textbf{aucun sous-produit}};
% Polycondensation
\node[font=\bfseries, poporange] at (2.2,-0.9) {Polycondensation};
\draw[rounded corners, draw=poporange, fill=poporangeL] (0,-2.2) rectangle (1.1,-1.2) node[midway] {\ce{A-A}};
\draw[rounded corners, draw=popgreen, fill=popgreenL] (1.5,-2.2) rectangle (2.6,-1.2) node[midway] {\ce{B-B}};
\draw[rounded corners, draw=poporange, fill=poporangeL] (3.0,-2.2) rectangle (4.1,-1.2) node[midway] {\ce{A-A}};
\draw[->, very thick, popdark] (4.4,-1.7) -- (5.6,-1.7);
\draw[rounded corners, draw=poppurple, fill=poppurpleL] (5.9,-2.2) rectangle (9.0,-1.2);
\node at (7.45,-1.7) {\ce{-A-B-A-B-}${}_n$};
\node[poppink, font=\bfseries] at (9.9,-1.7) {$+\ n\,\ce{H2O}$};
\node[align=center, popmuted] at (2.2,-2.9) {deux monomères bifonctionnels\\ \textbf{élimination d'une petite molécule}};
\end{tikzpicture}
\end{center}
\legende{Comparaison : en polyaddition, les monomères s'assemblent sans perte d'atomes ; en polycondensation, deux monomères bifonctionnels alternent dans la chaîne et chaque liaison formée libère une petite molécule (eau ou \ce{HCl}).}
\end{popfigure}

\begin{center}
\begin{tabularx}{\linewidth}{|l|X|X|}
\hline
\rowcolor{popblueL} \textbf{Critère} & \textbf{Polyaddition} & \textbf{Polycondensation} \\
\hline
Monomère & une double liaison \ce{C=C} & deux fonctions réactives \\
\hline
Sous-produit & aucun & \ce{H2O} ou \ce{HCl} éliminés \\
\hline
Masse du polymère & somme des masses des monomères & inférieure (perte des petites molécules) \\
\hline
Exemples & polyéthylène, PVC & nylon 6,6, tergal (PET) \\
\hline
\end{tabularx}
\end{center}

\courssub{Propriétés et usages des polymères de condensation}

Les chaînes de polyamides et de polyesters sont longues, régulières, et reliées entre elles par de nombreuses \cle{liaisons hydrogène} (entre groupes \ce{N-H} et \ce{C=O} pour le nylon). Il en résulte des matériaux à la fois \textbf{solides, souples et résistants à l'usure} :
\begin{itemize}
\item \textbf{Nylon 6,6} : fibres textiles (bas, vêtements de sport), cordes, fils de pêche très utilisés sur les côtes de Kribi et Limbé, brosses, pièces techniques ;
\item \textbf{Tergal (PET)} : fibres pour tissus (souvent mélangé au coton), bouteilles plastiques pour l'eau et les boissons, films d'emballage.
\end{itemize}

\begin{savaistu}[Du laboratoire aux uniformes camerounais]
Le nom « nylon » viendrait du lancement simultané de la fibre en 1938 à \textbf{N}ew \textbf{Y}ork et à \textbf{Lon}dres : NY + Lon = nylon ! Inventé par Wallace Carothers chez DuPont, il fut d'abord utilisé pour les bas des femmes, puis pour les parachutes pendant la Seconde Guerre mondiale. Aujourd'hui, le tergal, infroissable et bon marché, compose la plupart des uniformes scolaires portés chaque matin par les élèves du Cameroun : la polycondensation est littéralement sur ton dos.
\end{savaistu}

\begin{exercice}[À toi de jouer]
On fait réagir l'acide hexanedioïque \ce{HOOC-(CH2)4-COOH} avec l'octane-1,8-diamine \ce{H2N-(CH2)8-NH2}.
\begin{enumerate}
\item Écrire l'équation-bilan de la polycondensation et encadrer le motif.
\item Comment se nomme ce nylon ?
\item Calculer la masse molaire du motif. Données : $M(\ce{C}) = 12$, $M(\ce{H}) = 1$, $M(\ce{O}) = 16$, $M(\ce{N}) = \SI{14}{g.mol^{-1}}$.
\end{enumerate}
\end{exercice}

\begin{corrige}[Exercice 1]
\textbf{1.} $n\,\ce{H2N-(CH2)8-NH2} + n\,\ce{HOOC-(CH2)4-COOH} \ce{->[\Delta]} \ce{[-NH-(CH2)8-NH-CO-(CH2)4-CO-]}_n + n\,\ce{H2O}$.
\textbf{2.} Diamine à 8 C, diacide à 6 C : c'est le \textbf{nylon 8,6} (la diamine est citée en premier).
\textbf{3.} Le motif \ce{[-NH-(CH2)8-NH-CO-(CH2)4-CO-]} a pour formule brute \ce{C14H26N2O2} :
\[ M = 14 \times 12 + 26 \times 1 + 2 \times 14 + 2 \times 16 = 168 + 26 + 28 + 32 = \SI{254}{g.mol^{-1}}. \]
\end{corrige}

\begin{aretenir}[L'essentiel de la section]
\begin{itemize}
\item Un \textbf{amide} \ce{R-CO-NH2} s'obtient en chauffant le sel d'ammonium \ce{R-COO^{-}NH4+} (déshydratation), ou rapidement et totalement par action d'un chlorure d'acyle sur \ce{NH3} ou une amine ; on le nomme en « ...amide » (\emph{N}-alkyl si l'azote est substitué).
\item La \textbf{polycondensation} assemble des monomères \emph{bifonctionnels} avec élimination de petites molécules (\ce{H2O}, \ce{HCl}) : ne jamais oublier les $n\,\ce{H2O}$ dans l'équation.
\item \textbf{Nylon 6,6} = hexane-1,6-diamine (6 C) + acide hexanedioïque (6 C) ; les chiffres donnent les nombres de carbone de la diamine puis du diacide.
\item \textbf{Tergal (PET)} = acide téréphtalique + éthane-1,2-diol ; c'est un polyester à liaisons \ce{-CO-O-}.
\item Polycondensation (avec sous-produit, monomères bifonctionnels) $\neq$ polyaddition (sans sous-produit, doubles liaisons).
\end{itemize}
\end{aretenir}

\newpage

\coursec{Travaux pratiques}

\courssub{Objectifs du TP}

À l'issue de cette séance de travaux pratiques, tu seras capable de :

\begin{itemize}
\item réaliser la \cle{saponification} d'un corps gras naturel (l'huile de palme, mélange de triglycérides) par la soude concentrée, à chaud ;
\item isoler le savon formé par la technique du \cle{relargage} dans une solution saturée de chlorure de sodium ;
\item vérifier expérimentalement le \cle{caractère basique} d'une solution de savon (pH $> 7$) ;
\item mettre en évidence le \cle{pouvoir moussant} du savon et l'influence de la dureté de l'eau (eau distillée contre eau de puits) ;
\item relier chaque observation macroscopique à l'équation-bilan de la réaction étudiée en cours.
\end{itemize}

Ce TP illustre directement la situation de la famille « fabrication d'un savon » : il reproduit, avec du matériel de lycée, le procédé utilisé depuis des générations dans les savonneries artisanales de nos villages, où l'huile de palme rouge est chauffée avec une lessive de cendres ou de soude.

\courssub{Matériel et produits}

\begin{center}
\begin{tabularx}{\linewidth}{|l|X|}
\hline
\textbf{Matériel} & \textbf{Produits} \\
\hline
2 béchers de \SI{250}{mL} & Huile de palme du marché : \SI{20}{mL} \\
1 bécher de \SI{500}{mL} (bain-marie) & Soude \ce{NaOH} à environ \SI{6}{mol.L^{-1}} : \SI{20}{mL} \\
Agitateur en verre & Solution saturée de \ce{NaCl} (sel de cuisine) : \SI{50}{mL} \\
Chauffe-ballon ou réchaud à gaz & Eau distillée \\
Éprouvette graduée de \SI{50}{mL} & Eau de puits (eau dure) \\
Tamis fin ou entonnoir Büchner + fiole à vide & Papier pH \\
Thermomètre (facultatif) & Gants de ménage, lunettes de protection, blouse \\
2 tubes à essai + bouchons & Spatule, cristallisoir \\
\hline
\end{tabularx}
\end{center}

\textbf{Alternatives locales.} Si le laboratoire ne dispose pas de soude en pastilles, on peut utiliser une lessive concentrée de cendres de bois filtrée (riche en potasse \ce{KOH}) : le savon obtenu sera plus mou, c'est le savon noir traditionnel. À défaut de Büchner, un tamis de cuisine propre recouvert d'un tissu fin (pagne à trame serrée lavé) convient pour la filtration. Un réchaud à gaz ou même un feu de bois maîtrisé remplace le chauffe-ballon.

\courssub{Sécurité}

\begin{attention}[La soude concentrée est corrosive]
La solution de soude à \SI{6}{mol.L^{-1}} est \textbf{corrosive} (pictogramme : deux gouttes attaquant une main et une barre métallique — danger de brûlures chimiques graves). Gestes obligatoires :
\begin{itemize}
\item porter en permanence \textbf{gants, lunettes et blouse} ; attacher les cheveux longs ;
\item ne jamais pipetter la soude à la bouche : utiliser une éprouvette graduée ;
\item verser la soude \emph{lentement} dans le bécher, jamais l'inverse avec de l'eau chaude ;
\item en cas de contact avec la peau : rincer abondamment à l'eau froide pendant plusieurs minutes, puis prévenir l'enseignant ;
\item le chauffage se fait \textbf{au bain-marie} uniquement : l'huile chaude et la soude ne doivent jamais être chauffées à flamme directe (projections, risque d'emballement) ;
\item ne pas goûter, ne pas sentir directement les vapeurs au-dessus du bécher.
\end{itemize}
\end{attention}

\courssub{Schéma du dispositif}

\begin{popfigure}
\begin{center}
\begin{tikzpicture}[scale=0.95, every node/.style={font=\small}]
% Bain-marie (grand bécher)
\draw[thick, popdark] (-2.6,0) -- (-2.6,3.2) -- (-2.4,3.4);
\draw[thick, popdark] (2.6,0) -- (2.6,3.2) -- (2.4,3.4);
\draw[thick, popdark] (-2.6,0) -- (2.6,0);
% Eau du bain-marie
\fill[popblue!25] (-2.55,0.05) rectangle (2.55,1.6);
\draw[popblue, thick] (-2.55,1.6) .. controls (-1.5,1.7) and (-0.5,1.5) .. (0.3,1.6) .. controls (1.2,1.7) and (2.0,1.5) .. (2.55,1.6);
\node[popblue!70!black] at (1.9,0.6) {eau};
% Bécher intérieur
\draw[thick, popdark] (-1.4,0.35) -- (-1.4,4.0) -- (-1.2,4.2);
\draw[thick, popdark] (1.4,0.35) -- (1.4,4.0) -- (1.2,4.2);
\draw[thick, popdark] (-1.4,0.35) -- (1.4,0.35);
% Mélange huile + soude
\fill[poporange!45] (-1.35,0.4) rectangle (1.35,2.3);
\draw[poporange!80!black] (-1.35,2.3) .. controls (-0.6,2.4) and (0.6,2.2) .. (1.35,2.3);
\node[poporange!80!black, align=center] at (0,1.3) {huile de palme\\+ soude};
% Agitateur en verre
\draw[very thick, popgreen!60!black] (0.7,5.2) -- (0.15,1.0);
\node[popgreen!60!black, anchor=south] at (0.75,5.2) {agitateur};
% Thermomètre
\draw[thick, popdark] (-0.9,4.8) -- (-0.55,1.2);
\fill[popink] (-0.55,1.2) circle (0.09);
\node[anchor=south, popdark] at (-0.9,4.8) {thermomètre};
% Flamme / chauffage
\draw[popink, thick, decorate, decoration={coil, amplitude=2.5pt, segment length=6pt}] (-1.6,-0.5) -- (1.6,-0.5);
\node[popink] at (0,-1.0) {chauffage doux (bain-marie, $\approx$ \SI{80}{\celsius})};
% Vapeur
\draw[popmuted, decorate, decoration={snake, amplitude=1.5pt}] (-0.3,4.3) -- (-0.3,5.1);
\draw[popmuted, decorate, decoration={snake, amplitude=1.5pt}] (1.6,4.0) -- (1.6,4.8);
\end{tikzpicture}
\end{center}
\legende{Dispositif de saponification au bain-marie : le bécher contenant l'huile de palme et la soude est placé dans l'eau chaude ; on agite en permanence avec l'agitateur en verre.}
\end{popfigure}

\courssub{Structure du tripalmitate de glycéryle (palmitine)}

\begin{popfigure}
\begin{center}
\chemfig{
CH_2(-[0]O-C(=[2]O)-[6]C_{15}H_{31})
-CH(-[0]O-C(=[2]O)-[6]C_{15}H_{31})
-CH_2(-[0]O-C(=[2]O)-[6]C_{15}H_{31})
}
\end{center}
\legende{Formule semi-développée du tripalmitate de glycéryle (triglycéride majoritaire de l'huile de palme). Les trois chaînes latérales \ce{C15H31-} proviennent de l'acide palmitique \ce{C15H31COOH}.}
\end{popfigure}

\courssub{Protocole}

\begin{methode}[Fabrication et isolement du savon]
\begin{enumerate}
\item \textbf{Préparation du mélange.} Verser \SI{20}{mL} d'huile de palme dans un bécher de \SI{250}{mL}. Ajouter lentement, en agitant, \SI{20}{mL} de soude à \SI{6}{mol.L^{-1}}.
\item \textbf{Chauffage.} Placer le bécher dans le bain-marie maintenu aux environs de \SI{80}{\celsius}. Chauffer doucement pendant environ \SI{20}{min} en \textbf{agitant sans interruption} avec l'agitateur en verre, jusqu'à obtention d'une pâte épaisse et homogène, de couleur orangée. Si la pâte devient trop sèche, ajouter quelques millilitres d'eau distillée chaude.
\item \textbf{Test de fin de réaction (facultatif).} Prélever une goutte de pâte, la déposer dans un tube avec de l'eau distillée et agiter : si aucune gouttelette d'huile ne persiste en surface, la saponification est terminée.
\item \textbf{Relargage.} Verser la pâte chaude dans un bécher contenant \SI{50}{mL} de solution saturée de \ce{NaCl}. Agiter, puis laisser refroidir : le savon, \textbf{insoluble dans l'eau salée}, précipite et surnage.
\item \textbf{Filtration et lavage.} Filtrer sur Büchner (ou sur tamis recouvert d'un tissu fin). Laver le résidu deux fois avec un peu d'eau salée \emph{froide} (jamais d'eau pure, qui redissoudrait une partie du savon). Presser le savon entre les parois pour l'essorer.
\item \textbf{Séchage.} Étaler le savon sur du papier absorbant et laisser sécher à l'air libre jusqu'à la séance suivante.
\item \textbf{Tests du savon.} Dissoudre un fragment de savon (environ \SI{1}{g}) dans \SI{20}{mL} d'eau distillée :
\begin{itemize}
\item déposer une goutte de la solution sur du papier pH, noter la valeur ;
\item répartir la solution dans deux tubes à essai ; ajouter de l'eau distillée dans le tube 1 et de l'eau de puits dans le tube 2 ; boucher et agiter chaque tube 10 fois de la même façon ; comparer les hauteurs de mousse.
\end{itemize}
\end{enumerate}
\end{methode}

\courssub{Observations et résultats attendus}

\begin{center}
\begin{tabularx}{\linewidth}{|X|X|}
\hline
\textbf{Étape / test} & \textbf{Observation attendue} \\
\hline
Chauffage huile + soude & Le mélange, d'abord hétérogène (deux phases), devient une pâte épaisse orangée au bout de 15 à \SI{20}{min} \\
Relargage dans l'eau salée saturée & Un précipité granuleux de savon apparaît et surnage ; le liquide surnageant (eau, glycérol, excès de soude) est limpide \\
Aspect du savon sec & Solide cassant, jaune-orangé (coloré par les caroténoïdes de l'huile de palme), d'odeur caractéristique \\
pH de la solution de savon & pH compris entre 9 et 10 : le savon est \textbf{basique} \\
Mousse dans l'eau distillée & Mousse abondante et stable (hauteur typique : 3 à \SI{4}{cm}) \\
Mousse dans l'eau de puits & Mousse faible (moins de \SI{1}{cm}) avec formation éventuelle de grumeaux blanchâtres \\
\hline
\end{tabularx}
\end{center}

\courssub{Exploitation}

\begin{exoresolu}[Questions d'exploitation du TP]
\textbf{Q1.} Écrire l'équation-bilan de la saponification de la palmitine (tripalmitate de glycéryle, triglycéride majoritaire de l'huile de palme) par la soude. Nommer les produits obtenus.

\textbf{Q2.} Pourquoi la saponification est-elle réalisée à chaud et sous agitation constante ?

\textbf{Q3.} Quel est le rôle du relargage ? Pourquoi lave-t-on le savon à l'eau \emph{salée froide} et non à l'eau distillée ?

\textbf{Q4.} La solution de savon a un pH voisin de 9,5. Écrire l'équation de la réaction qui justifie ce caractère basique, en utilisant l'ion palmitate \ce{C15H31COO-}.

\textbf{Q5.} L'eau de puits contient des ions \ce{Ca^2+} et \ce{Mg^2+}. Proposer une explication à la faible production de mousse dans le tube 2.

\textbf{Q6.} Pourquoi le glycérol formé ne se trouve-t-il pas dans le savon isolé ?
\tcblower
\textbf{R1.} L'équation-bilan s'écrit :
\[ \ce{C3H5(COOC15H31)3 + 3 NaOH ->[chaleur] 3 C15H31COONa + C3H5(OH)3} \]
Les produits sont le \textbf{palmitate de sodium} (le savon) et le \textbf{glycérol} (propane-1,2,3-triol). La réaction est lente mais \textbf{totale}.

\textbf{R2.} La saponification est une réaction \textbf{lente} à température ambiante : la chaleur (bain-marie à \SI{80}{\celsius}) l'accélère. L'huile et la soude ne sont pas miscibles ; l'agitation permanente augmente la surface de contact entre les deux phases et homogénéise le milieu, ce qui permet à la réaction d'atteindre son terme.

\textbf{R3.} Le savon est \textbf{soluble dans l'eau} mais \textbf{insoluble dans l'eau salée saturée} : l'excès d'ions \ce{Na+} et \ce{Cl-} diminue fortement sa solubilité (effet de \cle{relargage} ou \cle{salage}), il précipite donc et surnage, ce qui permet de le séparer du glycérol et de la soude restant en solution. Le lavage à l'eau salée froide élimine les traces de soude et de glycérol sans dissoudre le savon ; un lavage à l'eau distillée entraînerait une perte de produit.

\textbf{R4.} L'ion palmitate est la base conjuguée de l'acide palmitique (acide faible). Il réagit avec l'eau selon l'hydrolyse de base :
\[ \ce{C15H31COO- + H2O <=> C15H31COOH + HO-} \]
La formation d'ions hydroxyde \ce{HO-} rend la solution basique, d'où un pH supérieur à 7 (ici 9 à 10).

\textbf{R5.} Les ions \ce{Ca^2+} et \ce{Mg^2+} des eaux dures forment avec les ions palmitate des sels de calcium et de magnésium \textbf{insolubles} (les grumeaux blanchâtres observés, appelés « savons calciques » ou « savons magnésiens ») :
\[ \ce{2 C15H31COO- + Ca^2+ -> (C15H31COO)2Ca v} \]
Une partie du savon est ainsi « consommée » et ne peut plus mousser : le pouvoir détergent diminue dans les eaux dures.

\textbf{R6.} Le glycérol est \textbf{très soluble dans l'eau}, y compris dans l'eau salée : il reste dans le filtrat lors du relargage et de la filtration. C'est d'ailleurs un sous-produit valorisable (pharmacie, cosmétique) dans les savonneries industrielles.
\end{exoresolu}

\courssub{Conclusion}

Au cours de ce TP, nous avons transformé un corps gras naturel, l'huile de palme, en savon par action de la soude à chaud : c'est la \cle{saponification}, réaction \textbf{lente mais totale} qui produit un carboxylate de sodium (le savon) et du glycérol. Le \cle{relargage} dans l'eau salée saturée a permis d'isoler le savon par précipitation, puis les tests ont confirmé ses deux propriétés essentielles : une solution de savon est \textbf{basique} (pH $\approx 9{,}5$) car l'ion carboxylate est une base faible, et le savon possède un \textbf{pouvoir moussant et détergent}, fortement réduit dans les eaux dures à cause de la précipitation des carboxylates de calcium et de magnésium. Ce protocole, réalisé avec du matériel simple et des produits du marché local, reproduit fidèlement le procédé ancestral des savonneries artisanales camerounaises : la chimie organique est bien une science vivante, au service du quotidien.

\newpage

\coursec{Méthodes et exercices résolus}

\courssub{Méthode 1 — Écrire la formule semi-développée d'un acide carboxylique}

\begin{methode}[Écrire la formule d'un acide carboxylique (saturé, insaturé, polyacide)]
\begin{enumerate}
\item Repérer le \cle{groupe carboxyle} \ce{-COOH} : il contient le carbone de rang 1, toujours placé en bout de chaîne.
\item Compter le nombre total d'atomes de carbone $n$ de la chaîne principale (le carbone du carboxyle est inclus).
\item Pour un \cle{monoacide saturé} : la chaîne carbonée est \ce{C_{n-1}H_{2n-1}-}, d'où la formule générale \ce{C_nH_{2n}O2} ou \ce{C_{n-1}H_{2n-1}-COOH}. Exemple : $n=3$ donne \ce{CH3-CH2-COOH}.
\item Pour un \cle{acide insaturé} : chaque double liaison \ce{C=C} retire 2 H. Placer la double liaison au rang indiqué (compté depuis le carboxyle).
\item Pour un \cle{polyacide} : placer un groupe \ce{-COOH} à chaque extrémité (diacide) et écrire \ce{HOOC-(CH2)_m-COOH}.
\item Vérifier la tétravalence de chaque carbone et compter les H pour contrôle.
\end{enumerate}
\textbf{Réflexe :} le carbone du groupe carboxyle ne porte jamais d'hydrogène supplémentaire : il est lié à \ce{=O}, \ce{-OH} et à la chaîne \ce{R}.
\end{methode}

\begin{exoresolu}[Formules d'acides carboxyliques]
Écrire les formules semi-développées des acides suivants : (a) acide pentanoïque ; (b) acide prop-2-énoïque (acide acrylique) ; (c) acide hexanedioïque (acide adipique, utilisé dans la fabrication du nylon 6,6) ; (d) acide 2-méthylpropanoïque.
\tcblower
\textbf{(a) Acide pentanoïque :} 5 carbones au total, chaîne saturée. Le carboxyle porte le carbone 1 ; il reste 4 carbones saturés : \ce{C4H9-} soit \ce{CH3-CH2-CH2-CH2-}.
\[ \chemfig{CH_3-CH_2-CH_2-CH_2-C(=[1]O)-[7]OH} \]
Vérification : \ce{C5H10O2}, conforme à \ce{C_nH_{2n}O2} avec $n=5$.

\textbf{(b) Acide prop-2-énoïque :} 3 carbones, double liaison entre C2 et C3 (numérotation depuis \ce{-COOH}) :
\[ \chemfig{CH_2=CH-C(=[1]O)-[7]OH} \]
soit \ce{CH2=CH-COOH} (\ce{C3H4O2} : la double liaison a retiré 2 H par rapport à \ce{C3H6O2}).

\textbf{(c) Acide hexanedioïque :} diacide à 6 carbones, un \ce{-COOH} à chaque extrémité, donc 4 \ce{CH2} entre eux :
\[ \ce{HOOC-CH2-CH2-CH2-CH2-COOH} \quad \text{soit} \quad \ce{HOOC-(CH2)4-COOH} \]

\textbf{(d) Acide 2-méthylpropanoïque :} chaîne principale de 3 carbones avec un méthyle sur C2 :
\[ \chemfig{H_3C-CH(-[2]CH_3)-C(=[1]O)-[7]OH} \]
Conclusion : les quatre formules sont \ce{CH3(CH2)3COOH}, \ce{CH2=CH-COOH}, \ce{HOOC-(CH2)4-COOH} et \ce{(CH3)2CH-COOH}.
\end{exoresolu}

\courssub{Méthode 2 — Nommer les acides carboxyliques et leurs dérivés}

\begin{methode}[Nomenclature des acides carboxyliques et dérivés]
\begin{enumerate}
\item \textbf{Acide carboxylique :} identifier la plus longue chaîne contenant \ce{-COOH} ; nommer : \emph{acide} + nom de l'alcane dont le « e » final est remplacé par « \cle{oïque} ». Le carbone du carboxyle est le C1 (jamais précisé). Exemple : \ce{CH3-CH2-COOH} = acide propanoïque.
\item \textbf{Diacide :} suffixe « \cle{dioïque} » : \ce{HOOC-(CH2)4-COOH} = acide hexanedioïque.
\item \textbf{Chlorure d'acyle} \ce{R-COCl} : « chlorure de » + nom de l'acyle (acide ...oïque $\to$ ...\cle{anoyle}). Exemple : \ce{CH3-COCl} = chlorure d'éthanoyle.
\item \textbf{Anhydride d'acide} \ce{(R-CO)2O} : « anhydride » + nom de l'acide dont « acide » est remplacé par « anhydride ». Exemple : \ce{(CH3-CO)2O} = anhydride éthanoïque.
\item \textbf{Ester} \ce{R-CO-O-R$'$} : nom de l'acide avec « oïque » $\to$ « \cle{oate} » + « de » + nom du groupe alkyle \ce{R$'$} (chaîne issue de l'alcool). Exemple : \ce{CH3-COO-CH2-CH3} = éthanoate d'éthyle.
\item \textbf{Amide} \ce{R-CO-NR$'$R$''$} : nom de l'acide avec « oïque » $\to$ « \cle{amide} », précédé de N-alkyle si l'azote est substitué. Exemple : \ce{CH3-CO-NH-CH3} = N-méthyléthanamide.
\end{enumerate}
\textbf{Réflexe :} dans un ester, la partie « oate » vient de l'acide (contient le \ce{C=O}), la partie « alkyle » vient de l'alcool (liée à l'oxygène).
\end{methode}

\begin{exoresolu}[Nommer des acides et leurs dérivés]
Nommer les composés suivants :
(a) \ce{CH3-CH2-CH2-COOH} ; (b) \ce{CH3-CH2-COCl} ; (c) \ce{(CH3-CH2-CO)2O} ; (d) \ce{CH3-CH2-COO-CH3} ; (e) \ce{CH3-CO-NH2} ; (f) \ce{HOOC-CH2-CH2-COOH}.
\tcblower
\textbf{(a)} Chaîne de 4 carbones incluant le carboxyle : butane $\to$ \textbf{acide butanoïque}.

\textbf{(b)} Chlorure d'acyle dérivé de l'acide propanoïque (3 C) : \textbf{chlorure de propanoyle}.

\textbf{(c)} Anhydride symétrique dérivé de l'acide propanoïque : \textbf{anhydride propanoïque}.

\textbf{(d)} Ester : partie acyle \ce{CH3-CH2-CO-} (3 C, propanoate) et alkyle \ce{-CH3} (méthyle) : \textbf{propanoate de méthyle}.

\textbf{(e)} Amide non substitué dérivé de l'acide éthanoïque : \textbf{éthanamide}.

\textbf{(f)} Diacide à 4 carbones : \textbf{acide butanedioïque} (acide succinique).

Conclusion : (a) acide butanoïque ; (b) chlorure de propanoyle ; (c) anhydride propanoïque ; (d) propanoate de méthyle ; (e) éthanamide ; (f) acide butanedioïque.
\end{exoresolu}

\courssub{Méthode 3 — Expliquer les propriétés physiques par les liaisons hydrogène}

\begin{methode}[Justifier températures d'ébullition et solubilité]
\begin{enumerate}
\item Repérer dans la molécule la présence d'un H lié à O (groupe \ce{-OH} du carboxyle) : condition d'existence de \cle{liaisons hydrogène} intermoléculaires.
\item Préciser que deux molécules d'acide carboxylique s'associent en \cle{dimère} par deux liaisons hydrogène (même en phase vapeur) : la masse molaire « apparente » double.
\item Conséquence 1 : les températures de fusion et d'ébullition sont \textbf{plus élevées} que celles des alcools, aldéhydes ou alcanes de masse molaire voisine.
\item Conséquence 2 : les acides à \textbf{courte chaîne} ($n \leq 4$) sont \textbf{solubles dans l'eau} car ils forment des liaisons hydrogène avec les molécules d'eau ; la solubilité \textbf{diminue} quand la chaîne carbonée (hydrophobe) s'allonge.
\item Comparer deux composés : à chaîne égale, celui qui forme le plus de liaisons hydrogène a la température d'ébullition la plus élevée.
\end{enumerate}
\textbf{Schéma du dimère :}
\end{methode}

\begin{center}
\begin{tikzpicture}[scale=0.9]
\draw (0,0) node {\chemfig{R-C(=[1]O)-[7]O-H}};
\draw (5.2,0) node {\chemfig{H-O(-[1]C(=[3]O)-R)-[7]}};
\draw[dashed,thick,popblue] (1.55,-0.35) -- (4.0,-0.35);
\draw[dashed,thick,popblue] (1.55,0.75) -- (4.0,0.75);
\draw (2.8,-1.2) node[popblue,font=\small] {liaisons hydrogène};
\end{tikzpicture}
\end{center}

\begin{exoresolu}[Comparaison de températures d'ébullition]
On donne les températures d'ébullition : propane \SI{-42}{\celsius} ; éthanal \SI{21}{\celsius} ; éthanol \SI{78}{\celsius} ; acide méthanoïque \SI{101}{\celsius}. Ces quatre composés ont des masses molaires voisines ($\SI{44}{g.mol^{-1}}$ à $\SI{46}{g.mol^{-1}}$). Expliquer cette évolution. Pourquoi l'acide méthanoïque est-il miscible à l'eau alors que l'acide stéarique \ce{C17H35COOH} ne l'est pas ?
\tcblower
\textbf{Étape 1 — Le propane} ne possède ni atome d'oxygène ni liaison polarisée : seules des forces de Van der Waals faibles agissent ; il bout à \SI{-42}{\celsius}.

\textbf{Étape 2 — L'éthanal} possède un groupe carbonyle polaire \ce{C=O} (interactions dipôle-dipôle) mais pas de H lié à O : pas de liaison hydrogène ; $T_{\text{éb}} = \SI{21}{\celsius}$, plus élevée que le propane.

\textbf{Étape 3 — L'éthanol} possède un groupe \ce{-OH} : ses molécules s'associent par liaisons hydrogène ; $T_{\text{éb}} = \SI{78}{\celsius}$.

\textbf{Étape 4 — L'acide méthanoïque} possède le groupe carboxyle (donneur \emph{et} accepteur de liaisons hydrogène) : il forme des \textbf{dimères} stables par deux liaisons hydrogène, ce qui double sa masse apparente ; d'où $T_{\text{éb}} = \SI{101}{\celsius}$, la plus élevée.

\textbf{Solubilité :} l'acide méthanoïque possède une chaîne carbonée quasi nulle ; son groupe carboxyle forme des liaisons hydrogène avec l'eau : il est miscible. L'acide stéarique possède une longue chaîne hydrocarbonée \ce{C17H35-} hydrophobe qui l'emporte sur le caractère hydrophile du \ce{-COOH} : il est insoluble dans l'eau.

Conclusion : les liaisons hydrogène (dimères) expliquent les températures d'ébullition élevées des acides carboxyliques, et la solubilité décroît quand la chaîne carbonée s'allonge.
\end{exoresolu}

\courssub{Méthode 4 — Écrire l'action de l'eau et la formation des dérivés (chlorure d'acyle, anhydride)}

\begin{methode}[Équations-bilans des acides carboxyliques]
\begin{enumerate}
\item \textbf{Action de l'eau (réaction acido-basique) :} l'acide carboxylique est un acide \textbf{faible} ; écrire l'équilibre avec une double flèche :
\[ \ce{RCOOH + H2O <=> RCOO- + H3O+} \]
et citer le couple \ce{RCOOH}/\ce{RCOO-}.
\item \textbf{Chlorure d'acyle} (réactions \textbf{totales}) :
\[ \ce{RCOOH + PCl5 -> RCOCl + POCl3 + HCl} \]
\[ \ce{RCOOH + SOCl2 -> RCOCl + SO2 + HCl} \]
\ce{SOCl2} est préféré car les sous-produits \ce{SO2} et \ce{HCl} sont gazeux (faciles à éliminer).
\item \textbf{Anhydride d'acide} par déshydratation entre deux molécules d'acide en présence d'un déshydratant (\ce{P4O10}) ou par chauffage :
\[ \ce{2 RCOOH ->[P4O10] (RCO)2O + H2O} \]
\item Vérifier la conservation des atomes et préciser l'état physique des produits gazeux (\ce{HCl}, \ce{SO2}).
\end{enumerate}
\textbf{Réflexe :} dans les trois cas, le groupe \ce{-OH} du carboxyle est remplacé par \ce{-Cl}, \ce{-O-CO-R} (ou \ce{-OR} pour l'ester) : le dérivé est toujours \ce{R-CO-X}.
\end{methode}

\begin{exoresolu}[Synthèse d'un chlorure d'acyle]
Un laboratoire de Yaoundé prépare le chlorure d'éthanoyle à partir de $m = \SI{12,0}{g}$ d'acide éthanoïque et d'un excès de chlorure de thionyle \ce{SOCl2}. (1) Écrire l'équation-bilan. (2) Calculer la masse maximale de chlorure d'éthanoyle obtenue. (3) Pourquoi utilise-t-on \ce{SOCl2} plutôt que \ce{PCl5} ? Données : $M(\ce{CH3COOH}) = \SI{60,0}{g.mol^{-1}}$ ; $M(\ce{CH3COCl}) = \SI{78,5}{g.mol^{-1}}$.
\tcblower
\textbf{(1) Équation-bilan :}
\[ \ce{CH3COOH + SOCl2 -> CH3COCl + SO2 + HCl} \]
Vérification : à gauche C2, H4, O2, S1, Cl2 ; à droite C2, H4, O2, S1, Cl2 : l'équation est équilibrée.

\textbf{(2) Tableau d'avancement} (\ce{SOCl2} en excès, l'acide est limitant) :
\begin{center}
\begin{tabularx}{\linewidth}{|l|X|X|X|}
\hline
\rowcolor{popblueL} État & \ce{CH3COOH} & \ce{SOCl2} & \ce{CH3COCl} \\
\hline
Initial (mol) & $n_0$ & excès & 0 \\
\hline
Final (mol) & $n_0 - x_{\max}=0$ & excès & $x_{\max}$ \\
\hline
\end{tabularx}
\end{center}
Quantité d'acide :
\[ n_0 = \frac{m}{M} = \frac{12,0}{60,0} = \SI{0,200}{mol} \]
Stœchiométrie 1:1 : $n(\ce{CH3COCl}) = x_{\max} = \SI{0,200}{mol}$.
\[ m(\ce{CH3COCl}) = n \times M = 0,200 \times 78,5 = \SI{15,7}{g} \]

\textbf{(3)} Avec \ce{SOCl2}, les sous-produits \ce{SO2} et \ce{HCl} sont \textbf{gazeux} : ils s'échappent du milieu, ce qui déplace l'équilibre et facilite la purification du chlorure d'éthanoyle ; avec \ce{PCl5}, \ce{POCl3} reste liquide et doit être séparé.

Conclusion : la masse maximale de chlorure d'éthanoyle est $\boxed{m = \SI{15,7}{g}}$.
\end{exoresolu}

\courssub{Méthode 5 — Synthèse des esters : estérification directe et voies rapides}

\begin{methode}[Choisir et écrire une synthèse d'ester]
\begin{enumerate}
\item \textbf{Estérification directe} (acide + alcool) :
\[ \ce{RCOOH + R$'$OH <=> RCOOR$'$ + H2O} \]
Réaction \textbf{lente}, \textbf{limitée} (équilibre) et \textbf{athermique} ; catalysée par \ce{H2SO4} ou \ce{H3O+}, accélérée par chauffage à reflux. Limite : $\approx \SI{67}{\%}$ (alcool primaire, mélange équimolaire).
\item \textbf{Voies rapides et totales} (à privilégier pour un rendement maximal) :
\[ \ce{RCOCl + R$'$OH -> RCOOR$'$ + HCl} \]
\[ \ce{(RCO)2O + R$'$OH -> RCOOR$'$ + RCOOH} \]
Ces réactions sont \textbf{rapides} et \textbf{totales} ; l'anhydride est moins violent que le chlorure (pas de \ce{HCl} corrosif dégagé).
\item Pour un calcul de rendement : $\eta = \dfrac{n_{\text{ester obtenu}}}{n_{\text{ester théorique}}} \times 100$.
\item Identifier le réactif limitant avant tout calcul (tableau d'avancement).
\end{enumerate}
\textbf{Réflexe :} dans l'ester formé, \ce{R-CO-} vient de l'acide (ou de son dérivé) et \ce{-O-R$'$} vient de l'alcool.
\end{methode}

\begin{exoresolu}[Rendement d'une estérification]
Pour parfumer des bonbons, on synthétise l'éthanoate d'isoamyle (odeur de banane). On chauffe à reflux $n_1 = \SI{0,10}{mol}$ d'acide éthanoïque avec $n_2 = \SI{0,10}{mol}$ de 3-méthylbutan-1-ol en présence d'acide sulfurique. À l'équilibre, on obtient $m = \SI{8,6}{g}$ d'ester. (1) Écrire l'équation-bilan. (2) Calculer le rendement. (3) Proposer deux modifications pour améliorer le rendement, puis une méthode de synthèse totale de cet ester. Donnée : $M(\text{ester}) = \SI{130}{g.mol^{-1}}$.
\tcblower
\textbf{(1) Équation-bilan :}
\[ \ce{CH3COOH + (CH3)2CH-CH2-CH2OH <=> CH3COO-CH2-CH2-CH(CH3)2 + H2O} \]

\textbf{(2) Tableau d'avancement} (mélange équimolaire) :
\begin{center}
\begin{tabularx}{\linewidth}{|l|X|X|X|X|}
\hline
\rowcolor{popblueL} État & \ce{CH3COOH} & alcool & ester & \ce{H2O} \\
\hline
Initial (mol) & 0,10 & 0,10 & 0 & 0 \\
\hline
Équilibre (mol) & $0{,}10 - x$ & $0{,}10 - x$ & $x$ & $x$ \\
\hline
\end{tabularx}
\end{center}
Quantité d'ester obtenue :
\[ n_{\text{ester}} = \frac{m}{M} = \frac{8,6}{130} = \SI{0,066}{mol} \]
Quantité théorique maximale : $x_{\max} = \SI{0,10}{mol}$. Rendement :
\[ \eta = \frac{0,066}{0,10} \times 100 = \SI{66}{\%} \]
C'est bien la limite attendue pour un alcool primaire en mélange équimolaire ($\approx \SI{67}{\%}$).

\textbf{(3)} Pour améliorer le rendement : (i) mettre l'un des réactifs (le moins cher, l'acide) \textbf{en excès} ; (ii) \textbf{éliminer l'eau} au fur et à mesure (déplacement d'équilibre). Pour une synthèse \textbf{totale} : faire réagir le chlorure d'éthanoyle ou l'anhydride éthanoïque sur le 3-méthylbutan-1-ol :
\[ \ce{(CH3CO)2O + (CH3)2CH-CH2-CH2OH -> CH3COO-CH2-CH2-CH(CH3)2 + CH3COOH} \]
Conclusion : le rendement de l'estérification directe est $\boxed{\eta = \SI{66}{\%}}$ ; la voie par anhydride permettrait un rendement proche de \SI{100}{\%}.
\end{exoresolu}

\courssub{Méthode 6 — Hydrolyse, saponification et fabrication du savon}

\begin{methode}[Écrire une hydrolyse ou une saponification]
\begin{enumerate}
\item \textbf{Hydrolyse d'un ester} (réaction inverse de l'estérification) : lente et limitée :
\[ \mathrm{RCOOR'} + \ce{H2O} \rightleftharpoons \mathrm{RCOOH} + \mathrm{R'OH} \]
\item \textbf{Saponification} (hydrolyse en milieu basique, \ce{NaOH} ou \ce{KOH}, à chaud) : réaction \textbf{lente mais totale} :
\[ \mathrm{RCOOR'} + \ce{Na+} + \ce{HO-} \rightarrow \mathrm{RCOO^-}\ce{Na+} + \mathrm{R'OH} \]
Le produit est le \textbf{carboxylate de sodium} (savon si R est une longue chaîne grasse), pas l'acide carboxylique.
\item \textbf{Fabrication du savon} à partir d'un triglycéride (huile de palme) : 3 moles de soude pour 1 mole de triglycéride :
\[ \text{triglycéride} + 3\,(\ce{Na+} + \ce{HO-}) \rightarrow 3\,\mathrm{RCOONa} + \text{glycérol} \]
\item Protocole : chauffer huile + soude concentrée, puis \textbf{relargage} (ajout de saumure, solution concentrée de \ce{NaCl}) : le savon, insoluble dans l'eau salée, précipite ; on filtre, on rince, on sèche.
\item Stœchiométrie : toujours vérifier les coefficients (3 pour un triglycéride).
\end{enumerate}
\end{methode}

\begin{exoresolu}[Fabrication artisanale de savon à l'huile de palme]
Une coopérative de savonnerie artisanale de la région du Littoral fait réagir $m = \SI{806}{g}$ de palmitine (tripalmitate de glycéryle, $M = \SI{806}{g.mol^{-1}}$) avec un excès de soude. (1) Écrire l'équation-bilan de la saponification. (2) Calculer la masse de savon (palmitate de sodium, $M = \SI{278}{g.mol^{-1}}$) et la masse de glycérol ($M = \SI{92}{g.mol^{-1}}$) obtenues, la réaction étant totale. (3) Quel est le rôle du relargage ?
\tcblower
\textbf{(1) Équation-bilan :}
\[ \ce{(C15H31COO)3C3H5 + 3 (Na+ + HO-) -> 3 (C15H31COO- + Na+) + C3H5(OH)3} \]

\textbf{(2) Quantité de triglycéride :}
\[ n_{\text{tri}} = \frac{m}{M} = \frac{806}{806} = \SI{1,00}{mol} \]
Tableau d'avancement (soude en excès) :
\begin{center}
\begin{tabularx}{\linewidth}{|l|X|X|X|X|}
\hline
\rowcolor{popblueL} État & palmitine & \ce{HO-} & savon & glycérol \\
\hline
Initial (mol) & 1,00 & excès & 0 & 0 \\
\hline
Final (mol) & 0 & excès & 3,00 & 1,00 \\
\hline
\end{tabularx}
\end{center}
D'après la stœchiométrie : $n_{\text{savon}} = 3 \times n_{\text{tri}} = \SI{3,00}{mol}$ et $n_{\text{glycérol}} = \SI{1,00}{mol}$.
\[ m_{\text{savon}} = 3,00 \times 278 = \SI{834}{g} \qquad m_{\text{glycérol}} = 1,00 \times 92 = \SI{92}{g} \]

\textbf{(3)} Le relargage consiste à verser le mélange dans une saumure concentrée : le savon y est \textbf{insoluble} et précipite en surface, ce qui permet de le séparer du glycérol et de la soude restante.

Conclusion : on obtient $\boxed{m_{\text{savon}} = \SI{834}{g}}$ de palmitate de sodium et $\boxed{m_{\text{glycérol}} = \SI{92}{g}}$.
\end{exoresolu}

\courssub{Méthode 7 — Amides et polycondensation : nylon 6,6 et tergal}

\begin{methode}[Écrire une amidification et une polycondensation]
\begin{enumerate}
\item \textbf{Formation d'un amide} : un dérivé d'acide (chlorure d'acyle ou anhydride) réagit sur l'ammoniac ou une amine ; réaction rapide et totale :
\[ \ce{RCOCl + NH3 -> RCO-NH2 + HCl} \]
\[ \ce{RCOCl + R$'$NH2 -> RCO-NH-R$'$ + HCl} \]
\item \textbf{Polycondensation} : réaction en chaîne entre \textbf{monomères bifonctionnels} avec élimination d'une \textbf{petite molécule} (\ce{H2O} ou \ce{HCl}).
\item \textbf{Nylon 6,6} (polyamide) : acide hexanedioïque (6 C) + hexane-1,6-diamine (6 C) :
\[ \ce{n HOOC-(CH2)4-COOH + n H2N-(CH2)6-NH2 -> [-CO-(CH2)4-CO-NH-(CH2)6-NH-]_n + 2n H2O} \]
Le motif contient la liaison \cle{amide} \ce{-CO-NH-}.
\item \textbf{Tergal (PET)} (polyester) : acide téréphtalique + éthane-1,2-diol :
\[ \ce{n HOOC-C6H4-COOH + n HO-CH2-CH2-OH -> [-CO-C6H4-CO-O-CH2-CH2-O-]_n + 2n H2O} \]
Le motif contient la liaison \cle{ester} \ce{-CO-O-}.
\item Toujours écrire le \textbf{motif} entre crochets avec l'indice $n$ et équilibrer la petite molécule éliminée.
\end{enumerate}
\end{methode}

\begin{exoresolu}[Le nylon 6,6]
Le nylon 6,6 est obtenu par polycondensation de l'acide hexanedioïque et de l'hexane-1,6-diamine. (1) Écrire les formules semi-développées des deux monomères. (2) Écrire l'équation-bilan de la polycondensation et encadrer le motif du polymère. (3) Justifier le nom « 6,6 ». (4) Calculer la masse de nylon obtenue à partir de $n = \SI{1,0}{mol}$ de chaque monomère (réaction totale), sachant que le motif a pour masse molaire $M_{\text{motif}} = \SI{226}{g.mol^{-1}}$.
\tcblower
\textbf{(1) Monomères :}
\[ \ce{HOOC-(CH2)4-COOH} \quad \text{(acide hexanedioïque)} \qquad \ce{H2N-(CH2)6-NH2} \quad \text{(hexane-1,6-diamine)} \]

\textbf{(2) Équation-bilan :}
\[ \ce{n HOOC-(CH2)4-COOH + n H2N-(CH2)6-NH2 -> [-CO-(CH2)4-CO-NH-(CH2)6-NH-]_n + 2n H2O} \]
Chaque jonction entre deux monomères crée une \textbf{liaison amide} \ce{-CO-NH-} avec élimination d'une molécule d'eau : c'est bien une polycondensation.

\textbf{(3)} Le nom « nylon 6,6 » indique le nombre d'atomes de carbone de chaque monomère : le diacide en comporte \textbf{6} et la diamine \textbf{6}.

\textbf{(4)} À partir de \SI{1,0}{mol} de chaque monomère, on forme \SI{1,0}{mol} de motifs (et $2 \times 1{,}0 = \SI{2,0}{mol}$ d'eau éliminée) :
\[ m_{\text{nylon}} = n_{\text{motif}} \times M_{\text{motif}} = 1,0 \times 226 = \SI{226}{g} \]
Vérification par conservation de la masse : $m_{\text{monomères}} = 146 + 116 = \SI{262}{g}$ ; $m_{\text{eau}} = 2{,}0 \times 18 = \SI{36}{g}$ ; $262 - 36 = \SI{226}{g}$ : cohérent.

Conclusion : on obtient $\boxed{m = \SI{226}{g}}$ de nylon 6,6, polymère dont le motif contient des liaisons amide.
\end{exoresolu}

\begin{attention}[Les 5 erreurs qui coûtent des points au Bac]
\begin{enumerate}
\item \textbf{Confondre hydrolyse et saponification.} L'hydrolyse d'un ester est lente et \emph{limitée} et redonne l'acide carboxylique ; la saponification (soude ou potasse, à chaud) est \emph{totale} et donne l'ion \textbf{carboxylate} \ce{RCOO-}, jamais l'acide \ce{RCOOH} en milieu basique.
\item \textbf{Écrire l'estérification avec une flèche simple.} La réaction acide + alcool est un \textbf{équilibre} : double flèche \ce{<=>} obligatoire. En revanche, les synthèses à partir du chlorure d'acyle ou de l'anhydride sont totales : flèche simple.
\item \textbf{Inverser les deux parties du nom d'un ester.} Dans « éthanoate d'éthyle », la partie « éthanoate » vient de l'\emph{acide} (celle qui porte le \ce{C=O}) et « éthyle » vient de l'\emph{alcool}. \ce{CH3COOCH2CH3} n'est pas de l'éthanoate de méthyle !
\item \textbf{Oublier les coefficients stœchiométriques du savon.} Un triglycéride consomme \textbf{3} moles de soude et produit \textbf{3} moles de savon et \textbf{1} mole de glycérol. Oublier le facteur 3 fausse tout le calcul de masse.
\item \textbf{Mal écrire le motif d'un polymère de polycondensation.} Le motif doit être écrit entre crochets avec l'indice $n$, et il faut équilibrer la petite molécule éliminée ($2n$ \ce{H2O} pour le nylon 6,6 et le tergal). Ne pas confondre polycondensation (élimination d'une molécule) et polyaddition (sans élimination).
\end{enumerate}
\end{attention}

\newpage

\coursec{Exercices}

\courssub{Je vérifie mes connaissances}

\begin{aretenir}[Rappels indispensables]
\begin{itemize}
\item \textbf{Formule générale} des acides carboxyliques saturés non ramifiés : \ce{C_nH_{2n}O_2} (ou \ce{C_{n-1}H_{2n-1}COOH}).
\item \textbf{Fonctions dérivées} : chlorure d'acyle (\ce{-COCl}), anhydride d'acide (\ce{-CO-O-CO-}), ester (\ce{-COO-}), amide (\ce{-CONH2}).
\item \textbf{Estérification directe} : lente, limitée, équilibre (\ce{<=>}), catalysée par \ce{H+}.
\item \textbf{Synthèses via dérivés réactifs} (chlorure d'acyle, anhydride) : rapides, totales (\ce{->}).
\item \textbf{Saponification} : hydrolyse basique d'un ester, rapide, totale, irréversible.
\item \textbf{Polycondensation} : élimination d'une petite molécule (eau, \ce{HCl}...), croissance par étapes.
\end{itemize}
\end{aretenir}

\begin{exercice}[QCM : à toi de choisir]
Pour chaque question, une seule réponse est exacte. Recopie la lettre correspondante.
\begin{enumerate}
\item La formule générale d'un acide carboxylique saturé à chaîne non ramifiée comportant $n$ atomes de carbone est :
\begin{enumerate}
\item[a)] \ce{C_{n}H_{2n}O_2} \quad b)] \ce{C_{n}H_{2n+2}O_2} \quad c)] \ce{C_{n}H_{2n-2}O_2}
\end{enumerate}
\item Le composé \ce{CH3-CO-Cl} est :
\begin{enumerate}
\item[a)] un anhydride d'acide \quad b)] un chlorure d'acyle \quad c)] un ester
\end{enumerate}
\item L'action de \ce{PCl5} sur l'acide éthanoïque donne :
\begin{enumerate}
\item[a)] l'anhydride éthanoïque \quad b)] l'éthanoate d'éthyle \quad c)] le chlorure d'éthanoyle
\end{enumerate}
\item La saponification est une réaction :
\begin{enumerate}
\item[a)] lente et limitée \quad b)] rapide et totale \quad c)] lente et totale
\end{enumerate}
\item Le nylon 6,6 est obtenu par :
\begin{enumerate}
\item[a)] polyaddition \quad b)] polycondensation \quad c)] estérification
\end{enumerate}
\end{enumerate}
\end{exercice}

\begin{exercice}[Vrai ou faux ? Justifie]
Réponds par \textbf{vrai} ou \textbf{faux}, puis justifie ta réponse en une ou deux phrases.
\begin{enumerate}
\item Les acides carboxyliques sont des acides forts dans l'eau.
\item Les acides carboxyliques ont des températures d'ébullition élevées à cause de la formation de dimères par liaisons hydrogène.
\item L'estérification directe entre un acide carboxylique et un alcool est une réaction rapide et totale.
\item La réaction entre un chlorure d'acyle et un alcool est rapide et totale.
\item L'hydrolyse d'un ester en milieu basique (soude) est une réaction limitée.
\item Le savon obtenu à partir de potasse \ce{KOH} est un savon dur.
\end{enumerate}
\end{exercice}

\begin{exercice}[Texte à trous]
Recopie et complète le texte suivant avec les mots ou expressions de la liste : \emph{carboxyle ; oxonium ; carboxylate ; déshydratant ; total ; limité ; glycérol ; chlorure d'acyle ; anhydride d'acide ; relargage}.

Le groupe fonctionnel des acides carboxyliques est le groupe \ldots, de formule \ce{-COOH}. Dans l'eau, un acide carboxylique cède un proton : il se forme des ions \ldots{} \ce{H3O+} et des ions \ldots{} \ce{RCOO-}. L'action de \ce{SOCl2} sur un acide carboxylique conduit à un \ldots, tandis que l'action d'un agent \ldots{} comme \ce{P4O10} conduit à un \ldots. L'estérification directe est une réaction lente et \ldots, alors que l'action d'un chlorure d'acyle sur un alcool est rapide et \ldots. La saponification d'un triglycéride par la soude produit du savon et du \ldots ; on sépare ensuite le savon par \ldots{} dans une solution concentrée de chlorure de sodium.
\end{exercice}

\begin{exercice}[Définitions]
Définis en une phrase claire et précise chacun des termes suivants :
\begin{enumerate}
\item Acide carboxylique ;
\item Ester ;
\item Saponification ;
\item Polycondensation ;
\item Amide.
\end{enumerate}
\end{exercice}

\courssub{Je m'entraîne}

\begin{methode}[Nomenclature des dérivés d'acides carboxyliques]
\begin{enumerate}
\item Identifier l'acide parent (plus longue chaîne portant le \ce{-COOH}).
\item Remplacer le suffixe \og -oïque \fg{} de l'acide par :
  \begin{itemize}
  \item \og -oyle \fg{} + \og chlorure \fg{} pour le chlorure d'acyle (\ce{-COCl}) ;
  \item \og -oïque \fg{} + \og anhydride \fg{} pour l'anhydride (\ce{-CO-O-CO-}) ;
  \item \og -oate de \fg{} + \emph{nom de l'alkyle} pour l'ester (\ce{-COOR}) ;
  \item \og -amide \fg{} pour l'amide (\ce{-CONH2}, \ce{-CONHR}, \ce{-CONR2}).
  \end{itemize}
\item Numéroter la chaîne à partir du carbone du groupe carbonyl.
\end{enumerate}
\end{methode}

\begin{exercice}[Nomenclature des dérivés d'acide]
On considère les quatre composés suivants :
\[ \ce{CH3-CO-Cl} \qquad \ce{(CH3-CO)2O} \qquad \ce{CH3-COO-CH2-CH3} \qquad \ce{CH3-CO-NH2} \]
\begin{enumerate}
\item Nomme chacun de ces composés et précise la famille de dérivés d'acide à laquelle il appartient.
\item Écris l'équation-bilan de la préparation de chacun d'eux à partir de l'acide éthanoïque (précise le réactif utilisé : \ce{SOCl2}, \ce{P4O10}, éthanol, ammoniac).
\item Pour la préparation de l'ester, propose une méthode rapide et totale et écris l'équation correspondante.
\end{enumerate}
\end{exercice}

\begin{exercice}[Températures d'ébullition et liaisons hydrogène]
On donne les températures d'ébullition de trois composés de masses molaires voisines :
\begin{center}
\begin{tabularx}{\linewidth}{|l|X|X|}
\hline
\rowcolor{popblueL} \textbf{Composé} & \textbf{Masse molaire (\si{g.mol^{-1}})} & \textbf{Température d'ébullition (\si{\celsius})} \\
\hline
Pentane & 72 & 36 \\
\hline
Butan-1-ol & 74 & 118 \\
\hline
Acide propanoïque & 74 & 141 \\
\hline
\end{tabularx}
\end{center}
\begin{enumerate}
\item Écris la formule semi-développée de chacun de ces trois composés.
\item Explique pourquoi le pentane a la température d'ébullition la plus faible.
\item Explique, en faisant intervenir les liaisons hydrogène et la formation de dimères, pourquoi l'acide propanoïque bout à une température plus élevée que le butan-1-ol.
\item Représente par un schéma le dimère de l'acide propanoïque en faisant apparaître les liaisons hydrogène.
\end{enumerate}
\end{exercice}

\begin{popfigure}
\centering
\chemfig{CH_3-CH_2-C(=[2]O)-[6]O-H}
\quad\textcolor{popblue}{\textbf{$\cdots$ 2 Liaisons H $\cdots$}}\quad
\chemfig{CH_3-CH_2-C(=[6]O)-[2]O-H}
\legende{Dimère cyclique de l'acide propanoïque : deux liaisons hydrogène \ce{O-H\cdots O} unissent les deux molécules.}
\end{popfigure}

\begin{exercice}[Synthèse d'un ester par le chlorure d'acyle]
On fait réagir \SI{0,10}{mol} de chlorure d'éthanoyle avec \SI{0,10}{mol} d'éthanol. La réaction est rapide et totale.
\begin{enumerate}
\item Écris l'équation-bilan de la réaction et nomme l'ester obtenu.
\item Construis le tableau d'avancement et détermine la quantité d'ester formée.
\item Calcule la masse d'ester obtenue.
\item La même quantité d'ester préparée par estérification directe entre l'acide éthanoïque et l'éthanol conduirait à un rendement de \SI{67}{\%}. Calcule la masse d'ester qu'on obtiendrait alors et conclus sur l'intérêt de la voie par chlorure d'acyle.
\end{enumerate}
\emph{Données :} $M(\ce{C}) = \SI{12}{g.mol^{-1}}$ ; $M(\ce{H}) = \SI{1}{g.mol^{-1}}$ ; $M(\ce{O}) = \SI{16}{g.mol^{-1}}$.
\end{exercice}

\begin{attention}[Pièges classiques]
\begin{itemize}
\item \textbf{Équilibrage \ce{PCl5} :} \ce{RCOOH + PCl5 -> RCOCl + POCl3 + HCl} (pas \ce{PCl3O}).
\item \textbf{Rendement estérification :} Ne pas confondre rendement (expérimental) et avancement maximal (théorique, $x_{\text{max}}$).
\item \textbf{Savon mou/dur :} Cation \ce{K+} $\rightarrow$ savon mou (potasse) ; Cation \ce{Na+} $\rightarrow$ savon dur (soude).
\item \textbf{Relargage :} Ajout de \ce{NaCl} saturé $\rightarrow$ diminution de la solubilité du savon (effet salin) $\rightarrow$ précipitation.
\end{itemize}
\end{attention}

\begin{exercice}[Fabrication artisanale du savon]
Dans une savonnerie artisanale de la région du Littoral, on prépare un détergent à partir d'huile de palme et d'une solution de potasse \ce{KOH}.
\begin{enumerate}
\item Quel est le nom de la réaction mise en jeu ? Donne ses caractéristiques (cinétique, bilan).
\item Sachant que l'huile de palme contient essentiellement des triglycérides (triesters du glycérol), écris l'équation-bilan générale de la réaction entre un triglycéride et la potasse.
\item Explique pourquoi le savon au potassium obtenu est un savon \emph{mou}, contrairement au savon au sodium qui est dur.
\item Après la réaction, on verse le mélange dans une solution concentrée de chlorure de sodium. Quel est le nom de cette opération ? Quel est son rôle ?
\end{enumerate}
\end{exercice}

\begin{savaistu}[Le savon au Cameroun]
La fabrication artisanale du savon (appelé \og savon noir \fg{} ou \og savon de ménage \fg{}) utilise l'huile de palme rouge (riche en acide palmitique et oléique) et la potasse issue de la cendre de bois ou de coques de cacao (\ce{K2CO3} $\xrightarrow{\text{chaux}}$ \ce{KOH}). Le \textbf{relargage} à l'eau de mer ou au \ce{NaCl} saturé permet de récupérer le savon mou, malléable, utilisé pour le lessivage. Les savonneries industrielles (ex. \emph{Savonnerie du Cameroun}) produisent du savon dur (\ce{Na+}) pour la toilette par relargage à la saumure.
\end{savaistu}

\courssub{Je me prépare au Bac}

\begin{exercice}[Type Bac : identification d'un acide carboxylique]
Un acide carboxylique saturé \textbf{A}, à chaîne carbonée non cyclique, a pour masse molaire $M(\textbf{A}) = \SI{74}{g.mol^{-1}}$.
\begin{enumerate}
\item Donne la formule générale d'un acide carboxylique saturé comportant $n$ atomes de carbone.
\item Détermine la valeur de $n$ et déduis-en la formule brute de \textbf{A}.
\item Écris les formules semi-développées de tous les isomères de \textbf{A} et nomme-les.
\item On dissout \SI{1,48}{g} de \textbf{A} dans de l'eau distillée pour obtenir \SI{200}{mL} de solution.
\begin{enumerate}
\item[a)] Calcule la concentration molaire $C$ de la solution obtenue.
\item[b)] Écris l'équation-bilan de la réaction de \textbf{A} avec l'eau et précise le couple acide/base mis en jeu.
\item[c)] Le pH de la solution vaut $2{,}9$. Calcule les concentrations en ions \ce{H3O+} et \ce{HO-}. L'acide \textbf{A} est-il fort ou faible ? Justifie.
\end{enumerate}
\item On fait réagir \textbf{A} avec le pentachlorure de phosphore \ce{PCl5}. Écris l'équation-bilan de la réaction et nomme le produit organique obtenu.
\item Le produit obtenu à la question 5 réagit avec le méthanol. Écris l'équation-bilan, nomme l'ester formé et donne deux caractéristiques de cette réaction.
\end{enumerate}
\emph{Données :} $M(\ce{C}) = \SI{12}{g.mol^{-1}}$ ; $M(\ce{H}) = \SI{1}{g.mol^{-1}}$ ; $M(\ce{O}) = \SI{16}{g.mol^{-1}}$ ; $K_e = \SI{1,0e-14}{}$ à \SI{25}{\celsius}.
\end{exercice}

\begin{exercice}[Type Bac : saponification de la tristéarine]
La tristéarine est un triglycéride présent dans les graisses animales ; sa formule brute est \ce{C57H110O6} et sa masse molaire vaut $M = \SI{890}{g.mol^{-1}}$. On la fait réagir avec un excès de soude (\ce{Na+ + HO-}) à chaud. Il se forme du stéarate de sodium \ce{C17H35COO-Na+} (savon) et du glycérol \ce{C3H8O3}.
\begin{enumerate}
\item Définis la saponification et donne deux caractéristiques de cette réaction.
\item Écris la formule semi-développée de la tristéarine à partir de celle du glycérol et de l'acide stéarique \ce{C17H35COOH}.
\item Écris l'équation-bilan de la saponification de la tristéarine par la soude.
\item On traite $m = \SI{890}{g}$ de tristéarine par la soude en excès.
\begin{enumerate}
\item[a)] Construis le tableau d'avancement de la réaction.
\item[b)] Calcule la quantité de matière de savon obtenue, puis la masse de savon produite.
\item[c)] Calcule la masse minimale de soude nécessaire pour saponifier cette masse de tristéarine.
\item[d)] Calcule la masse de glycérol formée. Citer une application industrielle du glycérol.
\end{enumerate}
\item Pourquoi utilise-t-on la soude en excès ? Quelle précaution faut-il prendre lors de la manipulation de la soude concentrée ?
\end{enumerate}
\emph{Données :} $M(\ce{C}) = \SI{12}{g.mol^{-1}}$ ; $M(\ce{H}) = \SI{1}{g.mol^{-1}}$ ; $M(\ce{O}) = \SI{16}{g.mol^{-1}}$ ; $M(\ce{Na}) = \SI{23}{g.mol^{-1}}$.
\end{exercice}

\begin{exercice}[Type Bac : le nylon 6,6 et le tergal]
Une usine textile souhaite produire deux fibres synthétiques : le nylon 6,6 (un polyamide) et le tergal (un polyester, encore appelé PET).
\begin{enumerate}
\item \textbf{Le nylon 6,6.} Il est obtenu par polycondensation de l'acide hexanedioïque (acide adipique) \ce{HOOC-(CH2)4-COOH} et de l'hexane-1,6-diamine \ce{H2N-(CH2)6-NH2}.
\begin{enumerate}
\item[a)] Définis la polycondensation. En quoi se distingue-t-elle d'une polyaddition ?
\item[b)] Écris l'équation-bilan de la réaction entre une molécule d'acide adipique et une molécule d'hexane-1,6-diamine. Quelle petite molécule est éliminée ?
\item[c)] Écris le motif du nylon 6,6 en l'encadrant. Justifie le nom \og 6,6 \fg{} donné à ce polymère.
\item[d)] On désire produire \SI{1,13}{kg} de nylon 6,6. Sachant que la masse molaire du motif vaut \SI{226}{g.mol^{-1}}, calcule le nombre de motifs contenus dans cette masse de polymère.
\end{enumerate}
\item \textbf{Le tergal.} Il est obtenu à partir de l'acide benzène-1,4-dicarboxylique (acide téréphtalique) \ce{HOOC-C6H4-COOH} et de l'éthane-1,2-diol (éthylèneglycol) \ce{HO-CH2-CH2-OH}.
\begin{enumerate}
\item[a)] Écris l'équation-bilan de la réaction entre les deux monomères.
\item[b)] Écris le motif du tergal en l'encadrant et entoure la fonction ester.
\item[c)] Pourquoi le tergal est-il classé parmi les polyesters et le nylon parmi les polyamides ?
\end{enumerate}
\item Cite une utilisation courante de chacun de ces deux polymères dans la vie quotidienne au Cameroun.
\end{enumerate}
\end{exercice}



\newpage

\coursec{Corrigés des exercices}

\coursec{Exercices corrigés — Acides carboxyliques et leurs dérivés}

\begin{corrige}[Exercice 1 — QCM : à toi de choisir]
\begin{enumerate}
\item \textbf{a)} $C_nH_{2n}O_2$. Un acide carboxylique saturé non ramifié dérive d'un alcane $C_nH_{2n+2}$ par remplacement de \ce{-CH3} par \ce{-COOH} : on perd 2 H.
\item \textbf{b)} un chlorure d'acyle. Le groupe \ce{-CO-Cl} caractérise cette famille.
\item \textbf{c)} le chlorure d'éthanoyle : \ce{CH3COOH + PCl5 -> CH3COCl + POCl3 + HCl}.
\item \textbf{b)} rapide et totale. La saponification est une hydrolyse basique irréversible.
\item \textbf{b)} polycondensation. Le nylon 6,6 se forme avec élimination d'eau entre un diacide et une diamine.
\end{enumerate}
\end{corrige}

\begin{corrige}[Exercice 2 — Vrai ou faux ? Justifie]
\begin{enumerate}
\item \textbf{Faux.} Les acides carboxyliques sont des acides \emph{faibles} : leur réaction avec l'eau est partielle ($K_a \approx 10^{-5}$), d'où l'équilibre $RCOOH + H_2O \rightleftharpoons RCOO^- + H_3O^+$.
\item \textbf{Vrai.} Deux molécules s'associent par deux liaisons hydrogène $O-H\cdots O$ pour former un dimère cyclique : la masse molaire effective double, il faut donc plus d'énergie pour vaporiser le liquide.
\item \textbf{Faux.} L'estérification directe est lente et limitée par un équilibre ($K \approx 4$ pour un alcool primaire, rendement $\approx \SI{67}{\%}$).
\item \textbf{Vrai.} La réaction $RCOCl + R'OH \rightarrow RCOOR' + HCl$ est rapide, totale et exothermique : c'est l'intérêt des dérivés réactifs.
\item \textbf{Faux.} L'hydrolyse basique (saponification) est totale : l'ion carboxylate formé est la base d'un acide faible, la réaction ne peut pas s'inverser.
\item \textbf{Faux.} La potasse \ce{KOH} donne un savon \emph{mou} (pâteux) ; c'est la soude \ce{NaOH} qui donne un savon \emph{dur}.
\end{enumerate}
\end{corrige}

\begin{corrige}[Exercice 3 — Texte à trous]
Le groupe fonctionnel des acides carboxyliques est le groupe \textbf{carboxyle}, de formule \ce{-COOH}. Dans l'eau, un acide carboxylique cède un proton : il se forme des ions \textbf{oxonium} \ce{H3O+} et des ions \textbf{carboxylate} \ce{RCOO-}. L'action de \ce{SOCl2} sur un acide carboxylique conduit à un \textbf{chlorure d'acyle}, tandis que l'action d'un agent \textbf{déshydratant} comme \ce{P4O10} conduit à un \textbf{anhydride d'acide}. L'estérification directe est une réaction lente et \textbf{limitée}, alors que l'action d'un chlorure d'acyle sur un alcool est rapide et \textbf{totale}. La saponification d'un triglycéride par la soude produit du savon et du \textbf{glycérol} ; on sépare ensuite le savon par \textbf{relargage} dans une solution concentrée de chlorure de sodium.
\end{corrige}

\begin{corrige}[Exercice 4 — Définitions]
\begin{enumerate}
\item \textbf{Acide carboxylique} : composé organique possédant le groupe carboxyle $-\ce{COOH}$ (formule générale $R\text{-COOH}$).
\item \textbf{Ester} : dérivé d'acide carboxylique dans lequel l'hydrogène du groupe $-\ce{COOH}$ est remplacé par un groupe alkyle $R'$ (fonction $R\text{-COO-}R'$).
\item \textbf{Saponification} : hydrolyse d'un ester par une base forte (\ce{NaOH} ou \ce{KOH}), réaction rapide et totale donnant un ion carboxylate (savon) et un alcool.
\item \textbf{Polycondensation} : réaction de polymérisation par étapes entre monomères bifonctionnels, accompagnée de l'élimination d'une petite molécule (eau, \ce{HCl}\ldots).
\item \textbf{Amide} : dérivé d'acide carboxylique dans lequel le groupe $-\ce{OH}$ du carboxyle est remplacé par $-\ce{NH2}$, $-\ce{NH}R$ ou $-\ce{N}R_2$ (fonction $-\ce{C(=O)}NR_2$).
\end{enumerate}
\end{corrige}

\begin{corrige}[Exercice 5 — Nomenclature des dérivés d'acide]
\textbf{1. Noms et familles :}
\begin{center}
\begin{tabularx}{\linewidth}{|l|X|X|}
\hline
\rowcolor{popblueL} \textbf{Formule} & \textbf{Nom} & \textbf{Famille} \\
\hline
\ce{CH3-CO-Cl} & Chlorure d'éthanoyle & Chlorure d'acyle \\
\hline
\ce{(CH3-CO)2O} & Anhydride éthanoïque & Anhydride d'acide \\
\hline
\ce{CH3-COO-CH2-CH3} & Éthanoate d'éthyle & Ester \\
\hline
\ce{CH3-CO-NH2} & Éthanamide & Amide \\
\hline
\end{tabularx}
\end{center}

\textbf{2. Équations-bilans à partir de l'acide éthanoïque :}
\[ \ce{CH3COOH + SOCl2 -> CH3COCl + SO2 ^ + HCl ^} \]
\[ \ce{2 CH3COOH ->[P4O10] (CH3CO)2O + H2O} \]
\[ \ce{CH3COOH + CH3CH2OH <=>[H+] CH3COOCH2CH3 + H2O} \quad \text{(lente, limitée)} \]
\[ \ce{CH3COOH + NH3 -> CH3COONH4 ->[\Delta] CH3CONH2 + H2O} \]

\textbf{3. Méthode rapide et totale pour l'ester :} utiliser le chlorure d'éthanoyle (ou l'anhydride éthanoïque) à la place de l'acide :
\[ \ce{CH3COCl + CH3CH2OH -> CH3COOCH2CH3 + HCl} \]
Réaction rapide, totale et exothermique.
\end{corrige}

\begin{corrige}[Exercice 6 — Températures d'ébullition et liaisons hydrogène]
\textbf{1. Formules semi-développées :}
\begin{itemize}
\item Pentane : \chemfig{CH_3-CH_2-CH_2-CH_2-CH_3}
\item Butan-1-ol : \chemfig{CH_3-CH_2-CH_2-CH_2-OH}
\item Acide propanoïque : \chemfig{CH_3-CH_2-C(=[1]O)-[7]OH}
\end{itemize}

\textbf{2.} Le pentane est une molécule \emph{apolaire} : seules de faibles interactions de Van der Waals (forces de London) existent entre ses molécules. Peu d'énergie suffit à les séparer, d'où sa température d'ébullition faible (\SI{36}{\celsius}).

\textbf{3.} Le butan-1-ol possède un groupe \ce{-OH} : ses molécules s'associent par des liaisons hydrogène \ce{O-H\cdots O} formant des chaînes, ce qui élève $T_{\text{éb}}$ à \SI{118}{\celsius}. L'acide propanoïque possède à la fois une liaison \ce{O-H} (donneur) et une liaison \ce{C=O} (accepteur) : deux molécules s'associent par \textbf{deux} liaisons hydrogène pour former un \emph{dimère cyclique} quasi rigide. La masse molaire effective double ($\approx \SI{148}{g.mol^{-1}}$) et les interactions sont plus fortes : il faut davantage d'énergie pour rompre ces associations, d'où $T_{\text{éb}} = \SI{141}{\celsius}$, supérieure à celle de l'alcool de masse molaire identique.

\textbf{4.} Schéma du dimère (voir la figure du cours) : les deux molécules \chemfig{CH_3-CH_2-C(=[1]O)-[7]OH} sont disposées tête-bêche, unies par deux liaisons hydrogène pointillées \ce{O-H\cdots O=C} formant un cycle à huit atomes.
\end{corrige}

\begin{corrige}[Exercice 7 — Synthèse d'un ester par le chlorure d'acyle]
\textbf{1. Équation-bilan :}
\[ \ce{CH3COCl + CH3CH2OH -> CH3COOCH2CH3 + HCl} \]
L'ester obtenu est l'\textbf{éthanoate d'éthyle}.

\textbf{2. Tableau d'avancement} (réaction totale, réactifs en proportions stœchiométriques) :
\begin{center}
\begin{tabular}{|l|c|c|c|c|}
\hline
\rowcolor{popblueL} \textbf{État (mol)} & \ce{CH3COCl} & \ce{CH3CH2OH} & \ce{CH3COOCH2CH3} & \ce{HCl} \\
\hline
Initial $x=0$ & 0,10 & 0,10 & 0 & 0 \\
\hline
En cours $x$ & $0{,}10-x$ & $0{,}10-x$ & $x$ & $x$ \\
\hline
Final $x_{\max}=0{,}10$ & 0 & 0 & 0,10 & 0,10 \\
\hline
\end{tabular}
\end{center}
Les deux réactifs sont limitants simultanément : $x_{\max} = \SI{0,10}{mol}$, donc $n_{\text{ester}} = \SI{0,10}{mol}$.

\textbf{3. Masse d'ester :} $M(\ce{C4H8O2}) = 4\times 12 + 8\times 1 + 2\times 16 = \SI{88}{g.mol^{-1}}$.
\[ m = n \times M = 0{,}10 \times 88 = \SI{8,8}{g} \]

\textbf{4. Par estérification directe :} $m_{\text{réel}} = 8{,}8 \times 0{,}67 \approx \SI{5,9}{g}$.
\textbf{Conclusion :} la voie par chlorure d'acyle est quantitative ($\SI{100}{\%}$, soit \SI{8,8}{g}) alors que l'estérification directe, limitée par l'équilibre, ne fournit que \SI{5,9}{g}. D'où l'intérêt industriel des dérivés réactifs, malgré leur coût et la corrosion due à \ce{HCl}.
\end{corrige}

\begin{corrige}[Exercice 8 — Fabrication artisanale du savon]
\textbf{1.} Il s'agit de la \textbf{saponification} : hydrolyse basique des triglycérides. Elle est \emph{rapide} (à chaud), \emph{totale} et \emph{irréversible}.

\textbf{2. Équation-bilan générale} (R = chaîne grasse, ex. \ce{C15H31} ou \ce{C17H35}) :
\[ \ce{C3H5(OOCR)3 + 3 KOH -> 3 RCOOK + C3H5(OH)3} \]
Triglycéride + potasse $\rightarrow$ 3 carboxylates de potassium (savon) + glycérol.

\textbf{3.} L'ion \ce{K+} est plus volumineux que \ce{Na+} : les carboxylates de potassium forment un réseau moins compact, moins cohésif, et sont plus solubles dans l'eau. Le savon obtenu reste pâteux : c'est un savon \emph{mou}. Avec \ce{Na+}, le réseau est plus compact et le savon est \emph{dur}.

\textbf{4.} Cette opération est le \textbf{relargage} (ou salage). L'ajout d'une solution concentrée de \ce{NaCl} diminue fortement la solubilité du savon (effet salin) : le savon précipite en surface et se sépare de la liqueur mère (eau, glycérol, excès de potasse), ce qui permet de le recueillir et de le purifier.
\end{corrige}

\begin{corrige}[Exercice 9 — Type Bac : identification d'un acide carboxylique]
\textbf{1.} Formule générale : $C_nH_{2n}O_2$.

\textbf{2.} $M = 12n + 2n \times 1 + 2 \times 16 = 14n + 32 = 74 \Rightarrow 14n = 42 \Rightarrow n = 3$. Formule brute : $C_3H_6O_2$.

\textbf{3.} Un seul isomère acide carboxylique : \chemfig{CH_3-CH_2-C(=[1]O)-[7]OH}, l'\textbf{acide propanoïque}. (Avec 3 C dont un dans \ce{-COOH}, aucune ramification n'est possible.)

\textbf{4. a)} $n_A = \dfrac{m}{M} = \dfrac{1{,}48}{74} = \SI{0,020}{mol}$ ; \quad $C = \dfrac{n_A}{V} = \dfrac{0{,}020}{0{,}200} = \SI{0,10}{mol.L^{-1}}$.

\textbf{b)} \[ \ce{CH3CH2COOH + H2O <=> CH3CH2COO- + H3O+} \]
Couple mis en jeu : \ce{CH3CH2COOH}/\ce{CH3CH2COO-}.

\textbf{c)} $\ce{[H3O+]} = 10^{-\text{pH}} = 10^{-2{,}9} \approx \SI{1,26 \times 10^{-3}}{mol.L^{-1}}$ ; \quad $\ce{[HO-]} = \dfrac{K_e}{\ce{[H3O+]}} = \dfrac{1{,}0\times 10^{-14}}{1{,}26\times 10^{-3}} \approx \SI{7,9 \times 10^{-12}}{mol.L^{-1}}$.

Comparaison : $\ce{[H3O+]} = \SI{1,26 \times 10^{-3}}{mol.L^{-1}} \ll C = \SI{0,10}{mol.L^{-1}}$ : la dissociation est partielle (environ \SI{1,3}{\%}), donc \textbf{A est un acide faible}.

\textbf{5.} \[ \ce{CH3CH2COOH + PCl5 -> CH3CH2COCl + POCl3 + HCl} \]
Produit organique : le \textbf{chlorure de propanoyle}.

\textbf{6.} \[ \ce{CH3CH2COCl + CH3OH -> CH3CH2COOCH3 + HCl} \]
Ester formé : le \textbf{propanoate de méthyle}. Caractéristiques : réaction \emph{rapide} et \emph{totale} (et exothermique).

\textbf{Barème indicatif (sur 10) :} Q1 : 0,5 ; Q2 : 1 ; Q3 : 1 ; Q4a : 1,5 ; Q4b : 1 ; Q4c : 2 ; Q5 : 1,5 ; Q6 : 1,5.

\textbf{Points de vigilance :} ne pas oublier la double flèche \ce{<=>} pour l'acide faible ; comparer $\ce{[H3O+]}$ à $C$ (et non conclure sur le seul pH) ; équilibrer correctement avec \ce{POCl3} et non \ce{PCl3O}.
\end{corrige}

\begin{corrige}[Exercice 10 — Type Bac : saponification de la tristéarine]
\textbf{1.} La saponification est l'hydrolyse d'un ester par une base forte (soude ou potasse) donnant un ion carboxylate (savon) et un alcool. Caractéristiques : réaction \emph{rapide} (à chaud) et \emph{totale} (irréversible).

\textbf{2.} Tristéarine (R = \ce{C17H35}) :
\chemfig{CH_2(-[1]O-C(=[2]O)-[6]R)-CH(-[1]O-C(=[2]O)-[6]R)-CH_2(-[1]O-C(=[2]O)-[6]R)}

\textbf{3. Équation-bilan :}
\[ \ce{C57H110O6 + 3 NaOH -> 3 C17H35COONa + C3H8O3} \]

\textbf{4.} $n_{\text{tristéarine}} = \dfrac{890}{890} = \SI{1,00}{mol}$.

\textbf{a) Tableau d'avancement :}
\begin{center}
\begin{tabular}{|l|c|c|c|c|}
\hline
\rowcolor{popblueL} \textbf{État (mol)} & \ce{C57H110O6} & \ce{NaOH} & \ce{C17H35COONa} & \ce{C3H8O3} \\
\hline
Initial & 1,00 & excès & 0 & 0 \\
\hline
En cours $x$ & $1{,}00-x$ & excès & $3x$ & $x$ \\
\hline
Final $x_{\max}=1{,}00$ & 0 & excès & 3,00 & 1,00 \\
\hline
\end{tabular}
\end{center}

\textbf{b)} $n_{\text{savon}} = 3 \times 1{,}00 = \SI{3,00}{mol}$. $M(\ce{C18H35O2Na}) = 18\times 12 + 35\times 1 + 2\times 16 + 23 = \SI{306}{g.mol^{-1}}$.
\[ m_{\text{savon}} = 3{,}00 \times 306 = \SI{918}{g} \]

\textbf{c)} $n_{\ce{NaOH},\min} = 3 \times 1{,}00 = \SI{3,00}{mol}$ ; $M(\ce{NaOH}) = 23 + 16 + 1 = \SI{40}{g.mol^{-1}}$.
\[ m_{\ce{NaOH}} = 3{,}00 \times 40 = \SI{120}{g} \]

\textbf{d)} $n_{\text{glycérol}} = \SI{1,00}{mol}$ ; $M(\ce{C3H8O3}) = 3\times 12 + 8\times 1 + 3\times 16 = \SI{92}{g.mol^{-1}}$.
\[ m_{\text{glycérol}} = \SI{92}{g} \]
Applications : hydratant en cosmétique et pharmacie, sirop, fabrication de la nitroglycérine (explosifs), antigel.

\textbf{5.} La soude en excès garantit que tout le triglycéride (réactif limitant) est saponifié : le rendement en savon est maximal. Précaution : la soude concentrée est \emph{corrosive} (brûlures de la peau et des yeux) : porter gants, lunettes et blouse, manipuler sous hotte, et en cas de contact rincer abondamment à l'eau.

\textbf{Barème indicatif (sur 10) :} Q1 : 1 ; Q2 : 1,5 ; Q3 : 1,5 ; Q4a : 1,5 ; Q4b : 1,5 ; Q4c : 1 ; Q4d : 1 ; Q5 : 1.

\textbf{Points de vigilance :} bien écrire le coefficient 3 devant \ce{NaOH} et devant le savon ; ne pas oublier que la saponification est totale (flèche simple \ce{->}) ; vérifier les masses molaires avant tout calcul de masse.
\end{corrige}

\begin{corrige}[Exercice 11 — Type Bac : le nylon 6,6 et le tergal]
\textbf{1. a)} La polycondensation est une polymérisation par étapes entre monomères polyfonctionnels, avec élimination d'une petite molécule (eau, \ce{HCl}\ldots). La polyaddition, au contraire, enchaîne les monomères \emph{sans} élimination de molécule.

\textbf{b)} \[ \ce{HOOC-(CH2)4-COOH + H2N-(CH2)6-NH2 -> HOOC-(CH2)4-CO-NH-(CH2)6-NH2 + H2O} \]
La petite molécule éliminée est l'\textbf{eau}.

\textbf{c)} Motif du nylon 6,6 :
\[ \fbox{\ce{-[CO-(CH2)4-CO-NH-(CH2)6-NH]-}} \]
Le nom \og 6,6 \fg{} indique que la diamine comporte \textbf{6} atomes de carbone et le diacide \textbf{6} atomes de carbone.

\textbf{d)} $n_{\text{motifs}} = \dfrac{m}{M_{\text{motif}}} = \dfrac{1130}{226} = \SI{5,0}{mol}$ de motifs, soit $5{,}0 \times 6{,}02\times 10^{23} = 3{,}01 \times 10^{24}$ motifs.

\textbf{2. a)} \[ \ce{n HOOC-C6H4-COOH + n HO-CH2-CH2-OH -> \text{tergal} + 2n H2O} \]

\textbf{b)} Motif du tergal (la fonction ester \ce{-C(=O)-O-} est entourée) :
\[ \fbox{\ce{-[CO-C6H4-\underset{\text{ester}}{\underline{CO-O}}-CH2-CH2-O]-}} \]

\textbf{c)} Le tergal est un \emph{polyester} car ses motifs sont reliés par des fonctions \textbf{ester} \ce{-CO-O-} ; le nylon est un \emph{polyamide} car ses motifs sont reliés par des fonctions \textbf{amide} \ce{-CO-NH-}.

\textbf{3.} Au Cameroun : \textbf{nylon 6,6} : fils et filets de pêche, cordages, tissus techniques, brosses ; \textbf{tergal (PET)} : bouteilles d'eau et de jus, emballages alimentaires, fibres textiles mélangées (pagnes polyester-coton).

\textbf{Barème indicatif (sur 10) :} 1a : 1 ; 1b : 1,5 ; 1c : 1,5 ; 1d : 1 ; 2a : 1,5 ; 2b : 1,5 ; 2c : 1 ; 3 : 1.

\textbf{Points de vigilance :} dans le motif, bien faire apparaître les liaisons libres aux extrémités (crochets) ; ne pas confondre polyaddition et polycondensation ; vérifier que l'équation élimine bien 2 \ce{H2O} par motif formé (deux fonctions par monomère).
\end{corrige}

\newpage

\coursec{Fiche bilan}

\begin{popfigure}
\begin{center}
\begin{tikzpicture}[
  mindmap/.style={font=\small, align=center, rounded corners=3pt, inner sep=3pt},
  racine/.style={mindmap, fill=popblue, text=white, font=\bfseries, minimum width=3.2cm, minimum height=1.1cm},
  branche/.style={mindmap, text=popdark, minimum width=2.9cm},
  lien/.style={thick, popmuted}
]
\node[racine] (C) at (0,0) {Acides\\carboxyliques\\et dérivés};
\node[branche, fill=popblueL] (N) at (-4.6,2.4) {\textbf{Nomenclature}\\R--COOH\\acide ...anoïque};
\node[branche, fill=poporangeL] (P) at (0,3.2) {\textbf{Propriétés physiques}\\liaisons H, dimères\\solubilité, $T_{\text{éb}}$ élevées};
\node[branche, fill=popgreenL] (A) at (4.6,2.4) {\textbf{Acidité}\\\ce{RCOOH + H2O <=>}\\\ce{RCOO- + H3O+}};
\node[branche, fill=poppinkL] (D) at (-4.6,-2.4) {\textbf{Dérivés}\\chlorures R--COCl\\anhydrides (R--CO)$_2$O};
\node[branche, fill=popgoldL] (E) at (0,-3.2) {\textbf{Esters}\\estérification lente, limitée\\synthèses rapides et totales};
\node[branche, fill=poppurpleL] (S) at (4.6,-2.4) {\textbf{Savon et amides}\\saponification\\nylon 6,6 ; tergal};
\draw[lien] (C) -- (N);
\draw[lien] (C) -- (P);
\draw[lien] (C) -- (A);
\draw[lien] (C) -- (D);
\draw[lien] (C) -- (E);
\draw[lien] (C) -- (S);
\end{tikzpicture}
\end{center}
\legende{Carte mentale de la leçon : les six compétences clés autour du groupe carboxyle \ce{-COOH}.}
\end{popfigure}

\begin{bilan}[L'essentiel de la leçon]
\textbf{Définitions clés}
\begin{itemize}
  \item \cle{Acide carboxylique} : composé de formule \ce{R-COOH} ; le groupe \cle{carboxyle} \ce{-C(=O)-OH} porte le carbone en position 1.
  \item \cle{Dérivés d'acides} : chlorure d'acyle \ce{R-COCl}, anhydride \ce{(R-CO)2O}, ester \ce{R-COO-}$R'$, amide \ce{R-CO-N}$R'R''$.
  \item \cle{Estérification} : réaction acide + alcool $\rightleftharpoons$ ester + eau : \emph{lente, limitée, athermique}.
  \item \cle{Saponification} : hydrolyse basique d'un ester : \emph{totale}, donne un carboxylate (savon si l'acide est gras).
  \item \cle{Polycondensation} : polymérisation avec élimination d'une petite molécule (eau) : nylon 6,6 (polyamide), tergal (polyester).
\end{itemize}

\textbf{Équations-types à connaître par c\oe ur}
\begin{align*}
&\ce{RCOOH + H2O <=> RCOO- + H3O+} &&\text{(couple \ce{RCOOH}/\ce{RCOO-})}\\
&\ce{RCOOH + PCl5 -> RCOCl + POCl3 + HCl} &&\text{(chlorure d'acyle)}\\
&\ce{2 RCOOH ->[P4O10] (RCO)2O + H2O} &&\text{(anhydride)}\\
&\ce{RCOOH + R^{\prime}OH <=> RCOOR^{\prime} + H2O} &&\text{(estérification)}\\
&\ce{RCOCl + R^{\prime}OH -> RCOOR^{\prime} + HCl} &&\text{(rapide et totale)}\\
&\ce{RCOOR^{\prime} + OH- -> RCOO- + R^{\prime}OH} &&\text{(saponification)}\\
&\ce{RCOOH + NH3 -> RCOO- + NH4+ ->[\Delta] RCONH2 + H2O} &&\text{(amide)}
\end{align*}

\textbf{Nomenclature express}
\begin{center}
\begin{tabularx}{\linewidth}{|l|X|X|}
\hline
\rowcolor{popblueL} \textbf{Famille} & \textbf{Suffixe / règle} & \textbf{Exemple} \\
\hline
Acide & acide ...\textbf{anoïque} & \ce{CH3-CH2-COOH} : acide propanoïque \\
\hline
Chlorure d'acyle & chlorure de ...\textbf{anoyle} & \ce{CH3-COCl} : chlorure d'éthanoyle \\
\hline
Anhydride & anhydride ...\textbf{anoïque} & \ce{(CH3-CO)2O} : anhydride éthanoïque \\
\hline
Ester & ...\textbf{anoate} de ...\textbf{yle} & \ce{CH3-COO-C2H5} : éthanoate d'éthyle \\
\hline
Amide & ...\textbf{anamide} & \ce{CH3-CONH2} : éthanamide \\
\hline
\end{tabularx}
\end{center}

\textbf{Méthodes express}
\begin{itemize}
  \item Chaîne principale = chaîne la plus longue contenant \ce{-COOH} ; le carbone du carboxyle est \textbf{toujours n° 1} ; suffixe « -dioïque » pour les diacides.
  \item Propriétés physiques : les liaisons hydrogène (dimères) expliquent les températures d'ébullition \emph{élevées} et la solubilité dans l'eau des acides à courte chaîne ($n \leq 4$).
  \item Pour un rendement d'estérification élevé : utiliser le \cle{chlorure d'acyle} ou l'\cle{anhydride} (réactions rapides et totales) plutôt que l'acide.
  \item Savon : triglycéride + 3 \ce{NaOH} $\rightarrow$ 3 carboxylates de sodium + glycérol ; relargage au sel, puis lavage et séchage.
\end{itemize}
\end{bilan}

\begin{aretenir}[Checklist avant le Bac]
\begin{itemize}
  \item[$\square$] Je sais écrire la formule semi-développée d'un acide carboxylique (saturé, insaturé, polyacide) à partir de son nom.
  \item[$\square$] Je nomme sans hésiter un acide, un chlorure d'acyle, un anhydride, un ester et un amide.
  \item[$\square$] Je sais que le carbone du groupe carboxyle est toujours numéroté 1.
  \item[$\square$] J'explique les températures d'ébullition élevées et la solubilité par les liaisons hydrogène (dimères).
  \item[$\square$] J'écris l'équation de l'action de l'eau sur un acide et je cite le couple \ce{RCOOH}/\ce{RCOO-}.
  \item[$\square$] J'écris les équations de formation des chlorures d'acyle (\ce{PCl5}, \ce{SOCl2}) et des anhydrides (\ce{P4O10}).
  \item[$\square$] Je connais les trois caractères de l'estérification : lente, limitée, athermique.
  \item[$\square$] Je sais que les synthèses d'esters via chlorures d'acyle ou anhydrides sont rapides et totales.
  \item[$\square$] Je définis la saponification et j'écris son équation-bilan avec un triglycéride.
  \item[$\square$] Je décris les étapes de la fabrication artisanale du savon (huile de palme, potasse, relargage).
  \item[$\square$] J'écris les équations de polycondensation du nylon 6,6 et du tergal.
  \item[$\square$] Je distingue hydrolyse (réversible) et saponification (totale) d'un ester.
\end{itemize}
\end{aretenir}

\findecours

\end{document}
